% !TEX program = lualatex
% ---------------------------------------------------------------------------------------------
%  Readback note 14: The height of a ratio of lift volumes: twisted theta forms on one line and Hecke values
%
%  Build:  python3 tools/rb.py note <slug>
%
%  One unit of a Lean artifact, rewritten as ordinary mathematics. The prose is hand-written. Lean names,
%  line ranges and counts come from data.tex; quoted Lean statements are copied from the artifact by the build.
%
%  Rules of the text: standard terminology; every symbol is explained where it first appears and listed again
%  in the notation at the end; no letter has two meanings; a proof is complete or says what it takes for granted;
%  Lean names appear only in the LEAN lines and in the concordance.
%
%  The signature rule (the build enforces it): a function, relation or predicate that the artifact defines is
%  written only through its notation macro, which the build makes from the spec (`notation`) with one slot for
%  each explicit argument of the Lean definition. No argument is fixed and dropped, also not "for brevity" and
%  not for one section. Every argument and every hypothesis of a stated declaration stands inside its
%  statement, or in a Setting that the statement cites by number; only implicit and instance arguments (the
%  ambient objects) may stand in the conventions.
% ---------------------------------------------------------------------------------------------
\documentclass[11pt,a4paper]{article}
\ifdefined\pdfvariable\pdfvariable suppressoptionalinfo \numexpr 1+2+4+8\relax\fi
\usepackage[a4paper,hmargin=35mm,top=22mm,bottom=27mm,footskip=12mm,nohead]{geometry}

\usepackage{fontspec}
\directlua{
  luaotfload.add_fallback("gfallback", {"STIXTwoMath-Regular.otf:mode=node;", "DejaVuSans.ttf:mode=node;"})
  luaotfload.add_fallback("gmonofallback", {"DejaVuSansMono.ttf:mode=node;", "STIXTwoMath-Regular.otf:mode=node;"})
}
\setmainfont{NewCM10}[Extension=.otf,UprightFont=*-Book,ItalicFont=*-BookItalic,
  BoldFont=*-Bold,BoldItalicFont=*-BoldItalic,SmallCapsFont=*-Book,SmallCapsFeatures={Letters=SmallCaps},
  RawFeature={fallback=gfallback}]
\setsansfont{NewCMSans10}[Extension=.otf,UprightFont=*-Book,ItalicFont=*-BookOblique,
  BoldFont=*-Bold,BoldItalicFont=*-BoldOblique,RawFeature={fallback=gfallback}]
\setmonofont{JetBrainsMono}[Extension=.otf,UprightFont=*-Regular,ItalicFont=*-Italic,
  BoldFont=*-SemiBold,BoldItalicFont=*-SemiBoldItalic,Scale=MatchLowercase,
  RawFeature={-calt;-liga;fallback=gmonofallback}]
\usepackage{amsmath}
\usepackage{amsthm}
\usepackage[math-style=TeX]{unicode-math}
\setmathfont{NewCMMath-Book.otf}
\emergencystretch=2.5em\relax

\usepackage{xcolor}
\usepackage{tikz}
\usetikzlibrary{calc,arrows.meta,positioning,backgrounds}
\usepackage[protrusion=true,expansion=true]{microtype}
\usepackage{enumitem}
\usepackage{needspace}
\usepackage{booktabs}
\usepackage{array}
\usepackage{longtable}
\usepackage{fancyhdr}
\usepackage{lastpage}
\usepackage{multicol}
\usepackage[hidelinks,pdfusetitle]{hyperref}

\definecolor{ink}{HTML}{0F172A}
\definecolor{inksoft}{HTML}{475569}
\definecolor{inkfaint}{HTML}{94A3B8}
\definecolor{rulegrey}{HTML}{CBD5E1}
\definecolor{arith}{HTML}{C2410C}
\definecolor{est}{HTML}{7C3AED}
\definecolor{ana}{HTML}{0891B2}
\definecolor{ok}{HTML}{047857}
\definecolor{leancol}{HTML}{3730A3}

\title{The height of a ratio of lift volumes: twisted theta forms on one line and Hecke values}
\author{grh-compact}
\date{}

% ------------------------------------------------------------------------------------- macros
\newcommand{\leanfont}{\ttfamily\fontsize{8.6}{10.5}\selectfont\color{leancol}}
\newcommand{\lean}[1]{\csname LN@#1\endcsname}
\newcommand{\fillin}[1]{{\color{arith}\itshape[#1]}}  % a place of the template that is not written yet; the build reports every one that is left
\newcommand{\leanquote}[1]{\csname LQ@#1\endcsname}
% generated by rb.py note; do not edit
\newcommand{\ConcDecls}{77}
\newcommand{\Pin}{72fead10}
\newcommand{\PinText}{that of commit {\leanfont 72fead10}}
\newcommand{\ArtifactName}{grh-compact}
\expandafter\newcommand\csname LN@onelemaxOpNorm\endcsname{\href{https://github.com/lsernau/grh-lean/blob/72fead102f663d5c7a407964a13dc0cf9d977207/grh-compact/src/New/Rows/SparseThetaSum.lean\#L21-L22}{{\leanfont one\_le\_maxOpNorm}}}
\expandafter\newcommand\csname LN@integralmulholMinors\endcsname{\href{https://github.com/lsernau/grh-lean/blob/72fead102f663d5c7a407964a13dc0cf9d977207/grh-compact/src/New/Rows/SparseThetaSum.lean\#L30-L56}{{\leanfont integral\_mul\_holMinorsPi\_le}}}
\expandafter\newcommand\csname LN@isolation\endcsname{\href{https://github.com/lsernau/grh-lean/blob/72fead102f663d5c7a407964a13dc0cf9d977207/grh-compact/src/New/Rows/SparseThetaSum.lean\#L62-L84}{{\leanfont isolation}}}
\expandafter\newcommand\csname LN@firstMinimummulle\endcsname{\href{https://github.com/lsernau/grh-lean/blob/72fead102f663d5c7a407964a13dc0cf9d977207/grh-compact/src/New/Rows/SparseThetaSum.lean\#L99-L112}{{\leanfont firstMinimum\_mul\_le\_hom}}}
\expandafter\newcommand\csname LN@dataInnertraceChar\endcsname{\href{https://github.com/lsernau/grh-lean/blob/72fead102f663d5c7a407964a13dc0cf9d977207/grh-compact/src/New/Rows/SparseThetaSum.lean\#L118-L128}{{\leanfont dataInner\_traceChar}}}
\expandafter\newcommand\csname LN@normmodKernelsub\endcsname{\href{https://github.com/lsernau/grh-lean/blob/72fead102f663d5c7a407964a13dc0cf9d977207/grh-compact/src/New/Rows/SparseThetaSum.lean\#L133-L171}{{\leanfont norm\_modKernel\_sub\_realInner\_le}}}
\expandafter\newcommand\csname LN@carrierFormspec\endcsname{\href{https://github.com/lsernau/grh-lean/blob/72fead102f663d5c7a407964a13dc0cf9d977207/grh-compact/src/New/Rows/NaturalFormNear.lean\#L23-L35}{{\leanfont carrierForm\_spec}}}
\expandafter\newcommand\csname LN@carrierFormnezero\endcsname{\href{https://github.com/lsernau/grh-lean/blob/72fead102f663d5c7a407964a13dc0cf9d977207/grh-compact/src/New/Rows/NaturalFormNear.lean\#L39-L50}{{\leanfont carrierForm\_ne\_zero}}}
\expandafter\newcommand\csname LN@negIndexposRowIndexspec\endcsname{\href{https://github.com/lsernau/grh-lean/blob/72fead102f663d5c7a407964a13dc0cf9d977207/grh-compact/src/New/Rows/NaturalFormNear.lean\#L54-L58}{{\leanfont Data.\allowbreak{}negIndex\_posRowIndex\_spec}}}
\expandafter\newcommand\csname LN@existsfNaturalapprox\endcsname{\href{https://github.com/lsernau/grh-lean/blob/72fead102f663d5c7a407964a13dc0cf9d977207/grh-compact/src/New/Rows/NaturalFormNear.lean\#L65-L146}{{\leanfont Data.\allowbreak{}exists\_fNatural\_approx}}}
\expandafter\newcommand\csname LN@cmSignstar\endcsname{\href{https://github.com/lsernau/grh-lean/blob/72fead102f663d5c7a407964a13dc0cf9d977207/grh-compact/src/New/Rows/Separation.lean\#L29-L30}{{\leanfont cmSign\_star}}}
\expandafter\newcommand\csname LN@isRowTypeauxPair\endcsname{\href{https://github.com/lsernau/grh-lean/blob/72fead102f663d5c7a407964a13dc0cf9d977207/grh-compact/src/New/Rows/Separation.lean\#L35-L41}{{\leanfont isRowType\_auxPair}}}
\expandafter\newcommand\csname LN@eSign\endcsname{\href{https://github.com/lsernau/grh-lean/blob/72fead102f663d5c7a407964a13dc0cf9d977207/grh-compact/src/New/Rows/Separation.lean\#L44-L44}{{\leanfont Data.\allowbreak{}eSign}}}
\expandafter\newcommand\csname LN@heckeCount\endcsname{\href{https://github.com/lsernau/grh-lean/blob/72fead102f663d5c7a407964a13dc0cf9d977207/grh-compact/src/New/Rows/Separation.lean\#L47-L47}{{\leanfont Data.\allowbreak{}heckeCount}}}
\expandafter\newcommand\csname LN@linePair\endcsname{\href{https://github.com/lsernau/grh-lean/blob/72fead102f663d5c7a407964a13dc0cf9d977207/grh-compact/src/New/Rows/Separation.lean\#L51-L59}{{\leanfont Data.\allowbreak{}linePair}}}
\expandafter\newcommand\csname LN@heckeValueProd\endcsname{\href{https://github.com/lsernau/grh-lean/blob/72fead102f663d5c7a407964a13dc0cf9d977207/grh-compact/src/New/Rows/Separation.lean\#L63-L65}{{\leanfont Data.\allowbreak{}heckeValueProd}}}
\expandafter\newcommand\csname LN@existsisFlat\endcsname{\href{https://github.com/lsernau/grh-lean/blob/72fead102f663d5c7a407964a13dc0cf9d977207/grh-compact/src/New/Rows/Separation.lean\#L73-L78}{{\leanfont Data.\allowbreak{}exists\_isFlat}}}
\expandafter\newcommand\csname LN@heightLenormtrivialLift\endcsname{\href{https://github.com/lsernau/grh-lean/blob/72fead102f663d5c7a407964a13dc0cf9d977207/grh-compact/src/New/Rows/Separation.lean\#L83-L95}{{\leanfont Data.\allowbreak{}heightLe\_norm\_trivialLift\_sq}}}
\expandafter\newcommand\csname LN@cosetKernelSumeqfNatural\endcsname{\href{https://github.com/lsernau/grh-lean/blob/72fead102f663d5c7a407964a13dc0cf9d977207/grh-compact/src/New/Rows/Separation.lean\#L107-L157}{{\leanfont Data.\allowbreak{}cosetKernelSum\_eq\_fNatural}}}
\expandafter\newcommand\csname LQ@cosetKernelSumeqfNatural\endcsname{\leanquoteline{0}{lemma Data.cosetKernelSum\_eq\_fNatural \{P\} \{D : Data P\} \{a\} [Fintype a] j : ∃ C, ∃ c : D.Twist → ℂ,}\leanquoteline{2}{∀ ϑ, ‖c ϑ‖ \textasciicircum{} 2 = C / heckeValueProd j ϑ ∧ ∀ (u : Matrix \_ a (D.untwistedModulus 2).ResRing) i g}\leanquoteline{2}{(w : Matrix a \_ (D.untwistedModulus 2).ResRing), u * w = i.1 → fNatural j i g = c ϑ *}\leanquoteline{2}{cosetKernelSum (D.auxPair ϑ j) (bigChar ϑ j) w (.trans .inl .inl) (lineMultiplier ϑ j i) g}}
\expandafter\newcommand\csname LN@goodFormsrapid\endcsname{\href{https://github.com/lsernau/grh-lean/blob/72fead102f663d5c7a407964a13dc0cf9d977207/grh-compact/src/New/Rows/Separation.lean\#L163-L165}{{\leanfont Data.\allowbreak{}goodForms\_mem\_rapidOfLevel}}}
\expandafter\newcommand\csname LN@mulfNaturalrapid\endcsname{\href{https://github.com/lsernau/grh-lean/blob/72fead102f663d5c7a407964a13dc0cf9d977207/grh-compact/src/New/Rows/Separation.lean\#L169-L182}{{\leanfont Data.\allowbreak{}mul\_fNatural\_mem\_rapidOfLevel\_and\_hol}}}
\expandafter\newcommand\csname LN@tableTheta\endcsname{\href{https://github.com/lsernau/grh-lean/blob/72fead102f663d5c7a407964a13dc0cf9d977207/grh-compact/src/New/Rows/Separation.lean\#L190-L193}{{\leanfont Data.\allowbreak{}tableTheta}}}
\expandafter\newcommand\csname LN@eqzeroofintegral\endcsname{\href{https://github.com/lsernau/grh-lean/blob/72fead102f663d5c7a407964a13dc0cf9d977207/grh-compact/src/New/Rows/Separation.lean\#L199-L236}{{\leanfont Data.\allowbreak{}eq\_zero\_of\_integral\_fNatural\_tableTheta}}}
\expandafter\newcommand\csname LQ@eqzeroofintegral\endcsname{\leanquoteline{0}{lemma Data.eq\_zero\_of\_integral\_fNatural\_tableTheta \{P\} \{D : Data P\} \{j : D.LiftIndex\} (F :}\leanquoteline{2}{D.goodForms (D.spaceW j)) (h : ∀ i ζ, ∫ g, conj (F.1.1 g * fNatural j i g) * tableTheta j ζ}\leanquoteline{2}{g ∂(D.spaceW j).quotMeasure = 0) : F = 0}}
\expandafter\newcommand\csname LN@existstwistedclassSum\endcsname{\href{https://github.com/lsernau/grh-lean/blob/72fead102f663d5c7a407964a13dc0cf9d977207/grh-compact/src/New/Rows/HeckeValueHeight.lean\#L28-L45}{{\leanfont exists\_twisted\_classSum\_ge}}}
\expandafter\newcommand\csname LN@existsfactorisation\endcsname{\href{https://github.com/lsernau/grh-lean/blob/72fead102f663d5c7a407964a13dc0cf9d977207/grh-compact/src/New/Rows/HeckeValueHeight.lean\#L52-L56}{{\leanfont Data.\allowbreak{}exists\_factorisation\_le}}}
\expandafter\newcommand\csname LN@xiSum\endcsname{\href{https://github.com/lsernau/grh-lean/blob/72fead102f663d5c7a407964a13dc0cf9d977207/grh-compact/src/New/Rows/HeckeValueHeight.lean\#L60-L62}{{\leanfont Data.\allowbreak{}Pair.\allowbreak{}xiSum}}}
\expandafter\newcommand\csname LN@xiSummem\endcsname{\href{https://github.com/lsernau/grh-lean/blob/72fead102f663d5c7a407964a13dc0cf9d977207/grh-compact/src/New/Rows/HeckeValueHeight.lean\#L64-L69}{{\leanfont Data.\allowbreak{}Pair.\allowbreak{}xiSum\_mem}}}
\expandafter\newcommand\csname LN@smulapplyxiSum\endcsname{\href{https://github.com/lsernau/grh-lean/blob/72fead102f663d5c7a407964a13dc0cf9d977207/grh-compact/src/New/Rows/HeckeValueHeight.lean\#L71-L82}{{\leanfont Data.\allowbreak{}Pair.\allowbreak{}smul\_apply\_xiSum}}}
\expandafter\newcommand\csname LN@existstestform\endcsname{\href{https://github.com/lsernau/grh-lean/blob/72fead102f663d5c7a407964a13dc0cf9d977207/grh-compact/src/New/Rows/HeckeValueHeight.lean\#L87-L148}{{\leanfont Data.\allowbreak{}exists\_test\_form\_inner\_kernelForm}}}
\expandafter\newcommand\csname LN@heightLeinvratio\endcsname{\href{https://github.com/lsernau/grh-lean/blob/72fead102f663d5c7a407964a13dc0cf9d977207/grh-compact/src/New/Rows/HeckeValueHeight.lean\#L154-L216}{{\leanfont Data.\allowbreak{}heightLe\_inv\_ratio}}}
\expandafter\newcommand\csname LQ@heightLeinvratio\endcsname{\leanquoteline{0}{lemma heightLe\_inv\_ratio \{θ\} (o : D.IsFitFamily θ) (j : D.LiftIndex) (p : D.Pair) (hE :}\leanquoteline{2}{p.IsRankOne j) : HeightLe (6 * D.B₀ - D.dirCost) fun ι ↦ (p.ratio ((θ ι).powTwist j))⁻¹}}
\expandafter\newcommand\csname LN@heightLeinvheckeValueProd\endcsname{\href{https://github.com/lsernau/grh-lean/blob/72fead102f663d5c7a407964a13dc0cf9d977207/grh-compact/src/New/Rows/HeckeValueHeight.lean\#L223-L237}{{\leanfont Data.\allowbreak{}heightLe\_inv\_heckeValueProd\_div}}}
\expandafter\newcommand\csname LN@compressedGramOp\endcsname{\href{https://github.com/lsernau/grh-lean/blob/72fead102f663d5c7a407964a13dc0cf9d977207/grh-compact/src/New/Rows/LiftVolumeHeight.lean\#L26-L27}{{\leanfont Data.\allowbreak{}compressedGramOp}}}
\expandafter\newcommand\csname LN@adjointcomp\endcsname{\href{https://github.com/lsernau/grh-lean/blob/72fead102f663d5c7a407964a13dc0cf9d977207/grh-compact/src/New/Rows/LiftVolumeHeight.lean\#L31-L32}{{\leanfont Data.\allowbreak{}adjoint\_comp\_eq\_conj\_smul}}}
\expandafter\newcommand\csname LN@detcompressedGramOp\endcsname{\href{https://github.com/lsernau/grh-lean/blob/72fead102f663d5c7a407964a13dc0cf9d977207/grh-compact/src/New/Rows/LiftVolumeHeight.lean\#L35-L37}{{\leanfont Data.\allowbreak{}det\_compressedGramOp}}}
\expandafter\newcommand\csname LN@normdetpadic\endcsname{\href{https://github.com/lsernau/grh-lean/blob/72fead102f663d5c7a407964a13dc0cf9d977207/grh-compact/src/New/Rows/LiftVolumeHeight.lean\#L47-L70}{{\leanfont Data.\allowbreak{}norm\_det\_compressedGramOp\_padic\_le}}}
\expandafter\newcommand\csname LN@normdetarch\endcsname{\href{https://github.com/lsernau/grh-lean/blob/72fead102f663d5c7a407964a13dc0cf9d977207/grh-compact/src/New/Rows/LiftVolumeHeight.lean\#L76-L95}{{\leanfont Data.\allowbreak{}norm\_det\_compressedGramOp\_arch\_le}}}
\expandafter\newcommand\csname LN@heightLeliftVolumeratio\endcsname{\href{https://github.com/lsernau/grh-lean/blob/72fead102f663d5c7a407964a13dc0cf9d977207/grh-compact/src/New/Rows/LiftVolumeHeight.lean\#L103-L127}{{\leanfont Data.\allowbreak{}heightLe\_liftVolume\_ratio}}}
\expandafter\newcommand\csname LQ@heightLeliftVolumeratio\endcsname{\leanquoteline{0}{lemma heightLe\_liftVolume\_ratio \{θ\} (O : D.IsFitFamily θ) (j : D.LiftIndex) :}\leanquoteline{2}{HeightLe (D.rowDim * (D.k * (D.logB * D.B₀ - D.dirCost))) fun ι ↦}\leanquoteline{2}{D.liftVolume 1 j / D.liftVolume (θ ι) j}}
\expandafter\newcommand\csname LN@untwistedModulus\endcsname{\href{https://github.com/lsernau/grh-lean/blob/72fead102f663d5c7a407964a13dc0cf9d977207/grh-compact/src/New/Rows/LineCoordinates.lean\#L24-L24}{{\leanfont Data.\allowbreak{}untwistedModulus}}}
\expandafter\newcommand\csname LN@idealuntwistedModulus\endcsname{\href{https://github.com/lsernau/grh-lean/blob/72fead102f663d5c7a407964a13dc0cf9d977207/grh-compact/src/New/Rows/LineCoordinates.lean\#L26-L27}{{\leanfont Data.\allowbreak{}ideal\_untwistedModulus\_eq\_span}}}
\expandafter\newcommand\csname LN@CarrierIndex\endcsname{\href{https://github.com/lsernau/grh-lean/blob/72fead102f663d5c7a407964a13dc0cf9d977207/grh-compact/src/New/Rows/LineCoordinates.lean\#L33-L33}{{\leanfont Data.\allowbreak{}CarrierIndex}}}
\expandafter\newcommand\csname LN@FlatIndex\endcsname{\href{https://github.com/lsernau/grh-lean/blob/72fead102f663d5c7a407964a13dc0cf9d977207/grh-compact/src/New/Rows/LineCoordinates.lean\#L36-L36}{{\leanfont Data.\allowbreak{}FlatIndex}}}
\expandafter\newcommand\csname LN@SmallIndex\endcsname{\href{https://github.com/lsernau/grh-lean/blob/72fead102f663d5c7a407964a13dc0cf9d977207/grh-compact/src/New/Rows/LineCoordinates.lean\#L39-L39}{{\leanfont Data.\allowbreak{}SmallIndex}}}
\expandafter\newcommand\csname LN@SmallClasses\endcsname{\href{https://github.com/lsernau/grh-lean/blob/72fead102f663d5c7a407964a13dc0cf9d977207/grh-compact/src/New/Rows/LineCoordinates.lean\#L43-L44}{{\leanfont Data.\allowbreak{}SmallClasses}}}
\expandafter\newcommand\csname LN@basisVec\endcsname{\href{https://github.com/lsernau/grh-lean/blob/72fead102f663d5c7a407964a13dc0cf9d977207/grh-compact/src/New/Rows/LineCoordinates.lean\#L49-L49}{{\leanfont Data.\allowbreak{}basisVec}}}
\expandafter\newcommand\csname LN@LineChoice\endcsname{\href{https://github.com/lsernau/grh-lean/blob/72fead102f663d5c7a407964a13dc0cf9d977207/grh-compact/src/New/Rows/LineCoordinates.lean\#L55-L56}{{\leanfont Data.\allowbreak{}LineChoice}}}
\expandafter\newcommand\csname LN@shortRows\endcsname{\href{https://github.com/lsernau/grh-lean/blob/72fead102f663d5c7a407964a13dc0cf9d977207/grh-compact/src/New/Rows/LineCoordinates.lean\#L60-L61}{{\leanfont Data.\allowbreak{}shortRows}}}
\expandafter\newcommand\csname LN@twistedChi\endcsname{\href{https://github.com/lsernau/grh-lean/blob/72fead102f663d5c7a407964a13dc0cf9d977207/grh-compact/src/New/Rows/LineCoordinates.lean\#L67-L67}{{\leanfont Data.\allowbreak{}twistedChi}}}
\expandafter\newcommand\csname LN@twistedChihasWeight\endcsname{\href{https://github.com/lsernau/grh-lean/blob/72fead102f663d5c7a407964a13dc0cf9d977207/grh-compact/src/New/Rows/LineCoordinates.lean\#L69-L70}{{\leanfont Data.\allowbreak{}twistedChi\_hasWeight}}}
\expandafter\newcommand\csname LN@PadIndex\endcsname{\href{https://github.com/lsernau/grh-lean/blob/72fead102f663d5c7a407964a13dc0cf9d977207/grh-compact/src/New/Rows/LineCoordinates.lean\#L74-L74}{{\leanfont Data.\allowbreak{}PadIndex}}}
\expandafter\newcommand\csname LN@BigIndex\endcsname{\href{https://github.com/lsernau/grh-lean/blob/72fead102f663d5c7a407964a13dc0cf9d977207/grh-compact/src/New/Rows/LineCoordinates.lean\#L77-L77}{{\leanfont Data.\allowbreak{}BigIndex}}}
\expandafter\newcommand\csname LN@padSigns\endcsname{\href{https://github.com/lsernau/grh-lean/blob/72fead102f663d5c7a407964a13dc0cf9d977207/grh-compact/src/New/Rows/LineCoordinates.lean\#L81-L81}{{\leanfont Data.\allowbreak{}padSigns}}}
\expandafter\newcommand\csname LN@negIndex\endcsname{\href{https://github.com/lsernau/grh-lean/blob/72fead102f663d5c7a407964a13dc0cf9d977207/grh-compact/src/New/Rows/LineCoordinates.lean\#L85-L86}{{\leanfont Data.\allowbreak{}negIndex}}}
\expandafter\newcommand\csname LN@posRowIndex\endcsname{\href{https://github.com/lsernau/grh-lean/blob/72fead102f663d5c7a407964a13dc0cf9d977207/grh-compact/src/New/Rows/LineCoordinates.lean\#L90-L91}{{\leanfont Data.\allowbreak{}posRowIndex}}}
\expandafter\newcommand\csname LN@untwistedAuxPair\endcsname{\href{https://github.com/lsernau/grh-lean/blob/72fead102f663d5c7a407964a13dc0cf9d977207/grh-compact/src/New/Rows/LineCoordinates.lean\#L94-L94}{{\leanfont Data.\allowbreak{}untwistedAuxPair}}}
\expandafter\newcommand\csname LN@DetIndex\endcsname{\href{https://github.com/lsernau/grh-lean/blob/72fead102f663d5c7a407964a13dc0cf9d977207/grh-compact/src/New/Rows/LineCoordinates.lean\#L98-L99}{{\leanfont Data.\allowbreak{}DetIndex}}}
\expandafter\newcommand\csname LN@linePoly\endcsname{\href{https://github.com/lsernau/grh-lean/blob/72fead102f663d5c7a407964a13dc0cf9d977207/grh-compact/src/New/Rows/LineCoordinates.lean\#L104-L106}{{\leanfont Data.\allowbreak{}linePoly}}}
\expandafter\newcommand\csname LN@bigCharExp\endcsname{\href{https://github.com/lsernau/grh-lean/blob/72fead102f663d5c7a407964a13dc0cf9d977207/grh-compact/src/New/Rows/LineCoordinates.lean\#L111-L111}{{\leanfont Data.\allowbreak{}bigCharExp}}}
\expandafter\newcommand\csname LN@carrierChar\endcsname{\href{https://github.com/lsernau/grh-lean/blob/72fead102f663d5c7a407964a13dc0cf9d977207/grh-compact/src/New/Rows/LineCoordinates.lean\#L115-L115}{{\leanfont Data.\allowbreak{}carrierChar}}}
\expandafter\newcommand\csname LN@carrierCharhasWeight\endcsname{\href{https://github.com/lsernau/grh-lean/blob/72fead102f663d5c7a407964a13dc0cf9d977207/grh-compact/src/New/Rows/LineCoordinates.lean\#L117-L118}{{\leanfont Data.\allowbreak{}carrierChar\_hasWeight}}}
\expandafter\newcommand\csname LN@bigChar\endcsname{\href{https://github.com/lsernau/grh-lean/blob/72fead102f663d5c7a407964a13dc0cf9d977207/grh-compact/src/New/Rows/LineCoordinates.lean\#L121-L121}{{\leanfont Data.\allowbreak{}bigChar}}}
\expandafter\newcommand\csname LN@MemberIndex\endcsname{\href{https://github.com/lsernau/grh-lean/blob/72fead102f663d5c7a407964a13dc0cf9d977207/grh-compact/src/New/Rows/LineCoordinates.lean\#L125-L126}{{\leanfont Data.\allowbreak{}MemberIndex}}}
\expandafter\newcommand\csname LN@fNatural\endcsname{\href{https://github.com/lsernau/grh-lean/blob/72fead102f663d5c7a407964a13dc0cf9d977207/grh-compact/src/New/Rows/LineCoordinates.lean\#L130-L134}{{\leanfont Data.\allowbreak{}fNatural}}}
\expandafter\newcommand\csname LN@carddetIndexaddbigCharExp\endcsname{\href{https://github.com/lsernau/grh-lean/blob/72fead102f663d5c7a407964a13dc0cf9d977207/grh-compact/src/New/Rows/LineCoordinates.lean\#L136-L138}{{\leanfont Data.\allowbreak{}card\_detIndex\_add\_bigCharExp}}}
\expandafter\newcommand\csname LN@LPlusIndex\endcsname{\href{https://github.com/lsernau/grh-lean/blob/72fead102f663d5c7a407964a13dc0cf9d977207/grh-compact/src/New/Rows/LineCoordinates.lean\#L141-L141}{{\leanfont Data.\allowbreak{}LPlusIndex}}}
\expandafter\newcommand\csname LN@regroupEquiv\endcsname{\href{https://github.com/lsernau/grh-lean/blob/72fead102f663d5c7a407964a13dc0cf9d977207/grh-compact/src/New/Rows/LineCoordinates.lean\#L145-L146}{{\leanfont Data.\allowbreak{}regroupEquiv}}}
\expandafter\newcommand\csname LN@posEquiv\endcsname{\href{https://github.com/lsernau/grh-lean/blob/72fead102f663d5c7a407964a13dc0cf9d977207/grh-compact/src/New/Rows/LineCoordinates.lean\#L150-L151}{{\leanfont Data.\allowbreak{}posEquiv}}}
\expandafter\newcommand\csname LN@detRows\endcsname{\href{https://github.com/lsernau/grh-lean/blob/72fead102f663d5c7a407964a13dc0cf9d977207/grh-compact/src/New/Rows/LineCoordinates.lean\#L155-L158}{{\leanfont Data.\allowbreak{}detRows}}}
\expandafter\newcommand\csname LN@linePolyCoords\endcsname{\href{https://github.com/lsernau/grh-lean/blob/72fead102f663d5c7a407964a13dc0cf9d977207/grh-compact/src/New/Rows/LineCoordinates.lean\#L163-L164}{{\leanfont Data.\allowbreak{}linePolyCoords}}}
\expandafter\newcommand\csname LN@forgetTwist\endcsname{\href{https://github.com/lsernau/grh-lean/blob/72fead102f663d5c7a407964a13dc0cf9d977207/grh-compact/src/New/Rows/LineCoordinates.lean\#L168-L168}{{\leanfont Data.\allowbreak{}forgetTwist}}}
\expandafter\newcommand\csname LN@lPlusDatum\endcsname{\href{https://github.com/lsernau/grh-lean/blob/72fead102f663d5c7a407964a13dc0cf9d977207/grh-compact/src/New/Rows/LineCoordinates.lean\#L172-L174}{{\leanfont Data.\allowbreak{}lPlusDatum}}}
\expandafter\newcommand\csname LN@lineMultiplier\endcsname{\href{https://github.com/lsernau/grh-lean/blob/72fead102f663d5c7a407964a13dc0cf9d977207/grh-compact/src/New/Rows/LineCoordinates.lean\#L179-L183}{{\leanfont Data.\allowbreak{}lineMultiplier}}}
\expandafter\newcommand\csname LN@carddetIndex\endcsname{\href{https://github.com/lsernau/grh-lean/blob/72fead102f663d5c7a407964a13dc0cf9d977207/grh-compact/src/New/Rows/LineCoordinates.lean\#L189-L192}{{\leanfont Data.\allowbreak{}card\_detIndex}}}
\expandafter\newcommand\csname LN@weylNormSqlinePolyCoords\endcsname{\href{https://github.com/lsernau/grh-lean/blob/72fead102f663d5c7a407964a13dc0cf9d977207/grh-compact/src/New/Rows/LineCoordinates.lean\#L201-L238}{{\leanfont Data.\allowbreak{}weylNormSq\_linePolyCoords\_le}}}
\expandafter\newcommand\csname LN@toPolylinePolyCoords\endcsname{\href{https://github.com/lsernau/grh-lean/blob/72fead102f663d5c7a407964a13dc0cf9d977207/grh-compact/src/New/Rows/LineCoordinates.lean\#L240-L243}{{\leanfont Data.\allowbreak{}toPoly\_linePolyCoords}}}
\expandafter\newcommand\csname LN@existsdenomlinePolyCoords\endcsname{\href{https://github.com/lsernau/grh-lean/blob/72fead102f663d5c7a407964a13dc0cf9d977207/grh-compact/src/New/Rows/LineCoordinates.lean\#L248-L272}{{\leanfont Data.\allowbreak{}exists\_denom\_linePolyCoords}}}
\expandafter\newcommand\csname LN@normmulgProd\endcsname{\href{https://github.com/lsernau/grh-lean/blob/72fead102f663d5c7a407964a13dc0cf9d977207/grh-compact/src/New/Rows/LineArithmetic.lean\#L20-L23}{{\leanfont Data.\allowbreak{}norm\_mul\_gProd\_pow\_mem\_modulus}}}
\expandafter\newcommand\csname LN@lineMultiplierintegral\endcsname{\href{https://github.com/lsernau/grh-lean/blob/72fead102f663d5c7a407964a13dc0cf9d977207/grh-compact/src/New/Rows/LineArithmetic.lean\#L30-L105}{{\leanfont Data.\allowbreak{}lineMultiplier\_integral\_and\_small}}}
\newcommand{\ConcordanceRows}{\crow{onelemaxOpNorm}{\href{https://github.com/lsernau/grh-lean/blob/72fead102f663d5c7a407964a13dc0cf9d977207/grh-compact/src/New/Rows/SparseThetaSum.lean\#L21-L22}{{\leanfont one\_le\_maxOpNorm}}}{SparseThetaSum.lean}{21--22}
\crow{integralmulholMinors}{\href{https://github.com/lsernau/grh-lean/blob/72fead102f663d5c7a407964a13dc0cf9d977207/grh-compact/src/New/Rows/SparseThetaSum.lean\#L30-L56}{{\leanfont integral\_mul\_holMinorsPi\_le}}}{SparseThetaSum.lean}{30--56}
\crow{isolation}{\href{https://github.com/lsernau/grh-lean/blob/72fead102f663d5c7a407964a13dc0cf9d977207/grh-compact/src/New/Rows/SparseThetaSum.lean\#L62-L84}{{\leanfont isolation}}}{SparseThetaSum.lean}{62--84}
\crow{firstMinimummulle}{\href{https://github.com/lsernau/grh-lean/blob/72fead102f663d5c7a407964a13dc0cf9d977207/grh-compact/src/New/Rows/SparseThetaSum.lean\#L99-L112}{{\leanfont firstMinimum\_mul\_le\_hom}}}{SparseThetaSum.lean}{99--112}
\crow{dataInnertraceChar}{\href{https://github.com/lsernau/grh-lean/blob/72fead102f663d5c7a407964a13dc0cf9d977207/grh-compact/src/New/Rows/SparseThetaSum.lean\#L118-L128}{{\leanfont dataInner\_traceChar}}}{SparseThetaSum.lean}{118--128}
\crow{normmodKernelsub}{\href{https://github.com/lsernau/grh-lean/blob/72fead102f663d5c7a407964a13dc0cf9d977207/grh-compact/src/New/Rows/SparseThetaSum.lean\#L133-L171}{{\leanfont norm\_modKernel\_sub\_realInner\_le}}}{SparseThetaSum.lean}{133--171}
\crow{carrierFormspec}{\href{https://github.com/lsernau/grh-lean/blob/72fead102f663d5c7a407964a13dc0cf9d977207/grh-compact/src/New/Rows/NaturalFormNear.lean\#L23-L35}{{\leanfont carrierForm\_spec}}}{NaturalFormNear.lean}{23--35}
\crow{carrierFormnezero}{\href{https://github.com/lsernau/grh-lean/blob/72fead102f663d5c7a407964a13dc0cf9d977207/grh-compact/src/New/Rows/NaturalFormNear.lean\#L39-L50}{{\leanfont carrierForm\_ne\_zero}}}{NaturalFormNear.lean}{39--50}
\crow{negIndexposRowIndexspec}{\href{https://github.com/lsernau/grh-lean/blob/72fead102f663d5c7a407964a13dc0cf9d977207/grh-compact/src/New/Rows/NaturalFormNear.lean\#L54-L58}{{\leanfont Data.\allowbreak{}negIndex\_posRowIndex\_spec}}}{NaturalFormNear.lean}{54--58}
\crow{existsfNaturalapprox}{\href{https://github.com/lsernau/grh-lean/blob/72fead102f663d5c7a407964a13dc0cf9d977207/grh-compact/src/New/Rows/NaturalFormNear.lean\#L65-L146}{{\leanfont Data.\allowbreak{}exists\_fNatural\_approx}}}{NaturalFormNear.lean}{65--146}
\crow{cmSignstar}{\href{https://github.com/lsernau/grh-lean/blob/72fead102f663d5c7a407964a13dc0cf9d977207/grh-compact/src/New/Rows/Separation.lean\#L29-L30}{{\leanfont cmSign\_star}}}{Separation.lean}{29--30}
\crow{isRowTypeauxPair}{\href{https://github.com/lsernau/grh-lean/blob/72fead102f663d5c7a407964a13dc0cf9d977207/grh-compact/src/New/Rows/Separation.lean\#L35-L41}{{\leanfont isRowType\_auxPair}}}{Separation.lean}{35--41}
\crow{eSign}{\href{https://github.com/lsernau/grh-lean/blob/72fead102f663d5c7a407964a13dc0cf9d977207/grh-compact/src/New/Rows/Separation.lean\#L44-L44}{{\leanfont Data.\allowbreak{}eSign}}}{Separation.lean}{44--44}
\crow{heckeCount}{\href{https://github.com/lsernau/grh-lean/blob/72fead102f663d5c7a407964a13dc0cf9d977207/grh-compact/src/New/Rows/Separation.lean\#L47-L47}{{\leanfont Data.\allowbreak{}heckeCount}}}{Separation.lean}{47--47}
\crow{linePair}{\href{https://github.com/lsernau/grh-lean/blob/72fead102f663d5c7a407964a13dc0cf9d977207/grh-compact/src/New/Rows/Separation.lean\#L51-L59}{{\leanfont Data.\allowbreak{}linePair}}}{Separation.lean}{51--59}
\crow{heckeValueProd}{\href{https://github.com/lsernau/grh-lean/blob/72fead102f663d5c7a407964a13dc0cf9d977207/grh-compact/src/New/Rows/Separation.lean\#L63-L65}{{\leanfont Data.\allowbreak{}heckeValueProd}}}{Separation.lean}{63--65}
\crow{existsisFlat}{\href{https://github.com/lsernau/grh-lean/blob/72fead102f663d5c7a407964a13dc0cf9d977207/grh-compact/src/New/Rows/Separation.lean\#L73-L78}{{\leanfont Data.\allowbreak{}exists\_isFlat}}}{Separation.lean}{73--78}
\crow{heightLenormtrivialLift}{\href{https://github.com/lsernau/grh-lean/blob/72fead102f663d5c7a407964a13dc0cf9d977207/grh-compact/src/New/Rows/Separation.lean\#L83-L95}{{\leanfont Data.\allowbreak{}heightLe\_norm\_trivialLift\_sq}}}{Separation.lean}{83--95}
\crow{cosetKernelSumeqfNatural}{\href{https://github.com/lsernau/grh-lean/blob/72fead102f663d5c7a407964a13dc0cf9d977207/grh-compact/src/New/Rows/Separation.lean\#L107-L157}{{\leanfont Data.\allowbreak{}cosetKernelSum\_eq\_fNatural}}}{Separation.lean}{107--157}
\crow{goodFormsrapid}{\href{https://github.com/lsernau/grh-lean/blob/72fead102f663d5c7a407964a13dc0cf9d977207/grh-compact/src/New/Rows/Separation.lean\#L163-L165}{{\leanfont Data.\allowbreak{}goodForms\_mem\_rapidOfLevel}}}{Separation.lean}{163--165}
\crow{mulfNaturalrapid}{\href{https://github.com/lsernau/grh-lean/blob/72fead102f663d5c7a407964a13dc0cf9d977207/grh-compact/src/New/Rows/Separation.lean\#L169-L182}{{\leanfont Data.\allowbreak{}mul\_fNatural\_mem\_rapidOfLevel\_and\_hol}}}{Separation.lean}{169--182}
\crow{tableTheta}{\href{https://github.com/lsernau/grh-lean/blob/72fead102f663d5c7a407964a13dc0cf9d977207/grh-compact/src/New/Rows/Separation.lean\#L190-L193}{{\leanfont Data.\allowbreak{}tableTheta}}}{Separation.lean}{190--193}
\crow{eqzeroofintegral}{\href{https://github.com/lsernau/grh-lean/blob/72fead102f663d5c7a407964a13dc0cf9d977207/grh-compact/src/New/Rows/Separation.lean\#L199-L236}{{\leanfont Data.\allowbreak{}eq\_zero\_of\_integral\_fNatural\_tableTheta}}}{Separation.lean}{199--236}
\crow{existstwistedclassSum}{\href{https://github.com/lsernau/grh-lean/blob/72fead102f663d5c7a407964a13dc0cf9d977207/grh-compact/src/New/Rows/HeckeValueHeight.lean\#L28-L45}{{\leanfont exists\_twisted\_classSum\_ge}}}{HeckeValueHeight.lean}{28--45}
\crow{existsfactorisation}{\href{https://github.com/lsernau/grh-lean/blob/72fead102f663d5c7a407964a13dc0cf9d977207/grh-compact/src/New/Rows/HeckeValueHeight.lean\#L52-L56}{{\leanfont Data.\allowbreak{}exists\_factorisation\_le}}}{HeckeValueHeight.lean}{52--56}
\crow{xiSum}{\href{https://github.com/lsernau/grh-lean/blob/72fead102f663d5c7a407964a13dc0cf9d977207/grh-compact/src/New/Rows/HeckeValueHeight.lean\#L60-L62}{{\leanfont Data.\allowbreak{}Pair.\allowbreak{}xiSum}}}{HeckeValueHeight.lean}{60--62}
\crow{xiSummem}{\href{https://github.com/lsernau/grh-lean/blob/72fead102f663d5c7a407964a13dc0cf9d977207/grh-compact/src/New/Rows/HeckeValueHeight.lean\#L64-L69}{{\leanfont Data.\allowbreak{}Pair.\allowbreak{}xiSum\_mem}}}{HeckeValueHeight.lean}{64--69}
\crow{smulapplyxiSum}{\href{https://github.com/lsernau/grh-lean/blob/72fead102f663d5c7a407964a13dc0cf9d977207/grh-compact/src/New/Rows/HeckeValueHeight.lean\#L71-L82}{{\leanfont Data.\allowbreak{}Pair.\allowbreak{}smul\_apply\_xiSum}}}{HeckeValueHeight.lean}{71--82}
\crow{existstestform}{\href{https://github.com/lsernau/grh-lean/blob/72fead102f663d5c7a407964a13dc0cf9d977207/grh-compact/src/New/Rows/HeckeValueHeight.lean\#L87-L148}{{\leanfont Data.\allowbreak{}exists\_test\_form\_inner\_kernelForm}}}{HeckeValueHeight.lean}{87--148}
\crow{heightLeinvratio}{\href{https://github.com/lsernau/grh-lean/blob/72fead102f663d5c7a407964a13dc0cf9d977207/grh-compact/src/New/Rows/HeckeValueHeight.lean\#L154-L216}{{\leanfont Data.\allowbreak{}heightLe\_inv\_ratio}}}{HeckeValueHeight.lean}{154--216}
\crow{heightLeinvheckeValueProd}{\href{https://github.com/lsernau/grh-lean/blob/72fead102f663d5c7a407964a13dc0cf9d977207/grh-compact/src/New/Rows/HeckeValueHeight.lean\#L223-L237}{{\leanfont Data.\allowbreak{}heightLe\_inv\_heckeValueProd\_div}}}{HeckeValueHeight.lean}{223--237}
\crow{compressedGramOp}{\href{https://github.com/lsernau/grh-lean/blob/72fead102f663d5c7a407964a13dc0cf9d977207/grh-compact/src/New/Rows/LiftVolumeHeight.lean\#L26-L27}{{\leanfont Data.\allowbreak{}compressedGramOp}}}{LiftVolumeHeight.lean}{26--27}
\crow{adjointcomp}{\href{https://github.com/lsernau/grh-lean/blob/72fead102f663d5c7a407964a13dc0cf9d977207/grh-compact/src/New/Rows/LiftVolumeHeight.lean\#L31-L32}{{\leanfont Data.\allowbreak{}adjoint\_comp\_eq\_conj\_smul}}}{LiftVolumeHeight.lean}{31--32}
\crow{detcompressedGramOp}{\href{https://github.com/lsernau/grh-lean/blob/72fead102f663d5c7a407964a13dc0cf9d977207/grh-compact/src/New/Rows/LiftVolumeHeight.lean\#L35-L37}{{\leanfont Data.\allowbreak{}det\_compressedGramOp}}}{LiftVolumeHeight.lean}{35--37}
\crow{normdetpadic}{\href{https://github.com/lsernau/grh-lean/blob/72fead102f663d5c7a407964a13dc0cf9d977207/grh-compact/src/New/Rows/LiftVolumeHeight.lean\#L47-L70}{{\leanfont Data.\allowbreak{}norm\_det\_compressedGramOp\_padic\_le}}}{LiftVolumeHeight.lean}{47--70}
\crow{normdetarch}{\href{https://github.com/lsernau/grh-lean/blob/72fead102f663d5c7a407964a13dc0cf9d977207/grh-compact/src/New/Rows/LiftVolumeHeight.lean\#L76-L95}{{\leanfont Data.\allowbreak{}norm\_det\_compressedGramOp\_arch\_le}}}{LiftVolumeHeight.lean}{76--95}
\crow{heightLeliftVolumeratio}{\href{https://github.com/lsernau/grh-lean/blob/72fead102f663d5c7a407964a13dc0cf9d977207/grh-compact/src/New/Rows/LiftVolumeHeight.lean\#L103-L127}{{\leanfont Data.\allowbreak{}heightLe\_liftVolume\_ratio}}}{LiftVolumeHeight.lean}{103--127}
\crow{untwistedModulus}{\href{https://github.com/lsernau/grh-lean/blob/72fead102f663d5c7a407964a13dc0cf9d977207/grh-compact/src/New/Rows/LineCoordinates.lean\#L24-L24}{{\leanfont Data.\allowbreak{}untwistedModulus}}}{LineCoordinates.lean}{24--24}
\crow{idealuntwistedModulus}{\href{https://github.com/lsernau/grh-lean/blob/72fead102f663d5c7a407964a13dc0cf9d977207/grh-compact/src/New/Rows/LineCoordinates.lean\#L26-L27}{{\leanfont Data.\allowbreak{}ideal\_untwistedModulus\_eq\_span}}}{LineCoordinates.lean}{26--27}
\crow{CarrierIndex}{\href{https://github.com/lsernau/grh-lean/blob/72fead102f663d5c7a407964a13dc0cf9d977207/grh-compact/src/New/Rows/LineCoordinates.lean\#L33-L33}{{\leanfont Data.\allowbreak{}CarrierIndex}}}{LineCoordinates.lean}{33--33}
\crow{FlatIndex}{\href{https://github.com/lsernau/grh-lean/blob/72fead102f663d5c7a407964a13dc0cf9d977207/grh-compact/src/New/Rows/LineCoordinates.lean\#L36-L36}{{\leanfont Data.\allowbreak{}FlatIndex}}}{LineCoordinates.lean}{36--36}
\crow{SmallIndex}{\href{https://github.com/lsernau/grh-lean/blob/72fead102f663d5c7a407964a13dc0cf9d977207/grh-compact/src/New/Rows/LineCoordinates.lean\#L39-L39}{{\leanfont Data.\allowbreak{}SmallIndex}}}{LineCoordinates.lean}{39--39}
\crow{SmallClasses}{\href{https://github.com/lsernau/grh-lean/blob/72fead102f663d5c7a407964a13dc0cf9d977207/grh-compact/src/New/Rows/LineCoordinates.lean\#L43-L44}{{\leanfont Data.\allowbreak{}SmallClasses}}}{LineCoordinates.lean}{43--44}
\crow{basisVec}{\href{https://github.com/lsernau/grh-lean/blob/72fead102f663d5c7a407964a13dc0cf9d977207/grh-compact/src/New/Rows/LineCoordinates.lean\#L49-L49}{{\leanfont Data.\allowbreak{}basisVec}}}{LineCoordinates.lean}{49--49}
\crow{LineChoice}{\href{https://github.com/lsernau/grh-lean/blob/72fead102f663d5c7a407964a13dc0cf9d977207/grh-compact/src/New/Rows/LineCoordinates.lean\#L55-L56}{{\leanfont Data.\allowbreak{}LineChoice}}}{LineCoordinates.lean}{55--56}
\crow{shortRows}{\href{https://github.com/lsernau/grh-lean/blob/72fead102f663d5c7a407964a13dc0cf9d977207/grh-compact/src/New/Rows/LineCoordinates.lean\#L60-L61}{{\leanfont Data.\allowbreak{}shortRows}}}{LineCoordinates.lean}{60--61}
\crow{twistedChi}{\href{https://github.com/lsernau/grh-lean/blob/72fead102f663d5c7a407964a13dc0cf9d977207/grh-compact/src/New/Rows/LineCoordinates.lean\#L67-L67}{{\leanfont Data.\allowbreak{}twistedChi}}}{LineCoordinates.lean}{67--67}
\crow{twistedChihasWeight}{\href{https://github.com/lsernau/grh-lean/blob/72fead102f663d5c7a407964a13dc0cf9d977207/grh-compact/src/New/Rows/LineCoordinates.lean\#L69-L70}{{\leanfont Data.\allowbreak{}twistedChi\_hasWeight}}}{LineCoordinates.lean}{69--70}
\crow{PadIndex}{\href{https://github.com/lsernau/grh-lean/blob/72fead102f663d5c7a407964a13dc0cf9d977207/grh-compact/src/New/Rows/LineCoordinates.lean\#L74-L74}{{\leanfont Data.\allowbreak{}PadIndex}}}{LineCoordinates.lean}{74--74}
\crow{BigIndex}{\href{https://github.com/lsernau/grh-lean/blob/72fead102f663d5c7a407964a13dc0cf9d977207/grh-compact/src/New/Rows/LineCoordinates.lean\#L77-L77}{{\leanfont Data.\allowbreak{}BigIndex}}}{LineCoordinates.lean}{77--77}
\crow{padSigns}{\href{https://github.com/lsernau/grh-lean/blob/72fead102f663d5c7a407964a13dc0cf9d977207/grh-compact/src/New/Rows/LineCoordinates.lean\#L81-L81}{{\leanfont Data.\allowbreak{}padSigns}}}{LineCoordinates.lean}{81--81}
\crow{negIndex}{\href{https://github.com/lsernau/grh-lean/blob/72fead102f663d5c7a407964a13dc0cf9d977207/grh-compact/src/New/Rows/LineCoordinates.lean\#L85-L86}{{\leanfont Data.\allowbreak{}negIndex}}}{LineCoordinates.lean}{85--86}
\crow{posRowIndex}{\href{https://github.com/lsernau/grh-lean/blob/72fead102f663d5c7a407964a13dc0cf9d977207/grh-compact/src/New/Rows/LineCoordinates.lean\#L90-L91}{{\leanfont Data.\allowbreak{}posRowIndex}}}{LineCoordinates.lean}{90--91}
\crow{untwistedAuxPair}{\href{https://github.com/lsernau/grh-lean/blob/72fead102f663d5c7a407964a13dc0cf9d977207/grh-compact/src/New/Rows/LineCoordinates.lean\#L94-L94}{{\leanfont Data.\allowbreak{}untwistedAuxPair}}}{LineCoordinates.lean}{94--94}
\crow{DetIndex}{\href{https://github.com/lsernau/grh-lean/blob/72fead102f663d5c7a407964a13dc0cf9d977207/grh-compact/src/New/Rows/LineCoordinates.lean\#L98-L99}{{\leanfont Data.\allowbreak{}DetIndex}}}{LineCoordinates.lean}{98--99}
\crow{linePoly}{\href{https://github.com/lsernau/grh-lean/blob/72fead102f663d5c7a407964a13dc0cf9d977207/grh-compact/src/New/Rows/LineCoordinates.lean\#L104-L106}{{\leanfont Data.\allowbreak{}linePoly}}}{LineCoordinates.lean}{104--106}
\crow{bigCharExp}{\href{https://github.com/lsernau/grh-lean/blob/72fead102f663d5c7a407964a13dc0cf9d977207/grh-compact/src/New/Rows/LineCoordinates.lean\#L111-L111}{{\leanfont Data.\allowbreak{}bigCharExp}}}{LineCoordinates.lean}{111--111}
\crow{carrierChar}{\href{https://github.com/lsernau/grh-lean/blob/72fead102f663d5c7a407964a13dc0cf9d977207/grh-compact/src/New/Rows/LineCoordinates.lean\#L115-L115}{{\leanfont Data.\allowbreak{}carrierChar}}}{LineCoordinates.lean}{115--115}
\crow{carrierCharhasWeight}{\href{https://github.com/lsernau/grh-lean/blob/72fead102f663d5c7a407964a13dc0cf9d977207/grh-compact/src/New/Rows/LineCoordinates.lean\#L117-L118}{{\leanfont Data.\allowbreak{}carrierChar\_hasWeight}}}{LineCoordinates.lean}{117--118}
\crow{bigChar}{\href{https://github.com/lsernau/grh-lean/blob/72fead102f663d5c7a407964a13dc0cf9d977207/grh-compact/src/New/Rows/LineCoordinates.lean\#L121-L121}{{\leanfont Data.\allowbreak{}bigChar}}}{LineCoordinates.lean}{121--121}
\crow{MemberIndex}{\href{https://github.com/lsernau/grh-lean/blob/72fead102f663d5c7a407964a13dc0cf9d977207/grh-compact/src/New/Rows/LineCoordinates.lean\#L125-L126}{{\leanfont Data.\allowbreak{}MemberIndex}}}{LineCoordinates.lean}{125--126}
\crow{fNatural}{\href{https://github.com/lsernau/grh-lean/blob/72fead102f663d5c7a407964a13dc0cf9d977207/grh-compact/src/New/Rows/LineCoordinates.lean\#L130-L134}{{\leanfont Data.\allowbreak{}fNatural}}}{LineCoordinates.lean}{130--134}
\crow{carddetIndexaddbigCharExp}{\href{https://github.com/lsernau/grh-lean/blob/72fead102f663d5c7a407964a13dc0cf9d977207/grh-compact/src/New/Rows/LineCoordinates.lean\#L136-L138}{{\leanfont Data.\allowbreak{}card\_detIndex\_add\_bigCharExp}}}{LineCoordinates.lean}{136--138}
\crow{LPlusIndex}{\href{https://github.com/lsernau/grh-lean/blob/72fead102f663d5c7a407964a13dc0cf9d977207/grh-compact/src/New/Rows/LineCoordinates.lean\#L141-L141}{{\leanfont Data.\allowbreak{}LPlusIndex}}}{LineCoordinates.lean}{141--141}
\crow{regroupEquiv}{\href{https://github.com/lsernau/grh-lean/blob/72fead102f663d5c7a407964a13dc0cf9d977207/grh-compact/src/New/Rows/LineCoordinates.lean\#L145-L146}{{\leanfont Data.\allowbreak{}regroupEquiv}}}{LineCoordinates.lean}{145--146}
\crow{posEquiv}{\href{https://github.com/lsernau/grh-lean/blob/72fead102f663d5c7a407964a13dc0cf9d977207/grh-compact/src/New/Rows/LineCoordinates.lean\#L150-L151}{{\leanfont Data.\allowbreak{}posEquiv}}}{LineCoordinates.lean}{150--151}
\crow{detRows}{\href{https://github.com/lsernau/grh-lean/blob/72fead102f663d5c7a407964a13dc0cf9d977207/grh-compact/src/New/Rows/LineCoordinates.lean\#L155-L158}{{\leanfont Data.\allowbreak{}detRows}}}{LineCoordinates.lean}{155--158}
\crow{linePolyCoords}{\href{https://github.com/lsernau/grh-lean/blob/72fead102f663d5c7a407964a13dc0cf9d977207/grh-compact/src/New/Rows/LineCoordinates.lean\#L163-L164}{{\leanfont Data.\allowbreak{}linePolyCoords}}}{LineCoordinates.lean}{163--164}
\crow{forgetTwist}{\href{https://github.com/lsernau/grh-lean/blob/72fead102f663d5c7a407964a13dc0cf9d977207/grh-compact/src/New/Rows/LineCoordinates.lean\#L168-L168}{{\leanfont Data.\allowbreak{}forgetTwist}}}{LineCoordinates.lean}{168--168}
\crow{lPlusDatum}{\href{https://github.com/lsernau/grh-lean/blob/72fead102f663d5c7a407964a13dc0cf9d977207/grh-compact/src/New/Rows/LineCoordinates.lean\#L172-L174}{{\leanfont Data.\allowbreak{}lPlusDatum}}}{LineCoordinates.lean}{172--174}
\crow{lineMultiplier}{\href{https://github.com/lsernau/grh-lean/blob/72fead102f663d5c7a407964a13dc0cf9d977207/grh-compact/src/New/Rows/LineCoordinates.lean\#L179-L183}{{\leanfont Data.\allowbreak{}lineMultiplier}}}{LineCoordinates.lean}{179--183}
\crow{carddetIndex}{\href{https://github.com/lsernau/grh-lean/blob/72fead102f663d5c7a407964a13dc0cf9d977207/grh-compact/src/New/Rows/LineCoordinates.lean\#L189-L192}{{\leanfont Data.\allowbreak{}card\_detIndex}}}{LineCoordinates.lean}{189--192}
\crow{weylNormSqlinePolyCoords}{\href{https://github.com/lsernau/grh-lean/blob/72fead102f663d5c7a407964a13dc0cf9d977207/grh-compact/src/New/Rows/LineCoordinates.lean\#L201-L238}{{\leanfont Data.\allowbreak{}weylNormSq\_linePolyCoords\_le}}}{LineCoordinates.lean}{201--238}
\crow{toPolylinePolyCoords}{\href{https://github.com/lsernau/grh-lean/blob/72fead102f663d5c7a407964a13dc0cf9d977207/grh-compact/src/New/Rows/LineCoordinates.lean\#L240-L243}{{\leanfont Data.\allowbreak{}toPoly\_linePolyCoords}}}{LineCoordinates.lean}{240--243}
\crow{existsdenomlinePolyCoords}{\href{https://github.com/lsernau/grh-lean/blob/72fead102f663d5c7a407964a13dc0cf9d977207/grh-compact/src/New/Rows/LineCoordinates.lean\#L248-L272}{{\leanfont Data.\allowbreak{}exists\_denom\_linePolyCoords}}}{LineCoordinates.lean}{248--272}
\crow{normmulgProd}{\href{https://github.com/lsernau/grh-lean/blob/72fead102f663d5c7a407964a13dc0cf9d977207/grh-compact/src/New/Rows/LineArithmetic.lean\#L20-L23}{{\leanfont Data.\allowbreak{}norm\_mul\_gProd\_pow\_mem\_modulus}}}{LineArithmetic.lean}{20--23}
\crow{lineMultiplierintegral}{\href{https://github.com/lsernau/grh-lean/blob/72fead102f663d5c7a407964a13dc0cf9d977207/grh-compact/src/New/Rows/LineArithmetic.lean\#L30-L105}{{\leanfont Data.\allowbreak{}lineMultiplier\_integral\_and\_small}}}{LineArithmetic.lean}{30--105}}
% the notations of the defined functions: one slot for each explicit argument of the Lean declaration
\newcommand{\eSign}[2]{\mathsf{e}_{#1}(#2)}
\newcommand{\heckeCount}[1]{\mathsf{e}^{\ast}(#1)}
\newcommand{\linePair}[2]{\operatorname{lp}_D(#1,#2)}
\newcommand{\heckeValueProd}[2]{\operatorname{HV}_D(#1,#2)}
\newcommand{\tableTheta}[3]{\operatorname{tt}_D(#1,#2)(#3)}
\newcommand{\xiSum}[3]{\operatorname{xs}_{#1,#2}(#3)}
\newcommand{\compressedGramOp}[3]{\operatorname{cg}_{#1}(#2,#3)}
\newcommand{\untwistedModulus}[1]{\mathfrak{m}^{\circ}_{#1}}
\newcommand{\CarrierIndex}[1]{\mathsf{C}_{#1}}
\newcommand{\FlatIndex}[1]{\mathsf{F}_{#1}}
\newcommand{\SmallIndex}[1]{\mathsf{S}_{#1}}
\newcommand{\SmallClasses}[1]{\operatorname{Cl}^{\circ}(#1)}
\newcommand{\basisVec}[2]{\operatorname{bv}_{#1}(#2)}
\newcommand{\LineChoice}[1]{\operatorname{LC}(#1)}
\newcommand{\shortRows}[2]{\operatorname{sr}_{#1}(#2)}
\newcommand{\twistedChi}[2]{\chi_D^{(#2)}[#1]}
\newcommand{\PadIndex}[1]{\mathsf{P}_D(#1)}
\newcommand{\BigIndex}[1]{\mathsf{B}_D(#1)}
\newcommand{\padSigns}[1]{\operatorname{ps}_D(#1)}
\newcommand{\posRowIndex}[3]{\operatorname{pos}_{D,#1}(#2,#3)}
\newcommand{\untwistedAuxPair}[1]{O^{\circ}_D(#1)}
\newcommand{\bigCharExp}[1]{\mathsf{x}_D(#1)}
\newcommand{\bigChar}[2]{\mathbf{b}''_D(#1,#2)}
\newcommand{\MemberIndex}[1]{\operatorname{Mem}_D(#1)}
\newcommand{\fNatural}[3]{F^{\natural}_{D}(#1,#2)(#3)}
\newcommand{\LPlusIndex}[1]{\mathsf{L}^{+}_D(#1)}
\newcommand{\regroupEquiv}[1]{\operatorname{rg}_D(#1)}
\newcommand{\posEquiv}[1]{\operatorname{pe}_D(#1)}
\newcommand{\linePolyCoords}[3]{\operatorname{lc}_D(#1,#2,#3)}
\newcommand{\forgetTwist}[2]{\operatorname{fg}_{#1}(#2)}
\newcommand{\lPlusDatum}[4]{\operatorname{ld}_D(#1,#2;#3)(#4)}
\newcommand{\lineMultiplier}[5]{\operatorname{lm}_D(#1,#2,#3,#4)_{#5}}
\newcommand{\UnitMembers}{219}
\newcommand{\UnitInNote}{77}
\newcommand{\UnitProofs}{85}
\newcommand{\UnitSteps}{634}
\newcommand{\UnitAllSteps}{680}


\linespread{1.08}                       % 11 pt on 14.7 pt
\setlength{\parindent}{0pt}
\setlength{\parskip}{5.5pt plus 1.5pt minus 0.5pt}
\setlength{\abovedisplayskip}{7.5pt plus 2pt minus 3pt}
\setlength{\belowdisplayskip}{7.5pt plus 2pt minus 3pt}
\setlength{\abovedisplayshortskip}{3pt plus 2pt}
\setlength{\belowdisplayshortskip}{6pt plus 2pt minus 2pt}
\emergencystretch=1.5em
\clubpenalty=10000 \widowpenalty=10000 \displaywidowpenalty=10000
\relpenalty=7000 \binoppenalty=9000
\raggedbottom
\pagestyle{fancy}\fancyhf{}\renewcommand{\headrulewidth}{0pt}
\fancyfoot[C]{\sffamily\fontsize{7.6}{9}\selectfont\color{inkfaint}\thepage\,/\,\pageref*{LastPage}}
\setlist{topsep=2pt,partopsep=0pt,parsep=0pt,itemsep=4pt,leftmargin=2.1em,labelsep=0.55em}

\newcommand{\Q}{\mathbb{Q}}
\newcommand{\Z}{\mathbb{Z}}
\newcommand{\F}{\mathbb{F}}
\newcommand{\len}{\operatorname{len}}
\newcommand{\cok}{\operatorname{cok}}
\newcommand{\End}{\operatorname{End}}
\newcommand{\Hom}{\operatorname{Hom}}
\newcommand{\Gal}{\operatorname{Gal}}
\newcommand{\diag}{\operatorname{diag}}
\newcommand{\kw}[1]{\emph{#1}}
\newcommand{\m}[1]{\mbox{$#1$}}        % a short formula that must not be broken across lines
\newcommand{\head}[1]{\par\addvspace{17pt plus 4pt minus 2pt}\Needspace{6\baselineskip}%
  \noindent{\fontsize{13.2}{16}\selectfont\bfseries\color{ink}#1}\par\nopagebreak\vspace{-1pt}}
\newcommand{\sect}[2]{\refstepcounter{section}\label{#2}\par\addvspace{19pt plus 4pt minus 2pt}\Needspace{7\baselineskip}%
  \noindent\llap{{\fontsize{13.2}{16}\selectfont\bfseries\color{inkfaint}\thesection}\hspace{9pt}}%
  {\fontsize{13.2}{16}\selectfont\bfseries\color{ink}#1}\par\nopagebreak\vspace{-1pt}}
\newcommand{\subhead}[1]{\par\addvspace{9pt plus 2pt}\Needspace{4\baselineskip}\noindent\textbf{#1}\enspace\ignorespaces}
% where a statement lives in the Lean archive
\newcommand{\leantag}{{\sffamily\bfseries\fontsize{6.3}{8}\selectfont\color{leancol}\addfontfeatures{LetterSpace=7}LEAN}}
\newcommand{\sep}{\unskip\nobreak\hspace{5pt}{\color{inkfaint}\textperiodcentered}\hspace{5pt plus 2pt}\ignorespaces}
\newcommand{\leanline}[1]{\par\nopagebreak\vspace{-2.5pt}{\fontsize{8.9}{12.4}\selectfont\color{inksoft}\raggedright
  \noindent\llap{\leantag\hspace{9pt}}#1\par}\vspace{1pt}}
% a displayed statement in a tinted box: #1 colour, #2 text
\newcommand{\statement}[2]{\par\addvspace{7pt}\begingroup\setlength{\fboxsep}{8.5pt}\setlength{\fboxrule}{0.6pt}%
  \noindent\fcolorbox{#1!30!white}{#1!3!white}{\parbox{\dimexpr\linewidth-2\fboxsep-2\fboxrule\relax}{#2}}%
  \endgroup\par\addvspace{7pt}}
\newcommand{\nt}[2]{\par\noindent\hangindent=1.1em\hangafter=1 #1\enspace{\color{inksoft}#2}\par\vspace{1pt}}
\newcommand{\legdot}[1]{\tikz[baseline=-0.62ex]\fill[#1] (0,0) circle (2.5pt);\hspace{3pt}}
\newcommand{\crow}[4]{\csname CS@#1\endcsname & #2 & \textcolor{inksoft}{#3} & \textcolor{inksoft}{#4}\\}
\newcommand{\leanquoteline}[2]{\par\noindent\hangindent=\dimexpr#1\fontdimen2\font+2.2em\relax\hangafter=1
  \hspace*{#1\fontdimen2\font}#2}
\newcommand{\quoted}[1]{\par\vspace{3pt}{\ttfamily\fontsize{7.5}{10.4}\selectfont\color{leancol}\raggedright
  \setlength{\parskip}{0pt}\leanquote{#1}\par}\vspace{4pt}}
\newcommand{\bx}[3]{{\sffamily\bfseries\fontsize{7.3}{9.2}\selectfont\color{#1!82!black}#2}\\[1.5pt]#3}
\newcommand{\cs}[2]{\expandafter\newcommand\csname CS@#1\endcsname{#2}}

% theorems
\newtheoremstyle{plainx}{7pt plus 2pt}{7pt plus 2pt}{\itshape}{}{\bfseries}{.}{0.5em}{}
\newtheoremstyle{uprightx}{7pt plus 2pt}{7pt plus 2pt}{}{}{\bfseries}{.}{0.5em}{}
\theoremstyle{plainx}
\newtheorem{theorem}{Theorem}[section]
\newtheorem{lemma}[theorem]{Lemma}
\newtheorem{proposition}[theorem]{Proposition}
\newtheorem{corollary}[theorem]{Corollary}
\theoremstyle{uprightx}
\newtheorem{definition}[theorem]{Definition}
\newtheorem{setting}[theorem]{Setting}          % hypotheses shared by several statements; each of them cites it by number
\newtheorem{remark}[theorem]{Remark}
\newtheorem{example}[theorem]{Example}
\renewcommand{\qedsymbol}{{\setlength{\fboxsep}{0pt}\setlength{\fboxrule}{0.5pt}\fbox{\rule{0pt}{0.46em}\hspace{0.46em}}}}
% what holds throughout the note: the ambient objects only (implicit and instance arguments), never a hypothesis
\newenvironment{conventions}{\par\addvspace{5pt}\noindent\textbf{Conventions.}\enspace\ignorespaces}{\par\addvspace{3pt}}
\makeatletter
\renewenvironment{proof}[1][Proof]{\par\pushQED{\qed}\normalfont\topsep0pt\partopsep0pt
  \trivlist\item[\hskip\labelsep\itshape #1.]\ignorespaces}{\popQED\endtrivlist\@endpefalse}
\makeatother

% where each declaration of the concordance is treated: \cs{key}{place}, the place naming a \label or a lettered item
\cs{onelemaxOpNorm}{Lemma~\ref{l:mh}}
\cs{integralmulholMinors}{Lemma~\ref{l:mh}}
\cs{isolation}{Proposition~\ref{p:iso}}
\cs{firstMinimummulle}{Lemma~\ref{l:sparse}}
\cs{dataInnertraceChar}{Lemma~\ref{l:sparse}}
\cs{normmodKernelsub}{Lemma~\ref{l:sparse}}
\cs{carrierFormspec}{Lemma~\ref{l:carrier}}
\cs{carrierFormnezero}{Lemma~\ref{l:carrier}}
\cs{negIndexposRowIndexspec}{Lemma~\ref{l:counts}}
\cs{existsfNaturalapprox}{Proposition~\ref{p:approx}}
\cs{cmSignstar}{Lemma~\ref{l:counts}}
\cs{isRowTypeauxPair}{Lemma~\ref{l:counts}}
\cs{eSign}{Definition~\ref{d:hv}}
\cs{heckeCount}{Definition~\ref{d:hv}}
\cs{linePair}{Definition~\ref{d:hv}}
\cs{heckeValueProd}{Definition~\ref{d:hv}}
\cs{existsisFlat}{Lemma~\ref{l:flat}}
\cs{heightLenormtrivialLift}{Lemma~\ref{l:flat}}
\cs{cosetKernelSumeqfNatural}{Theorem C}
\cs{goodFormsrapid}{Lemma~\ref{l:rapid}}
\cs{mulfNaturalrapid}{Lemma~\ref{l:rapid}}
\cs{tableTheta}{Definition~\ref{d:hv}}
\cs{eqzeroofintegral}{Theorem D}
\cs{existstwistedclassSum}{Lemma~\ref{l:pigeon}}
\cs{existsfactorisation}{Lemma~\ref{l:pigeon}}
\cs{xiSum}{Definition~\ref{d:aux}}
\cs{xiSummem}{Lemma~\ref{l:test}}
\cs{smulapplyxiSum}{Lemma~\ref{l:test}}
\cs{existstestform}{Lemma~\ref{l:test}}
\cs{heightLeinvratio}{Theorem B}
\cs{heightLeinvheckeValueProd}{Corollary~\ref{c:hv}}
\cs{compressedGramOp}{Definition~\ref{d:aux}}
\cs{adjointcomp}{Lemma~\ref{l:gram}}
\cs{detcompressedGramOp}{Lemma~\ref{l:gram}}
\cs{normdetpadic}{Proposition~\ref{p:det}}
\cs{normdetarch}{Proposition~\ref{p:det}}
\cs{heightLeliftVolumeratio}{Theorem A}
\cs{untwistedModulus}{Definition~\ref{d:index}}
\cs{idealuntwistedModulus}{Definition~\ref{d:chars}}
\cs{CarrierIndex}{Definition~\ref{d:index}}
\cs{FlatIndex}{Definition~\ref{d:index}}
\cs{SmallIndex}{Definition~\ref{d:index}}
\cs{SmallClasses}{Definition~\ref{d:index}}
\cs{basisVec}{Definition~\ref{d:index}}
\cs{LineChoice}{Definition~\ref{d:line}}
\cs{shortRows}{Definition~\ref{d:line}}
\cs{twistedChi}{Definition~\ref{d:chars}}
\cs{twistedChihasWeight}{Definition~\ref{d:chars}}
\cs{PadIndex}{Definition~\ref{d:index}}
\cs{BigIndex}{Definition~\ref{d:index}}
\cs{padSigns}{Definition~\ref{d:index}}
\cs{negIndex}{Definition~\ref{d:index}}
\cs{posRowIndex}{Definition~\ref{d:index}}
\cs{untwistedAuxPair}{Definition~\ref{d:index}}
\cs{DetIndex}{Definition~\ref{d:line}}
\cs{linePoly}{Definition~\ref{d:line}}
\cs{bigCharExp}{Definition~\ref{d:chars}}
\cs{carrierChar}{Definition~\ref{d:chars}}
\cs{carrierCharhasWeight}{Definition~\ref{d:chars}}
\cs{bigChar}{Definition~\ref{d:chars}}
\cs{MemberIndex}{Definition~\ref{d:line}}
\cs{fNatural}{Definition~\ref{d:line}}
\cs{carddetIndexaddbigCharExp}{Lemma~\ref{l:counts}}
\cs{LPlusIndex}{Definition~\ref{d:index}}
\cs{regroupEquiv}{Definition~\ref{d:index}}
\cs{posEquiv}{Definition~\ref{d:index}}
\cs{detRows}{Definition~\ref{d:line}}
\cs{linePolyCoords}{Definition~\ref{d:line}}
\cs{forgetTwist}{Definition~\ref{d:index}}
\cs{lPlusDatum}{Definition~\ref{d:line}}
\cs{lineMultiplier}{Definition~\ref{d:line}}
\cs{carddetIndex}{Lemma~\ref{l:counts}}
\cs{weylNormSqlinePolyCoords}{Lemma~\ref{l:coords}}
\cs{toPolylinePolyCoords}{Lemma~\ref{l:coords}}
\cs{existsdenomlinePolyCoords}{Lemma~\ref{l:coords}}
\cs{normmulgProd}{Proposition~\ref{p:lineint}}
\cs{lineMultiplierintegral}{Proposition~\ref{p:lineint}}
\newcommand{\CC}{\mathbb{C}}
\newcommand{\ZZ}{\mathbb{Z}}
\newcommand{\NN}{\mathbb{N}}
\newcommand{\RR}{\mathbb{R}}
\newcommand{\AAA}{\mathbb{A}}
\newcommand{\ix}[1]{\mathsf{#1}}
\newcommand{\OEr}{\mathcal{O}_E}
\newcommand{\UU}{\mathrm{U}}
\newcommand{\GL}{\mathrm{GL}}
\newcommand{\KK}[1]{\mathbf{K}_{#1}}
\newcommand{\ip}[2]{\langle #1,#2\rangle}
\newcommand{\clsum}[2]{\sideset{}{^{\mathrm{cl}}}\sum_{#1\bmod #2}}
\newcommand{\tk}[7]{\theta^{#1,#2}_{#3,#4;#5}[#6,#7]}
\newcommand{\AutC}{\operatorname{Aut}(\CC)}
\newcommand{\transp}{^{\mathsf{T}}}
\newcommand{\inputbox}[1]{\statement{inksoft}{#1}}
\newcommand{\LP}[3]{\Pi_{#1}(#2,#3)}
% objects of the sister notes and of other units, written by hand with all their arguments
\newcommand{\FormVx}[1]{\mathcal{F}(#1)}
\newcommand{\gau}[5]{\operatorname{gauge}^{#1}_{#2}(#3;#4,#5)}
\newcommand{\nsq}[2]{\|#1\|^{2}_{#2}}
\newcommand{\cks}[5]{\operatorname{cks}_{#1,#2}(#3,#4;#5)}
\newcommand{\Mgood}[1]{\mathcal{M}^{\circ}_D(W_{#1}(D))}
\newcommand{\alift}{\operatorname{lift}^{*}}
\newcommand{\lvol}[2]{\operatorname{vol}_D(#1,#2)}
\newcommand{\carrier}[2]{\mathfrak{C}_{#1,#2}}
\newcommand{\cXi}[3]{\Xi^{\circ}_{#1,#2}(#3)}
\newcommand{\vf}{\operatorname{vf}}
\newcommand{\rat}[3]{\rho_{#1}(#2,#3)}      % Lean's Pair.ratio p ϑ s: the pair, the tame character, the sign
\newcommand{\idem}{\operatorname{idem}}
\newcommand{\clU}{h_{\UU}}                    % the class sum of the constant 1 over U(1) at the chosen level
\newcommand{\tm}{t_{\mathbf{m}}}              % the number of the depth of an indicator jet
\newcommand{\rowpar}[2]{\operatorname{rp}_D(#1,#2)}
\newcommand{\acct}{ \textup{(}an account of the formal proof\textup{)}}
\newcommand{\triv}[2]{\Xi_{#1}(#2)}
\newcommand{\ialg}[2]{I_{#1}(#2)}
% definitions of this note whose arguments the text reading misreads; written with the arguments their uses show
\newcommand{\negIx}[2]{\operatorname{neg}_{D,#1}(#2)}
\newcommand{\DetIx}[2]{\operatorname{Det}_D(#1,#2)}
\newcommand{\linePoly}[1]{\Lambda_D(#1)}
\newcommand{\carrierChar}[1]{\mathbf{c}_D(#1)}
\newcommand{\detRows}[4]{\eta_{D}(#1,#2,#3,#4)}

\begin{document}
\color{ink}

{\raggedright
{\sffamily\bfseries\fontsize{7.2}{9}\selectfont\color{leancol}\addfontfeatures{LetterSpace=8}READBACK NOTE · CHAPTER VII · UNIT 143 · PART 4 OF 4\par}
\vspace{5pt}
{\fontsize{21.5}{25.5}\selectfont\bfseries The height of a ratio of lift volumes: twisted theta forms on one line and
Hecke values\par}
\vspace{5pt}
{\fontsize{12.4}{16}\selectfont\itshape\color{inksoft} One unit of the Lean artifact \ArtifactName, read back into
mathematics\par}}
\vspace{8pt}
{\color{rulegrey}\hrule height 0.6pt}
\vspace{7pt}
\thispagestyle{empty}

\textbf{Summary.} For a fit family of twists, the volume of a single theta lift at the trivial twist divided by
its volume at the twist has height at most the dimension of the rows times
$k(D)\cdot(\log b\cdot B_0(D)-\mathsf{U}(D))$ (Theorem~A): each lift repays the cost \m{\mathsf{U}(D)} of one direction, $k(D)$ times, for every row. The proof
has three parts. First, an explicit theta form \m{F^{\natural}} on the indefinite group of the $j$-th lift is, for
every twist, a constant multiple of a class sum of twisted theta kernels, and the constant has absolute value a
product of Hecke $L$-values: the twisted forms lie on one line (Theorem~C). Second, the products of
\m{F^{\natural}} with theta kernels of a complementary definite block detect every good form (Theorem~D); with the
comparison forms of the sister note this turns the volume of a lift into a determinant with integral entries
times a ratio of Hecke values. Third, the inverse of each ratio of a twisted to an untwisted Hecke value has height
at most \m{6B_0(D)-\mathsf{U}(D)} (Theorem~B): its valuations are bounded by test forms and the congruence bound of a
fit family, and its size by the norm of an untwisted theta lift.

\statement{leancol}{\fontsize{9.6}{13.2}\selectfont\textbf{What this note is.} The artifact \ArtifactName{} is a
body of Lean code: the compact edition of a formal proof, due to L.~Alpöge with the AI system Claude, of the Riemann
hypothesis and of its generalisation to Dirichlet $L$-functions. That headline is a claim. The authors report that
every file compiles with Lean v4.35.0 against a pinned version of Mathlib and that Lean's kernel accepted the
proofs; an independent build is recorded in the handout's record of checks when it reports. No mathematician has
read the argument through, and nothing here says that the Riemann hypothesis is proved. This note is the last of
four on one unit of the artifact, the folder of the rows; it treats \UnitInNote{} of the unit's \UnitMembers{}
declarations, those of seven files. Every declaration is stated. Theorem~A is proved in full from Theorem~B with
Corollary~\ref{c:hv}, from Lemma~\ref{l:gram} and Proposition~\ref{p:det}, from the comparison forms of the sister
note and from the definition of a fit family; the proofs of Corollary~\ref{c:hv}, of Lemmas \ref{l:pigeon}
and~\ref{l:gram}, of Proposition~\ref{p:iso} and of Proposition~\ref{p:det}\,(b) are complete as well. All other
proofs are accounts of the formal proofs: they say which results combine and do not verify the estimates and the
normalising constants. Each of them is headed as an account where it stands, and the proof of
Proposition~\ref{p:approx} is an outline. Theorems B, C and~D and Proposition~\ref{p:det}\,(a), on which Theorem~A
rests, are among the accounts. The inputs from other units and from the background library are named and not
proved here. Every statement was compared with its Lean text, argument by argument (for the definitions of the
line see the last section).
The Lean text read is \PinText; Appendix~A is the concordance of
\ConcDecls{} declarations.}

\sect{The statements}{s:intro}
\statement{arith}{\phantomsection\label{t:A}\textbf{Theorem A} (the height of a ratio of lift volumes).\enspace
\emph{In Setting~\ref{set:data}, let $P$ be parameters, $D$ admissible data for~$P$, \m{\theta=(\theta_{\ell,\iota})} a
family of $m$-fold twists of~$D$ that is fit for~$D$, and $j$ a lift index. Then the family
\[(\ell,\iota)\longmapsto\frac{\lvol{1}{j}}{\lvol{\theta_{\ell,\iota}}{j}}\]
has height at most \m{\dim\operatorname{Row}(D)\cdot\bigl(k(D)\cdot(\log b\cdot B_0(D)-\mathsf{U}(D))\bigr)}.}

\vspace{4pt}\phantomsection\label{t:B}\textbf{Theorem B} (each ratio of Hecke values repays one direction).\enspace
\emph{In Setting~\ref{set:data}, let \m{\theta} be a family of $m$-fold twists of~$D$ that is fit for~$D$, $j$ a lift
index and $p$ a pair of~$D$ that is of rank one for~$j$. Then the family
\m{(\ell,\iota)\mapsto\rat{p}{\theta_{\ell,\iota}^{\,j}}{1}^{-1}} of inverse ratios for the sign~$1$ has height at most
\m{6B_0(D)-\mathsf{U}(D)}.}}
\leanline{\lean{heightLeliftVolumeratio}\sep \lean{heightLeinvratio}}

\statement{arith}{\phantomsection\label{t:C}\textbf{Theorem C} (twisted theta forms lie on one line).\enspace
\emph{In Setting~\ref{set:data}, let \m{\ix a} be a finite index set and $j$ a lift index. Then there are a real
number~$C$ and a function~$c$ from the twists of~$D$ to~$\CC$ such that for every twist~$\vartheta$:
\m{|c(\vartheta)|^2=C/\heckeValueProd{j}{\vartheta}}, and for every matrix~$u$ with columns~$\ix a$ over
\m{\OEr/\bar\pi_g^{2}}, every choice \m{i\in\MemberIndex{j}}, every point \m{g\in\UU(W_j(D))(\AAA)} and every
matrix~$w$ with rows~$\ix a$ and columns \m{\BigIndex{j}} over \m{\OEr/\bar\pi_g^{2}} with \m{uw=i_1}:
\[\fNatural{j}{i}{g}=c(\vartheta)\cdot\cks{O_D(\vartheta,j)}{\bigChar{\vartheta}{j}}{w}{\phi}{\lineMultiplier{\vartheta}{j}{i}{\cdot}{\cdot}}(g),\]
with $\phi$ the inclusion of the row indices in \m{\BigIndex{j}}.}

\vspace{4pt}\phantomsection\label{t:D}\textbf{Theorem D} (products of theta forms detect all good forms).\enspace
\emph{In Setting~\ref{set:data}, let $j$ be a lift index and \m{F\in\Mgood{j}} a good form such that
\[\int_{[\UU(W_j(D))]}\overline{F(g)\,\fNatural{j}{i}{g}}\cdot\tableTheta{j}{\zeta}{g}\,dg=0\]
for every choice \m{i\in\MemberIndex{j}} and every function~$\zeta$ from \m{\SmallClasses{D}} to coordinate vectors of
weight~$\kappa^{-}$. Then \m{F=0}.}}
\leanline{\lean{cosetKernelSumeqfNatural}\sep \lean{eqzeroofintegral}}

Here \m{\lvol{\vartheta}{j}} is the volume of the $j$-th single lift at the twist~$\vartheta^{j}$ (the note \emph{The
volume of the chain of theta lifts as a product of lift volumes}), heights are those of Chapter~IV,
\m{B_0(D)=\operatorname{budget}(D)/(m\,k(D))} is the budget for congruences and \m{\mathsf{U}(D)} the cost of one direction
(Setting~\ref{set:data}); \m{\rat{p}{\vartheta}{s}} is the ratio of a twisted to an untwisted Hecke $L$-value of the
pair~$p$, for the tame character~$\vartheta$ and the sign~$s$ (Setting~\ref{set:data}); the function
\m{\fNatural{j}{i}{\cdot}}, the line multiplier and the product \m{\heckeValueProd{j}{\vartheta}} of Hecke values are
defined in §\ref{s:setting}. In Theorem~C the left side does not depend on the twist. The lettered theorems are
those of this note; ``Theorem~C of the handout'' is the theorem of the overview. Since
\m{h_{\mathrm{main}}(D)=m\,k(D)\,\mathsf{U}(D)}, the sum over \m{j=1,\dots,m} of the bounds of Theorem~A, which is the
bound for the product of the $m$ families, is
\m{\dim\operatorname{Row}(D)\cdot(\log b\cdot\operatorname{budget}(D)-h_{\mathrm{main}}(D))}: with the factorisation of the
volume of the chain in Theorem~A of the first note of this unit, this gives the second height bound in the proof of
Theorem~C of the handout.

\sect{The setting}{s:setting}
\begin{conventions}\phantomsection\label{c:conventions}
Finite index sets are written in sans-serif, as in the other notes of this unit; \m{|\ix n|} is the number of
elements. The ambient objects of the declarations are the parameters~$P$ and the data~$D$ in the statements about
data; a number field or CM field~$E$ with a CM type and Weil data, finite non-empty index sets and sign vectors in
the general statements; a prime~$\ell$ and an embedding into~$\CC_\ell$ where gauges occur; and finite-dimensional
inner product spaces in §\ref{s:volume}. Objects that the formal text introduces as implicit arguments of a
statement are introduced in the statement.
\end{conventions}

\begin{setting}[the data and what other units and the sister notes attach to it]\label{set:data}
The general notions are those of the background Setting of the note \emph{Auxiliary dual pairs, determinant
polynomials and the carrier}: a CM field~$E$ with CM type and Weil data, the character~$\upsilon$, moduli, characters with weight and
tame part, tameness, the partial $L$-functions \m{L^{T}(z,\nu)} (of the base change of~$\nu$, without the Euler
factors at the places of~$T$ and at their complex conjugates) and the products \m{\LP{T}{\nu}{d}} of their values,
hermitian spaces, class sums, theta kernels
\m{\tk{S}{T}{\alpha}{\beta}{\mathfrak{m}}{\Psi}{P}(g,h)}, weights, coordinate vectors with their polynomials
\m{\operatorname{pol}_\kappa(\xi)} and Gram norms, auxiliary pairs~$O$ with their forms \m{\FormVx{O}}, and the carrier
\m{\carrier{E}{m}} with the class sum \m{\cXi{m}{c}{h}}. Gauges, integrality at an embedding, the numbers
\m{\nsq{\phi}{\sigma}}, ``tested against a subspace for an idempotent'' and the coset kernel sum
\m{\cks{O}{b}{w}{e}{y}(g)} are those of the note \emph{Glued
lattices, product forms and the integral comparison forms}.

\emph{Data.} $P=(N_0,d)$ are parameters, with \m{K=\mathbb{Q}(\sqrt{-d})}, and $D$ is admissible data for~$P$: the
field \m{E(D)} with its infinite places~$\tau$, the scales $x$, $x_1$, $b$, the number~$m$ of lifts, an even
number, the ranks \m{k(D)} and~$r$ with \m{r+1=k(D)+m}, the step \m{t(D)=5r}, the number \m{f(D)} of infinite places,
the level \m{g(D)=g_0g_1} and the integer \m{\bar\pi_g} with \m{\bar\pi_g\overline{\bar\pi_g}=g(D)}, whose prime
divisors are the chosen primes, the nominal weights \m{w_D(\tau)}, the prime sums \m{\operatorname{Ask}_D(z,\tau)}, the
weights \m{w_j(\tau)} of the lifts for \m{0\le j\le m}, the gap \m{\mathsf{g}(D)} and the reduced half gap
\m{\mathsf{g}'(D)}, the fixed characters \m{\chi_D}, of weight~$0$ and tame part~$\omega_D$, and \m{\mu^{\flat}_D}, and
\m{\mu^{0}_j=\mu^{\flat}_D(\upsilon^{-\mathsf{g}(D)}\chi_D)^{j}}. Further
\begin{align*}
B_0(D)&=\frac{\operatorname{budget}(D)}{m\,k(D)},\qquad\operatorname{js}_D(n)=\log n!-n\log\frac{\Omega}{g(D)},\\
\mathsf{U}(D)&=\sum_\tau\Bigl(\operatorname{js}_D(w_0(\tau))-\operatorname{Ask}_D(w_D(\tau),\tau)\Bigr),
\end{align*}
with $\Omega$ the period constant of \m{E(D)}; \m{\mathsf{U}(D)} is the \kw{cost of one direction} and
\m{\operatorname{js}_D} the \kw{jet size}. A lift index is an integer~$j$ with \m{1\le j\le m}. Tame characters and
$m$-fold twists with their sets \m{Q(\vartheta)}, heights, the good idempotent \m{e_0(D)} and the idempotents
\m{\idem(\mathbf{c})} of the tuples~$\mathbf{c}$ of labels, the spaces $W_j(D)$, the good forms \m{\Mgood{j}}, the rows,
the auxiliary pairs \m{O_D(\vartheta,j)}, with index sets \m{\ix r} of size~$r$ and \m{\ix n} of size \m{r+t(D)}, the
adjoint lifts \m{\alift_D(\vartheta,j)} and the lift volumes \m{\lvol{\vartheta}{j}} are as in the first note of this
unit. A \kw{twist} is a tame character whose tame part is non-trivial at every prime of \m{Q(\vartheta)} (the Setting
on the data of the note \emph{Glued lattices, product forms and the integral comparison forms}); the powers
\m{\vartheta^{j}}, \m{1\le j\le m}, of an $m$-fold twist are twists. For a tame character~$\vartheta$ and an integer
\m{c\ge0}, \m{\mathfrak{m}_D(\vartheta,c)} is the modulus \m{\bar\pi_g^{c}} times the primes of \m{Q(\vartheta)}.

\emph{Pairs.} A \kw{pair} of~$D$ (unit \emph{Exchange}) is a character~$\nu_p$ of weight~$-c_p$, with tame
part~$\psi_p$, together with an integer~$z_p$. It is \kw{flat} for a finite index set~$\ix n$ if \m{c_{p,\tau}=z_p}
for all~$\tau$ and \m{|\ix n|=2z_p\ge4}; \kw{of rank one} for a lift index~$j$ if \m{m<z_p\le r}, its tame part
is~$\omega_D^{j}$ and \m{w_j(\tau)+r-t(D)\le c_{p,\tau}\le w_j(\tau)}; \kw{ramified} if its tame part is non-trivial
exactly at the chosen primes. A pair with a suitable index set and a twist~$\vartheta$ has an \kw{$L$-value datum},
with an auxiliary pair of a definite group \m{\UU(1)} and a definite space, jet degrees \m{c_{p,\tau}-z_p}, kernel
forms, and a \kw{trivial lift} \m{\triv{p}{\vartheta}}, the theta lift of the trivial character of \m{\UU(1)}; its
\kw{normalised value} \m{\ialg{p}{\vartheta}} is the squared norm of the trivial lift cleared of its denominators. The
labels of such an auxiliary pair have degree~$1$, and \m{\idem(X-1)}, the \kw{label idempotent}, is the idempotent of
its tuple of labels all equal to \m{X-1}. For a pair~$p$, a tame character~$\vartheta$ and a sign~$s$, the \kw{ratio}
of the twisted to the untwisted Hecke $L$-value of the pair is
\[\rat{p}{\vartheta}{s}=\bigl|L^{T_\vartheta}\bigl(z_p,(\nu_p\vartheta^{s})^{-1}\bigr)\bigr|\Big/\bigl|L^{T_1}\bigl(z_p,\nu_p^{-1}\bigr)\bigr|,\]
with \m{T_\vartheta} the set of the chosen primes and of the primes of \m{Q(\vartheta)}, so that $T_1$ is the set of the
chosen primes: the sign is the exponent of the twist. The ratio is real and not zero, and
\m{\rat{p}{\vartheta}{1}=\ialg{p}{\vartheta}/\ialg{p}{1}} for an index set that fits~$p$ (unit \emph{Exchange}).

\emph{Denominators.} Let $O$ be an auxiliary pair whose first space has rank~$r_0$, $\mathcal{Z}$ a linear subspace
of \m{\FormVx{O}}, $e$ an idempotent of the group algebra of the finite group of~$O$, \m{(\ell,\iota)} a prime with a
field isomorphism from~$\CC$ to an algebraic closure of~$\mathbb{Q}_\ell$, and $\beta$ a real number. The
subspace~$\mathcal{Z}$ has \kw{denominators at most~$\beta$ at~$\iota$ for~$e$} if, for all families \m{(y_a)},
\m{(y'_b)} of \m{\dim\mathcal{Z}} forms, the~$y_a$ integral at every embedding \m{\CC\to\CC_\ell} and the~$y'_b$
integral at~$\iota$ and tested against~\m{\mathcal{Z}^{\perp}} for~$e$, the determinant of the numbers
\m{\ip{y_a}{\operatorname{pr}_{\mathcal{Z}}y'_b}} has \m{|\iota(\,\cdot\,)|_\ell\le\exp(\dim\mathcal{Z}\cdot r_0\,\beta)}; here
\m{\operatorname{pr}_{\mathcal{Z}}} is the orthogonal projection to~$\mathcal{Z}$.

\emph{Fitness.} A family~$\theta$ of $m$-fold twists is \kw{fit for~$D$} if: (1) for every \m{(\ell,\iota)}, the
prime~$\ell$ is prime to the norms of the primes in \m{Q(\theta_{\ell,\iota})}; (2) for every lift index~$j$ there are
real numbers~$\beta_\ell$, whose sum over every finite set of primes is at most \m{B_0(D)}, such that for every
\m{(\ell,\iota)} the image of the adjoint lift for \m{\theta_{\ell,\iota}^{\,j}} on the good forms has denominators at
most~$\beta_\ell$ at~$\iota$ for the good idempotent \m{e_0(D)}, and, for every pair~$p$ of rank one for~$j$, the line
\m{\CC\,\triv{p}{\theta_{\ell,\iota}^{\,j}}} has denominators at most~$\beta_\ell$ at~$\iota$ for the label idempotent
\m{\idem(X-1)}; (3) for every lift index~$j$, every ramified pair~$p$ and every sign~$s$, the family
\m{(\ell,\iota)\mapsto\rat{p}{\theta_{\ell,\iota}^{\,j}}{s}^{-1}} has height at most~$x^{999}$. In~(3) the sign stands in
the third slot of the ratio: the twist enters the numerator as \m{(\theta_{\ell,\iota}^{\,j})^{s}}.
\end{setting}

\subsection*{Definitions: index sets and characters}
\begin{definition}[index sets of the $j$-th lift]\label{d:index}
In Setting~\ref{set:data}, let $j$ be a lift index.
\begin{enumerate}[label=\textup{(\arabic*)},leftmargin=2.3em,itemsep=2pt]
\item \m{\untwistedModulus{D}} is the function \m{c\mapsto\mathfrak{m}_D(1,c)}, the modulus~\m{\bar\pi_g^{c}}; for a tame
character~$\vartheta$ and \m{x\in\OEr/\mathfrak{m}_D(\vartheta,2)}, \m{\forgetTwist{\vartheta}{x}} is its reduction in
\m{\OEr/\bar\pi_g^{2}}; \m{\untwistedAuxPair{j}=O_D(1,j)} is the auxiliary pair at the trivial twist, with weights
\m{\kappa^{+}}, \m{\kappa^{-}}.
\item \m{\CarrierIndex{D}} is the disjoint union of \m{\{1,\dots,2m\}} and \m{\{1,\dots,m\}}; \m{\FlatIndex{D}} is
\m{\{1,\dots,6m+1\}}; \m{\SmallIndex{D}} is the disjoint union of the row index set~$\ix r$ and a set of
\m{t(D)-9m-1-r} elements; \m{\SmallClasses{D}} is the set of classes of \m{\UU(I_{\SmallIndex{D}})} at the principal
level modulo~\m{\bar\pi_g^{2}}.
\item \m{\PadIndex{j}} is the disjoint union of \m{\{1,\dots,j\}} and a set of \m{3m-j} elements;
\m{\BigIndex{j}} is the disjoint union of~$\ix r$, \m{\PadIndex{j}} and \m{\FlatIndex{D}}; \m{\LPlusIndex{j}} is the
disjoint union of a set of \m{3m-j} elements, \m{\FlatIndex{D}} and \m{\CarrierIndex{D}}. The sign vector
\m{\padSigns{j}} on \m{\ix r\sqcup\PadIndex{j}} has the signs of~$W_j(D)$ on~$\ix r$, $-1$ on the first $j$ padding
indices and $+1$ on the others.
\item \m{\basisVec{D}{i}}, for an integer \m{i\ge0}, is the row index with number \m{\min(i,r-1)} in the
numbering of~$\ix r$. For \m{0\le i<m}, the index \m{\negIx{j}{i}} of \m{\ix r\sqcup\PadIndex{j}} is the $i$-th of
the first padding indices if \m{i<j}, and \m{\basisVec{D}{k(D)+i-1}} otherwise. For integers \m{h\ge0} and
\m{0\le i<2m}, \m{\posRowIndex{j}{h}{i}} is \m{\basisVec{D}{h+i}} if \m{h+i<k(D)+j-1}, and the $i$-th of the other
padding indices if not.
\item \m{\regroupEquiv{j}} is the bijection that regroups the summands, and \m{\posEquiv{j}} is a chosen bijection:
\begin{align*}
\regroupEquiv{j}&\colon\ \ix r\sqcup\{1,\dots,j\}\sqcup\LPlusIndex{j}\ \longrightarrow\ \CarrierIndex{D}\sqcup\BigIndex{j},\\
\posEquiv{j}&\colon\ \{1,\dots,2m\}\sqcup\BigIndex{j}\ \longrightarrow\ \{1,\dots,k(D)+j-1\}\sqcup\LPlusIndex{j}.
\end{align*}
\end{enumerate}
\end{definition}
\leanline{\lean{untwistedModulus}\sep \lean{forgetTwist}\sep \lean{untwistedAuxPair}\sep \lean{CarrierIndex}\sep
\lean{FlatIndex}\sep \lean{SmallIndex}\sep \lean{SmallClasses}\sep \lean{PadIndex}\sep \lean{BigIndex}\sep
\lean{LPlusIndex}\sep \lean{padSigns}\sep \lean{basisVec}\sep \lean{negIndex}\sep \lean{posRowIndex}\sep
\lean{regroupEquiv}\sep \lean{posEquiv}\sep the fact that \m{m\ne0}, an unnamed instance of the same file, is used
throughout}

The space of the $j$-th auxiliary pair on the definite side has the index set
\m{\ix n=\BigIndex{j}\sqcup\SmallIndex{D}}: a big block~$V''$, which contains copies of the indices of~$W_j(D)$, of
a padding and of a flat block~$I^{\flat}$ of rank \m{6m+1}, and a small block~$V^{\circ}$. The carrier lives on a
space of signature \m{(2m,m)} with index set \m{\CarrierIndex{D}}.

\begin{definition}[characters]\label{d:chars}
In Setting~\ref{set:data}, let $j$ be a lift index and $\vartheta$ a tame character of~$D$.
\begin{enumerate}[label=\textup{(\roman*)},leftmargin=2.3em,itemsep=2pt]
\item \m{\twistedChi{\vartheta}{j}=\chi_D^{\,j}\,\vartheta}.
\item \m{\bigCharExp{j}=j+\lfloor(t(D)+k(D)-1)/2\rfloor}, and
\m{\bigChar{\vartheta}{j}=\mu^{0}_j\,\vartheta\,\upsilon^{\bigCharExp{j}}}.
\item \m{\carrierChar{j}=\mu^{\flat}_D\,\upsilon^{\,m+\bigCharExp{j}-j\,\mathsf{g}(D)}}.
\end{enumerate}
\end{definition}
\leanline{\lean{twistedChi}\sep \lean{bigCharExp}\sep \lean{bigChar}\sep \lean{carrierChar}\sep
\m{\twistedChi{\vartheta}{j}} has weight~$0$ and tame part \m{\omega_D^{j}} times that of~$\vartheta$,
\lean{twistedChihasWeight}\sep \m{\carrierChar{j}} has weight \m{\tau\mapsto m+\bigCharExp{j}-w_j(\tau)} and trivial tame
part, \lean{carrierCharhasWeight}\sep the ideal of \m{\untwistedModulus{D}(c)} is generated by \m{\bar\pi_g^{c}},
\lean{idealuntwistedModulus}}

\subsection*{Definitions: the line}
\begin{definition}[choices, determinants, the line multiplier, the function F-natural]\label{d:line}
In Setting~\ref{set:data}, let $j$ be a lift index.
\begin{enumerate}[label=\textup{(\arabic*)},leftmargin=2.3em,itemsep=2pt]
\item A \kw{line choice}, an element of \m{\LineChoice{D}}, is a pair \m{i=(i_c,i_l)} of a class~$i_c$ of
\m{\UU(I_{\FlatIndex{D}})} at the principal level modulo~\m{\bar\pi_g^{2}} and, for each infinite place~$\tau$, a family
$i_l$ of \m{\mathsf{g}'(D)} lists of $2m$ indices in \m{\FlatIndex{D}}. For a line choice~$i$, \m{\shortRows{D}{i}} is
the transpose of a chosen invertible matrix over \m{E(D)} whose columns lie in the lattice of the
class~$i_c$ and whose entries are small at every complex embedding (Input~3).
\item A \kw{choice}, an element of \m{\MemberIndex{j}}, is a pair \m{i=(i_1,i_2)} of a \m{\CarrierIndex{D}\times\BigIndex{j}}
matrix~$i_1$ over \m{\OEr/\bar\pi_g^{2}} and a line choice~$i_2$.
\item For an infinite place~$\tau$, \m{\DetIx{j}{\tau}} is the index set with one element \m{(h,c)} for every
\m{h<k(D)+j-1} and every \m{c<\kappa^{+}_\tau(h)}, and \m{\mathsf{g}'(D)} further elements.
\item The \kw{line polynomial} is
\[\linePoly{j}=\prod_\tau\det\bigl(y^{-}_\tau\bigr)^{|\DetIx{j}{\tau}|+|\BigIndex{j}|+m}\cdot\prod_\tau\ \prod_{h<k(D)+j-1}\det\bigl(y_{\tau,h}\bigr)^{\kappa^{+}_\tau(h)},\]
where $y^{-}_\tau$ is the minor of the matrix of variables at~$\tau$ in the rows \m{\negIx{j}{\cdot}} and the negative
columns, and $y_{\tau,h}$ the minor in the rows \m{\posRowIndex{j}{h}{\cdot}} and the positive columns.
\item For a line choice~$i$, a matrix~$Y$ over \m{E(D)}, by default \m{Y=\shortRows{D}{i}}, an infinite
place~$\tau$ and \m{l\in\DetIx{j}{\tau}}, \m{\detRows{j}{i}{Y}{\tau}_l} is the complex matrix with rows
\m{\{1,\dots,2m\}} and columns \m{\{1,\dots,k(D)+j-1\}\sqcup\LPlusIndex{j}} that is, for \m{l=(h,c)}, the transposed
inclusion matrix whose row~$a$ is the basis vector \m{h+a} of the first block if \m{h+a<k(D)+j-1} and the vector~$a$
of the block of \m{3m-j} indices otherwise, and, for the other~$l$, the rows \m{i_l(\tau,l)} of \m{\tau(Y)}, placed
in the block \m{\FlatIndex{D}}.
\item The \kw{line coordinates} \m{\linePolyCoords{j}{i}{Y}} are the coordinates in the first columns, in the
weight~$\kappa^{+}$, of the product of stacked determinants of the matrices \m{\detRows{j}{i}{Y}{\tau}_l} along
\m{\posEquiv{j}} (the sister note on auxiliary pairs, the definition of stacked determinants): for each standard filling~$l$ of
weight~$\kappa^{+}$ with rows \m{\BigIndex{j}}, a polynomial in variables indexed by \m{\LPlusIndex{j}}.
\item For a tame character~$\vartheta$, a matrix~$u$ over \m{\OEr/\bar\pi_g^{2}} and a
\m{\BigIndex{j}\times\LPlusIndex{j}} matrix~$z$ over \m{\OEr/\mathfrak{m}_D(\vartheta,2)},
\m{\lPlusDatum{\vartheta}{j}{u}{z}} is the sum, over the invertible \m{3m\times3m} matrices~$g$ over
\m{\OEr/\bar\pi_g^{2}}, of \m{\psi'(\operatorname{tr}(z'gu))} if the columns of~$z$ outside the last $3m$ are the
standard inclusion, and $0$ otherwise; here $z'$ consists of the last $3m$ columns of~$z$, reduced by
\m{\forgetTwist{\vartheta}{\cdot}}, and $\psi'$ is the additive character of the Weil data on \m{\OEr/\bar\pi_g^{2}}.
\item For a tame character~$\vartheta$, a choice \m{i\in\MemberIndex{j}}, \m{h\in\UU(I_{\BigIndex{j}})(\AAA_f)} and a
standard filling~$l$, the \kw{line multiplier} is
\begin{align*}
\lineMultiplier{\vartheta}{j}{i}{h}{l}&=\tk{I_j^{-}}{I_{\BigIndex{j}}}{\upsilon^{j}}{\beta}{\mathfrak{m}}{\mathbf{1}_{I}}{1}(1,h)\cdot\clsum{c}{\KK{3m}}\cXi{m}{\twistedChi{\vartheta}{j}}{c}\cdot\Theta^{+}_{l}(c,h),\\
\Theta^{+}_{l}(c,h)&=\tk{L^{+}}{I_{\BigIndex{j}}}{\beta}{\beta}{\mathfrak{m}}{\lPlusDatum{\vartheta}{j}{i_1}{\cdot}}{\linePolyCoords{j}{i_2}{Y}_l}\bigl((h_{i},c),h\bigr),
\end{align*}
where \m{\beta=\bigChar{\vartheta}{j}}, \m{\mathfrak{m}=\mathfrak{m}_D(\vartheta,2)}, \m{I_j^{-}} is the negative definite space of rank~$j$, \m{\mathbf{1}_I} the indicator of the inclusion
matrix, \m{L^{+}} the definite space with index set \m{\LPlusIndex{j}}, \m{Y=\shortRows{D}{i_2}}, and \m{(h_{i},c)} the
point of~$L^{+}$ with component~$1$ on the first block, a representative~$h_i$ of the class of~$i_2$ on
\m{I_{\FlatIndex{D}}} and $c$ on the carrier block.
\item For a choice \m{i\in\MemberIndex{j}} and \m{g\in\UU(W_j(D))(\AAA)},
\[\fNatural{j}{i}{g}=\int_{[\UU(S)]}\carrier{E(D)}{m}(x)\cdot\tk{S}{W'}{\upsilon^{m}}{\carrierChar{j}}{\bar\pi_g^{2}}{\psi'_{i_1}}{\linePoly{j}\cdot\Delta_{i_2}}\bigl(x,(g,1,h_{i})\bigr)\,dx,\]
where $S$ is the space of signature \m{(2m,m)} with index set \m{\CarrierIndex{D}}, \m{W'} is the orthogonal sum of
$W_j(D)$, of the padding \m{I_j^{-}\oplus I_{3m-j}} and of \m{I_{\FlatIndex{D}}}, \m{\psi'_{i_1}} is the function
\m{y\mapsto\psi'(\operatorname{tr}(y\,i_1))}, the line polynomial is read in the variables of the first two summands,
and \m{\Delta_{i_2}} is the product of the determinants of the rows \m{i_l(\tau,l)} of \m{\tau(\shortRows{D}{i_2})} times
the block of variables of the third summand.
\end{enumerate}
\end{definition}
\leanline{\lean{LineChoice}\sep \lean{shortRows}\sep \lean{MemberIndex}\sep \lean{DetIndex}\sep \lean{linePoly}\sep
\lean{detRows}\sep \lean{linePolyCoords}\sep \lean{lPlusDatum}\sep \lean{lineMultiplier}\sep \lean{fNatural}}

\begin{definition}[Hecke values and the theta function of a table]\label{d:hv}
In Setting~\ref{set:data}:
\begin{enumerate}[label=\textup{(\roman*)},leftmargin=2.3em,itemsep=2pt]
\item For a sign~$s$, \m{\eSign{D}{s}} is $m$ if \m{s=+1} and \m{m/2} if \m{s=-1}; \m{\heckeCount{D}=7m/2+\lfloor k(D)/2\rfloor}.
\item For a lift index~$j$, the three \kw{line pairs} \m{\linePair{j}{o}}, for $o$ either absent or a sign, are: for
$o$ absent, the character \m{\mu^{0}_j\,\upsilon^{\,n_j}}, \m{n_j=j+\lfloor(t(D)-11m-1)/2\rfloor}, of weight
\m{\tau\mapsto n_j-w_j(\tau)}, with the integer \m{\heckeCount{D}}; for a sign~$s$, the character
\m{\upsilon^{-\eSign{D}{s}}\chi_D^{\,sj}}, of weight \m{-\eSign{D}{s}}, with the integer \m{\eSign{D}{s}}.
\item For a lift index~$j$ and a twist~$\vartheta$, with \m{T_\vartheta} the set of the chosen primes and of the
primes in \m{Q(\vartheta)},
\[\heckeValueProd{j}{\vartheta}=\prod_{o}\LP{T_\vartheta}{(\nu_o\,\vartheta^{\,s_o})^{-1}}{z_o},\]
the product over the three line pairs \m{(\nu_o,z_o)=\linePair{j}{o}}, with \m{s_o=-1} for the sign~$-1$ and
\m{s_o=+1} otherwise.
\item For a lift index~$j$, a function~$\zeta$ from \m{\SmallClasses{D}} to coordinate vectors of
weight~$\kappa^{-}$ with rows \m{\SmallIndex{D}} and \m{g\in\UU(W_j(D))(\AAA)},
\[\tableTheta{j}{\zeta}{g}=\sum_{\mathfrak{c}\in\SmallClasses{D}}\overline{\tk{W_j(D)}{I_{\SmallIndex{D}}}{\mathbf{a}}{\upsilon^{m/2-j}}{\bar\pi_g^{2}}{\mathbf{1}_{I}}{\operatorname{pol}^{-}(\overline{\zeta(\mathfrak{c})})}(g,h_{\mathfrak{c}})},\]
with \m{\mathbf{a}} the first character of \m{\untwistedAuxPair{j}}, \m{h_{\mathfrak{c}}} a representative of the class and
the polynomial of the conjugate vector written in the negative letters.
\end{enumerate}
\end{definition}
\leanline{\lean{eSign}\sep \lean{heckeCount}\sep \lean{linePair}\sep \lean{heckeValueProd}\sep \lean{tableTheta}}

\begin{definition}[two auxiliary constructions]\label{d:aux}
\textup{(i)} In Setting~\ref{set:data}, let $q$ be a pair that is flat for \m{\{1,\dots,m'\}}, $X$ an $L$-value datum
with auxiliary pair~$O_X$, $\vartheta$ a twist and $F$ a function from \m{\UU(I_{m'})(\AAA_f)} to \m{\FormVx{O_X}}. Then
\[\xiSum{q}{\vartheta}{F}=M_D\cdot\clsum{h}{\Gamma}\ \frac{\delta_q}{N(Q(\vartheta))^{m'-1}}\,\bigl\langle1,\delta_q\,\triv{q}{\vartheta}(h)\bigr\rangle\cdot F(h),\]
where $M_D$ is the chosen group order of~$D$, $\delta_q$ the denominator bound of~$q$, \m{N(Q(\vartheta))} the product
of the norms of the primes in \m{Q(\vartheta)}, $\Gamma$ the level of the auxiliary pair of the $L$-value datum
of~$q$ at~$\vartheta$, and the bracket the Gram form with the vector~$1$.


\textup{(ii)} For finite-dimensional inner product spaces $T$ and~$V$, \m{\varkappa\in\CC}, a subspace~$L$ of~$V$
and a linear map \m{\operatorname{gl}\colon T\to V}, the \kw{compressed Gram operator}
\m{\compressedGramOp{\varkappa}{L}{\operatorname{gl}}} is the endomorphism
\m{z\mapsto\varkappa\cdot\operatorname{pr}_L(\operatorname{gl}(\operatorname{gl}^{*}z))} of~$L$, with \m{\operatorname{pr}_L} the
orthogonal projection.
\end{definition}
\leanline{\lean{xiSum}\sep \lean{compressedGramOp}}

\inputbox{\textbf{Inputs from the other notes of this unit.}

\vspace{4pt}\textbf{Input 1} (note \emph{Glued lattices, product forms and the integral comparison forms};
comparison forms).\enspace In Setting~\ref{set:data}, for a lift index~$j$ there are a function~$\gamma$ from the
twists of~$D$ to~$\CC$, a finite-dimensional inner product space \m{\operatorname{Tab}_D(j)} of tables
\m{\zeta=(\zeta_a)}, the index~$a$ a pair of a choice and a class \m{\mathfrak{c}\in\SmallClasses{D}}, a linear map
\m{A\colon\operatorname{Tab}_D(j)\to\Mgood{j}} and linear maps
\m{\operatorname{gl}_\vartheta\colon\operatorname{Tab}_D(j)\to\FormVx{O_D(\vartheta,j)}} such that, with $u_\vartheta$ the restriction
of \m{\alift_D(\vartheta,j)} to \m{\Mgood{j}}: (i) \m{\gamma(\vartheta)\ne0} and
\m{|\gamma(\vartheta)/\gamma(1)|^2=\heckeValueProd{j}{\vartheta}/\heckeValueProd{j}{1}}; (ii)
\m{u_\vartheta^{*}\circ\operatorname{gl}_\vartheta=\gamma(\vartheta)A}; (iii) $A$ is surjective; (iv)
\m{\operatorname{gl}_\vartheta(A^{*}y)} is tested against the orthogonal complement of the image of~$u_\vartheta$ for the
good idempotent \m{e_0(D)}, for every \m{y\in\Mgood{j}}; (v) if $\ell$ is prime to the norms of the primes in \m{Q(\vartheta)}, \m{\varphi\colon\CC\to\CC_\ell}
is an embedding and every \m{\overline{\zeta_a}} has gauge at most~$1$ at the representative of its class for
\m{\varphi\circ\mathrm{conj}}, then \m{\operatorname{gl}_\vartheta(\zeta)} is integral at~$\varphi$; (vi)
\m{\mathsf{h}_D(j)\cdot\nsq{\operatorname{gl}_1(\zeta)}{\sigma}\le\exp(2rB_0(D))\sum_a\|\sigma(\zeta_a)\|^2} for all
\m{\sigma\in\AutC}, with \m{\mathsf{h}_D(j)} the hook number of the weight~$\kappa^{-}$. Its proof uses Theorems C
and~D and Proposition~\ref{p:lineint} of this note; there is no circle, since these do not use it.

\vspace{3pt}\textbf{Input 2} (the same note).\enspace (a) Crude bounds: for finite index sets \m{\ix L}, \m{\ix N} and
an integer \m{\mathsf{d}\ge0} there is \m{c\ne0} such that $c$ times every kernel of the definite pair
\m{(I_{\ix L},I_{\ix N})} over \m{E(D)}, with a homogeneous polynomial of degree~$\mathsf{d}$, a
modulus~$\mathfrak{m}$ containing
\m{q\,g(D)^2} and a datum integral away from~$q$, is integral at the primes prime to~$q$ at points integral above
them, and has conjugates of size at most \m{e^{x_1}} times bounds for the datum and the polynomial, when
\m{q=1}, \m{\mathsf{d}\le g(D)} and the ranks are at most \m{40k(D)}. (b) The depth of an indicator jet: for orders
\m{\mathbf{m}=(\mathbf{m}_\tau)} with \m{w_D(\tau)\le\mathbf{m}_\tau\le w_D(\tau)+x^{42}} there is \m{\tm\in\CC},
\m{\tm\ne0}, with conjugates at most \m{\exp(f(D)mb/k(D))}, such that for every prime~$\ell$ and every embedding
\m{\varphi\colon\CC\to\CC_\ell} some Taylor coefficient of total orders \m{\mathbf{m}_\tau}, of an indicator datum of
\m{(I_1,I_{12})} modulo~\m{\bar\pi_g}, in the Taylor series scaled by \m{\bar\pi_g} times the jet scale, times
\m{\kappa_0^{12}\prod_\tau\mathbf{m}_\tau!/\tm}, has \m{|\varphi(\,\cdot\,)|_\ell\ge\ell^{-F_D(\ell)}}, with \m{F_D(\ell)} the
forced part and \m{\kappa_0} the vacuum scale. (c) The conjugates of the normalised value of a flat pair at the trivial twist are at
most \m{\exp(x^{34})}. (d) For an $L$-value datum and a monomial of the jet degrees there is an integral
coordinate vector whose harmonic polynomial is the monomial times a ratio of factorials, fixed by the matrices
that fix the columns of the monomial. (e) Point forms, their pairing and descent; the vanishing of periods for
forms that are not lifts; the coset kernel sum.

\vspace{3pt}\textbf{Input 3} (notes \emph{Auxiliary dual pairs, determinant polynomials and the carrier} and
\emph{The volume of the chain}).\enspace The explicit matching (its Theorem~A), the matching summed against a form
on a square block, the passage from the stabiliser to the full level, the carrier identity (its Theorem~C), the
equivariance of the product of stacked determinants, the lemma on a determinant with fixed columns, the product
rule for weights and norms, the coordinates in the first columns; the image of the adjoint lift on the good forms
is generated tuple by tuple.}

\inputbox{\textbf{Inputs from other units and from the background library.} Each is used where it is named.

\vspace{4pt}\textbf{Unit \emph{Heights}.}\enspace Heights add in products and powers; a bound may be increased; if
each $y'_{\ell,\iota}$ is $y_{\ell,\iota}$ times an integer power of the norm of the modulus of \m{\theta_{\ell,\iota}}, for a
fit family, a height bound for~$y$ holds for~$y'$ (unit \emph{Exchange}).
\textbf{Unit \emph{Exchange}.}\enspace The pairs \m{O_D(\vartheta,j)} fit their parameters; adjoint lifts are
injective; \m{\rat{p}{\vartheta}{s}} is real and not zero, and \m{\rat{p}{\vartheta}{1}=\ialg{p}{\vartheta}/\ialg{p}{1}}; the
logarithm of the conjugates of the squared norm of a
trivial lift; the pairing of a trivial lift with kernel forms; rank bounds; bounds for the weights of the kernels.
\textbf{Unit \emph{Selmer}.}\enspace The definition of a fit family; for a tested form, the orthogonal
projection to a subspace killed by a Hecke test does not increase the norms \m{\nsq{\cdot}{\sigma}}; the lifts of
eigenforms are killed by the Hecke test; an auxiliary pair has infinitely many split places off its level; the
descent of a form of an auxiliary pair \m{O_D(\vartheta,j)} is a form of principal level on \m{W_j(D)}. The
\emph{rank-one bound}: if the line \m{\CC\,\triv{p}{\vartheta}} of a trivial lift has denominators at most~$\beta$
at~$\iota$ for \m{\idem(X-1)}, the class stabilisers of \m{\UU(1)} at the chosen level are trivial, and $\Xi$,
$\Xi'$ are test forms at~$\ell$ (Lemma~\ref{l:test}\,(c)), then
\m{|\iota(\clU\ip{\Xi}{\operatorname{pr}\Xi'})|_\ell\le e^{\beta}}, with $\operatorname{pr}$ the orthogonal projection to the line and
\m{\clU} the number of classes of \m{\UU(1)} at that level.
\textbf{Units \emph{Data}, \emph{Columns}, \emph{Trace}, \emph{Filling}.}\enspace Comparison of monomials in the
scales; free levels and class numbers; \m{\operatorname{Ask}_D(y,\tau)} as a sum over primes of local terms; equal
dimensions of good forms; the line coordinates as coefficients of the product of determinants in the positive
letters.}

\inputbox{\textbf{Inputs from the background library.} Each is used where it is named. The references are those of
the handout's catalogue, with the relation that its second reading allows.

\vspace{4pt}A cusp form is large only away from the cusps: a non-zero rapidly
decreasing holomorphic automorphic form has a point~$g$ with \m{|F|\le2|F(g)|} everywhere and first minimum bounded
below wherever \m{|F|\ge|F(g)|}. (The authors call this Siegel's argument. The catalogue gives C.~Poor and D.~Yuen,
\emph{Linear dependence among Siegel modular forms}, Math.\ Ann.\ 318 (2000), §1, Theorems 1.2--1.3, as an analogue
for Siegel modular forms; the unitary estimate and its constant have not been located in the literature.) The
reproducing formula for theta kernels in the stable range is the artifact's own: the catalogue locates the
classical basis theorem (S.~Böcherer, \emph{Über die Fourier--Jacobi-Entwicklung Siegelscher Eisensteinreihen},
§1, Theorem~III, p.~22) and a unitary pairing in the stable range (J.-S.~Li, \emph{Non-vanishing theorems for the
cohomology of certain arithmetic quotients}, §2, Theorem~2.1, p.~184), and the finite unitary formula used here is
distinct from both. Auerbach's lemma, in the form used here (a field with an embedding into a normed field, a
compact separating family of linear forms, the coordinates chosen in the family), is a determinant-maximisation
extension of the classical lemma (Schlumprecht, course notes for Functional Analysis~I, Theorem~1.5.1); this
form has not been located in a published source. A short basis of a lattice (W.~Banaszczyk, Math.\ Ann.\ 296
(1993), Theorems (2.1)--(2.2), pp.~631--632; J.~C.~Lagarias, H.~W.~Lenstra and C.~P.~Schnorr, \emph{Korkin--Zolotarev
bases and successive minima of a lattice and its reciprocal lattice}, Theorems 2.4--2.5, pp.~336--337; the formal
radius is weaker than the one these give). The Cauchy--Schwarz inequality for conjugated pairings of forms,
\m{|\sigma\ip{y}{x}|\le\|x\|_{\sigma}\,\|y\|_{\tilde\sigma}} with
\m{\tilde\sigma=\mathrm{conj}\circ\sigma\circ\mathrm{conj}} and \m{\|x\|_\sigma} the square root of \m{\nsq{x}{\sigma}}.
Holomorphic minors as coefficients of discrete series, their circle mean and their reproducing identity; Poisson
summation for theta kernels; the dual isometry of Weyl modules (Fulton and Harris, \emph{Representation theory},
§15.5, as cited in the artifact); the Fock model. References to be confirmed where the catalogue has none.}

\subhead{Status of the proofs.} A proof headed ``Proof'' is complete relative to the inputs it names. A proof
headed ``an account of the formal proof'' says which results combine, in the order of the formal proof, and does
not verify the estimates, the exponents and the normalising constants; an ``outline'' does less and says what it
leaves out.

\sect{Proofs}{s:proofs}
\subsection*{The line: counts, coordinates, arithmetic}
\begin{lemma}[counts and index maps]\label{l:counts}
In Setting~\ref{set:data}, let $j$ be a lift index and $\tau$ an infinite place.
\begin{enumerate}[label=\textup{(\alph*)},leftmargin=2.3em,itemsep=2pt]
\item \m{|\DetIx{j}{\tau}|+\bigCharExp{j}=w_j(\tau)}.
\item \m{|\DetIx{j}{\tau}|=\sum_{h<|\BigIndex{j}|}\kappa^{+}_\tau(h)}, and
\m{|\DetIx{j}{\tau}|\le b_1+b_2+\mathsf{g}(D)+m\,\mathsf{g}(D)}, with $b_1$, $b_2$ the auxiliary bounds of the data.
\item The maps \m{i\mapsto\negIx{j}{i}} and, for every~$h$, \m{i\mapsto\posRowIndex{j}{h}{i}} are injective, the sign
\m{\padSigns{j}} is $-1$ on the values of the first and $+1$ on the values of the second.
\item The integers
\m{\lfloor t(D)/2\rfloor-\operatorname{sgn}(\sigma)\bigl(\rowpar{\tau_\sigma}{q}+w_j(\tau_\sigma)\bigr)}, \m{0\le q<r}, at the
embeddings~$\sigma$ of \m{E(D)}, with \m{\tau_\sigma} the infinite place of~$\sigma$, \m{\rowpar{\tau_\sigma}{q}} the row
parameters of the data and \m{\operatorname{sgn}(\sigma)=\pm1} according as $\sigma$ belongs to the CM type or not, satisfy
the jump conditions of a row type for \m{(r,t(D))} of the unit \emph{Filling}.
\end{enumerate}
\end{lemma}
\begin{proof}[Proof\acct]
(a) and~(b): the number of columns of height below \m{k(D)+j-1} is a telescoping sum of the row parameters of the
weight of~$W_j(D)$; adding the \m{\mathsf{g}'(D)} further elements and the exponent gives the weight, and the bound
follows from~(a). (c) and~(d) are case distinctions on the definitions, with the inequalities between the scales
for~(d).
\end{proof}
\leanline{\lean{carddetIndexaddbigCharExp}\sep \lean{carddetIndex}\sep \lean{negIndexposRowIndexspec}\sep
\lean{isRowTypeauxPair}\sep the sign of the conjugate embedding is the opposite one, \lean{cmSignstar}}

\begin{lemma}[the line coordinates]\label{l:coords}
In Setting~\ref{set:data}, let $j$ be a lift index, $i$ a line choice and $Y$ a matrix over \m{E(D)} of the size of
\m{\shortRows{D}{i}}. Write \m{\Delta(j,i,Y)} for the product of the stacked determinants of the matrices
\m{\detRows{j}{i}{Y}{\tau}_l} along \m{\posEquiv{j}}, read as a polynomial in the first columns.
\begin{enumerate}[label=\textup{(\alph*)},leftmargin=2.3em,itemsep=2pt]
\item \m{\operatorname{pol}_{\kappa^{+}}\bigl(\linePolyCoords{j}{i}{Y}\bigr)=\Delta(j,i,Y)}, as polynomials in the first
columns with polynomial coefficients.
\item Let \m{\varphi\colon\CC\to\CC} be a ring homomorphism and $z$ a point of the space of the variables indexed by
\m{\LPlusIndex{j}}, and let $y$ be the vector of the numbers \m{\varphi(\linePolyCoords{j}{i}{Y}_l(z))}. Then
\m{\operatorname{pol}_{\kappa^{+}}(y)} is the image of \m{\Delta(j,i,Y)} under evaluation at~$z$ followed by~$\varphi$; and if
\m{|\varphi(\tau(Y_{ab}))|\le g(D)^4} for all $\tau$, $a$, $b$ and \m{|\varphi(z_x)|\le g(D)^4} for all~$x$, then
\m{\|y\|^2_{\kappa^{+}}\le\bigl(g(D)^{9}\bigr)^{31m\sum_\tau|\DetIx{j}{\tau}|}}.
\item There are an integer \m{\mathsf{d}\ge0} with \m{\mathsf{d}\le31m\sum_\tau|\DetIx{j}{\tau}|} and an integer~$H$
with \m{0<H\le(\mathsf{d}^{\mathsf{d}})^3} such that for every line choice~$i$, every~$Y$ and every standard
filling~$l$ the polynomial \m{Q=H\cdot\linePolyCoords{j}{i}{Y}_l} is homogeneous of degree~$\mathsf{d}$ and has its
coefficients in every subring
of~$\CC$ that contains all the numbers \m{\tau(Y_{ab})}.
\end{enumerate}
\end{lemma}
\begin{proof}[Proof\acct]
(a) Compare coefficients after evaluating the coefficient polynomials at every point: the coordinates in the first
columns are defined by Fock duality so that the polynomial of the coordinate vector is the given polynomial, once
it is semi-invariant of weight~$\kappa^{+}$, which the product of determinants is (unit \emph{Filling}).
(b) The first assertion is (a) after evaluation. For the bound, each stacked determinant is, by the lemma on a
determinant with fixed columns (Input~3), a number~$\gamma$ times a single minor, whose coordinate vector has
norm~$1$, and \m{|\gamma|^2} is at most a Gram determinant of a matrix with entries bounded by \m{g(D)^4}. By the
product rule for weights and norms (Input~3) the norms multiply over the \m{\sum_\tau|\DetIx{j}{\tau}|} factors.
(c) $H$ is the hook number of~$\kappa^{+}$. By the dual isometry of Weyl modules (library), $H$ times a coordinate
is an integral combination of coefficients of the product of determinants, which are polynomials with integer
coefficients in the entries of the \m{\tau(Y)}; the degree and its bound come from the substitution law for the
coordinates (Input~3).
\end{proof}
\leanline{\lean{toPolylinePolyCoords}\sep \lean{weylNormSqlinePolyCoords}\sep \lean{existsdenomlinePolyCoords}}

\begin{proposition}[the line multipliers are integral and small]\label{p:lineint}
In Setting~\ref{set:data}, let $j$ be a lift index. Then there is \m{c\in\CC}, \m{c\ne0}, such that for every choice
\m{i\in\MemberIndex{j}}, every tame character~$\vartheta$ of~$D$ and every \m{h\in\UU(I_{\BigIndex{j}})(\AAA_f)}, the
vector \m{y=\overline{c\cdot\lineMultiplier{\vartheta}{j}{i}{h}{\cdot}}} satisfies:
\begin{enumerate}[label=\textup{(\roman*)},leftmargin=2.3em,itemsep=2pt]
\item for every prime~$p$ prime to the norms of the primes in \m{Q(\vartheta)} and every embedding
\m{\iota_p\colon\CC\to\CC_p}: \m{\gau{\iota_p}{\kappa^{+}}{0}{h}{y}\le1};
\item if \m{Q(\vartheta)=\emptyset}, then \m{\|\varphi(y)\|^2_{\kappa^{+}}\le\exp(r\,B_0(D))} for every \m{\varphi\in\AutC}.
\end{enumerate}
Moreover, for a tame character~$\vartheta$ and an integer \m{c'\ge0}, the integer \m{N(Q(\vartheta))\,g(D)^{c'}} lies in the
modulus \m{\mathfrak{m}_D(\vartheta,c')}.
\end{proposition}
\begin{proof}[Proof\acct]
The last assertion: \m{\bar\pi_g^{c'}} divides \m{g(D)^{c'}}, and the product of the norms of the primes of
\m{Q(\vartheta)} lies in their product. For the rest, the line multiplier is a class sum of products of three theta
kernels of definite pairs: the kernel inside \m{\cXi{m}{\twistedChi{\vartheta}{j}}{c}}, the constant of the negative
definite space~$I_j^{-}$, whose kernel is the conjugate of a kernel of~$I_j$ up to a number of absolute value~$1$,
and the kernel of~$L^{+}$ with the line coordinates. Take $\mathsf{d}$ and~$H$ from Lemma~\ref{l:coords}\,(c) and the three
constants of Input~2\,(a) for these pairs, with degrees $0$, $0$ and~$\mathsf{d}$; let $c$ be their product with $H$, with
the square of the chosen group order and with the number of absolute value~$1$. (i) The moduli contain
\m{N(Q(\vartheta))g(D)^2}, the data are integral away from \m{N(Q(\vartheta))}, and class representatives can be moved
by rational points to points integral above~$p$, the kernels changing by the rule for rational points; the
coefficients of \m{H\cdot\linePolyCoords{j}{i_2}{Y}_l} are $\iota_p$-integral for a suitable choice of the short basis
at~$p$. So each factor is $\iota_p$-integral, and the orders of class stabilisers divide the chosen group order.
(ii) If \m{Q(\vartheta)=\emptyset}, then \m{q=1}; the conjugates of the three kernels are bounded by Input~2\,(a), with the bound
of Lemma~\ref{l:coords}\,(b) for the polynomial, the entries of the short basis being bounded at every complex
embedding. The numbers of classes and all other factors are at most \m{\exp(rB_0(D))} by the inequalities
between the scales.
\end{proof}
\leanline{\lean{lineMultiplierintegral}\sep \lean{normmulgProd}}

\subsection*{Sparse theta sums and isolation}
The next results are archimedean. For a number field~$E$, a finite non-empty index set~$\ix n$ and a sign
vector~$s$, \m{\|h\|} is the maximum over the infinite places of the operator norms of the components of
\m{h\in G_\infty(s)}; it is at least~$1$. For a weight~$k$, a family of integer vectors \m{k_\tau} on~$\ix n$,
\m{\operatorname{mh}_k} is the product over~$\tau$ of the holomorphic minors of weight~$k_\tau$, a matrix coefficient of
a holomorphic discrete series when $k$ is dominant with a gap. A function on the adelic group is \kw{holomorphic of
order~$o$ and weight~$k$} if it transforms under the maximal compact subgroup and is annihilated by the
antiholomorphic derivatives accordingly (library); an order~$o$ on~$\ix n$ is \kw{positive first} for~$s$ if the
positive indices come first. The \kw{first minimum} \m{\lambda_1} of a hermitian space moved by a point is the infimum
of~$1$ and of the lengths of the non-zero vectors of its lattice: the length of a shortest vector, cut off at~$1$.

\begin{setting}[weights for isolation]\label{set:iso}
$E$ is a number field with \m{f_E} infinite places, \m{\ix n} a finite non-empty index set, $s$ a sign vector
on~$\ix n$, \m{\nu\ge0} an integer and $B$ a real number with \m{f_E\cdot2|\ix n|^3\le\nu}, \m{\nu\le B} and
\m{10^{13}\le B}; $o$ is an order that is positive first for~$s$; $w$ and~$k$ are weights that are dominant
for~$o$ and have holomorphic gap \m{300\nu}, and \m{|k_\tau(i)|\le B} for all $\tau$, $i$.
\end{setting}

\begin{lemma}[integrals against holomorphic minors]\label{l:mh}
In Setting~\ref{set:iso}, the integral \m{\int|\operatorname{mh}_k(h)|^2dh} over \m{G_\infty(s)} is positive; and for
every continuous function~$A$ on \m{G_\infty(s)}, every real number~$C$ and every integer \m{N\le6\nu} such that
\m{|A(h)|\le C\|h\|^{N}} for all~$h$, the product \m{A\cdot\operatorname{mh}_w} is integrable and
\[\Bigl|\int A(h)\operatorname{mh}_w(h)\,dh\Bigr|\le C\,B^{7\nu}\int|\operatorname{mh}_k(h)|^2dh .\]
\end{lemma}
\begin{proof}[Proof\acct]
The holomorphic minors are coefficients of discrete series; their $L^2$ norms are inverses of formal degrees,
hence positive and comparable for $w$ and~$k$ up to powers of~$B$, and their decay integrates the polynomial
growth of~$A$ (library).
\end{proof}
\leanline{\lean{integralmulholMinors}\sep the bound \m{1\le\|h\|} is \lean{onelemaxOpNorm}}

\begin{proposition}[isolation]\label{p:iso}
In Setting~\ref{set:iso}, let $E$ be a CM field, $F$ a continuous function on the adelic group that is holomorphic
of order~$o$ and weight~$k$, $z$ a point with \m{F(z)\ne0} and \m{|F(x)|\le2|F(z)|} for all~$x$, let $K_0$ be a
function on \m{G_\infty(s)} with \m{|K_0(h)|\le C} for all~$h$, for a real number~$C$, and let $\varepsilon$ be real.
Suppose: if \m{\operatorname{mh}_w} is integrable, then for every~$x$ with \m{|F(z)|\le|F(x)|} there is a function~$K$,
holomorphic of weight \m{k-w}, with \m{|K(h)-K_0(h)|\le\|h\|^{6\nu}\varepsilon} for all~$h$ and
\m{\int\overline{F(x\,(h,1))}\,K(h)\operatorname{mh}_w(h)\,dh=0}. Then \m{|K_0(1)|\le4\varepsilon B^{7\nu}+\varepsilon}.
\end{proposition}
\begin{proof}
By Lemma~\ref{l:mh}, \m{\operatorname{mh}_w} is integrable and \m{I=\int|\operatorname{mh}_k|^2>0}. For an angle~$\vartheta$ let
\m{z_\vartheta} be $z$ times the element \m{e^{i\vartheta o}} of the diagonal torus at infinity, with entries
\m{e^{i\vartheta o(i)}}, and let \m{\chi_\vartheta=\prod_\tau\prod_{i\in\ix n}e^{i\vartheta\,o(i)\,k_\tau(i)}}, a number of
absolute value~$1$: the conjugate of the factor of the library's circle mean. By holomorphy,
\m{|F(z_\vartheta)|=|F(z)|}, so the hypothesis gives for every angle a function \m{K^{(\vartheta)}}, holomorphic of weight
\m{k-w}, within \m{\|h\|^{6\nu}\varepsilon} of~$K_0$, with
\[\int W_\vartheta(h)\,K^{(\vartheta)}(h)\,dh=0,\qquad
W_\vartheta(h)=\chi_\vartheta\,\overline{F(z_\vartheta(h,1))}\,\operatorname{mh}_w(h);\]
the constant~\m{\chi_\vartheta} does not change the vanishing. Hence
\[\int W_\vartheta K^{(0)}=\int W_\vartheta\,\bigl(K^{(0)}-K^{(\vartheta)}\bigr),\]
and since \m{|F|\le2|F(z)|} and
\m{|K^{(0)}-K^{(\vartheta)}|\le2\|h\|^{6\nu}\varepsilon}, Lemma~\ref{l:mh} with \m{N=6\nu} bounds this by
\m{4\varepsilon|F(z)|B^{7\nu}I}, for every angle. On the other hand, by the circle mean for holomorphic functions of
weight~$k$ (library), the mean over the angles of \m{\chi_\vartheta\,\overline{F(z_\vartheta(h,1))}} is
\m{\overline{F(z)}\cdot\overline{\operatorname{mh}_k(h)}}, so that the mean of \m{W_\vartheta(h)} is
\[\overline{F(z)}\cdot\overline{\operatorname{mh}_k(h)}\cdot\operatorname{mh}_w(h).\]
The function \m{K^{(0)}\operatorname{mh}_w} is holomorphic of weight~$k$, so by the reproducing identity (library) its
integral against the conjugate of \m{\operatorname{mh}_k} is its value at~$1$ times~$I$. Exchanging the two integrals
gives \m{|F(z)|\,|K^{(0)}(1)|\,I\le4\varepsilon|F(z)|B^{7\nu}I}, that is \m{|K^{(0)}(1)|\le4\varepsilon B^{7\nu}}; and
\m{|K_0(1)|\le|K^{(0)}(1)|+\varepsilon}.
\end{proof}
\leanline{\lean{isolation}}

\begin{lemma}[first minima and matrix coefficients]\label{l:sparse}
Let $E$ be a field with a CM type, \m{\ix n}, \m{\ix n'} finite non-empty index sets, $S$ and~$T$ hermitian spaces
with index sets \m{\ix n}, \m{\ix n'}.
\begin{enumerate}[label=\textup{(\alph*)},leftmargin=2.3em,itemsep=2pt]
\item If \m{N\ge1} is an integer in the different of~$E$, then, with \m{c_N=6N(f_E|\ix n|)^2},
\[\frac{\lambda_1(S)\,\lambda_1(T)}{c_N\,|\ix n'|\,|\ix n|}\le\lambda_1(\operatorname{Hom}(S,T)).\]
\item Let $E$ have Weil data, $\mathfrak{m}$ a modulus, $t$ a matrix over \m{\OEr/\mathfrak{m}}, with trace character
\m{\psi'_t(y)=\psi'(\operatorname{tr}(yt))}, $v$ and~$w$ Fock vectors and $z$ an adelic vector of
\m{\operatorname{Hom}(S,T)}. Then the pairing of the datum \m{v\otimes\psi'_t} with the Heisenberg translate by~$z$ of
\m{w\otimes\psi'_t} is the value of~$\psi'_t$ at the finite part of~$z$ times the archimedean matrix coefficient of
$v$, $w$ at the real part of~$z$ times the vacuum function at the finite part.
\item Let $E$ have Weil data, $a$ and~$b$ characters with weights $\gamma$, $\gamma'$ and tame parts, $\mathfrak{m}$
a modulus, $g$, $h$ points, \m{c\in E} non-zero such that $c$ times the idele of the Weil data times~$z$ is
integral for every integral adele~$z$ that reduces to~$0$ modulo~$\mathfrak{m}$, \m{\mu>0} a real number with
\m{\mu\,|\tau(c)|\le\sqrt{2\pi\operatorname{Im}\tau(\jmath)}\cdot\lambda_1(\operatorname{Hom}(S^{g},T^{h}))} at every infinite
place~$\tau$, with $\jmath$ the purely imaginary element of the Weil data and \m{S^{g}}, \m{T^{h}} the spaces moved
by the points, \m{q>f_E|\ix n'||\ix n|} an integer, and $Q$ a polynomial, homogeneous of degree~$\mathsf{d}$ for real
scalars, with Fock vector \m{v=v_Q}. Then there is a function~$\Psi$ on matrices modulo~$\mathfrak{m}$ such that for
all \m{u\in G_\infty(s)}, \m{u'\in G_\infty(s')}
\begin{multline*}
\Bigl|\tk{S}{T}{a}{b}{\mathfrak{m}}{\Psi}{Q}\bigl(g\,(u,1),h\,(u',1)\bigr)-\bigl\langle v,\omega^{\infty}_{\gamma,\gamma'}(u,u')v\bigr\rangle\Bigr|\\
\le\|u'\|^{3q}\Bigl(\frac{100(\mathsf{d}+2q)}{\mu^2}\Bigr)^{q}|\ip{v}{v}|\,\|u\|^{3q}.
\end{multline*}
\end{enumerate}
\end{lemma}
\begin{proof}[Proof\acct]
(a) The first minima are cut off at~$1$. So it suffices to show that the left side is at most~$1$ and at most the
length of every non-zero vector~$z$ of the lattice of \m{\operatorname{Hom}(S,T)}. The first holds because
\m{\lambda_1(S)\le1}, \m{\lambda_1(T)\le1} and \m{c_N\,|\ix n'|\,|\ix n|\ge1}. For the second, a short basis~$Y$ of the
lattice of~$S$ (library), with vectors of length at most \m{c_N/\lambda_1(S)}, turns~$z$ into the vectors \m{zY_k}
of~$T$, one of them non-zero, of length at least \m{\lambda_1(T)} and at most \m{|z||Y_k|}; the factor
\m{|\ix n'|\,|\ix n|} allows for the comparison of the norms of vectors and of matrices in the formal text. (b) Lift~$t$,
write the Heisenberg translate of the named datum as a translate of the character, and compute place by place.
(c) Let $\Psi$ be the average over~$t$ of \m{\overline{\theta[\psi'_t,Q](g,h)}\,\psi'_t}. By Poisson summation the
kernel is a sum over the lattice of \m{\operatorname{Hom}(S^{g},T^{h})} of archimedean matrix coefficients, computed
by~(b); the term of the vector~$0$ is \m{\ip{v}{\omega^{\infty}_{\gamma,\gamma'}(u,u')v}}. The other lattice vectors have
length at least the first minimum, hence at least~$\mu$ in the normalised coordinates, and the Gaussian decay of
the matrix coefficients of a Fock vector of degree~$\mathsf{d}$, summed over the lattice with \m{q>f_E|\ix n'||\ix n|},
gives the error.
\end{proof}
\leanline{\lean{firstMinimummulle}\sep \lean{dataInnertraceChar}\sep \lean{normmodKernelsub}}

\subsection*{The carrier and the approximation by \m{F^{\natural}}}
\begin{lemma}[the carrier is rapid, holomorphic and not zero]\label{l:carrier}
Let $E$ be a CM field with a CM type and Weil data and \m{m\ge1} an integer.
\begin{enumerate}[label=\textup{(\alph*)},leftmargin=2.3em,itemsep=2pt]
\item Let $m$ be even, and let $\mathfrak{m}$ be a modulus, $c$ a character and $\chi$ a character of
\m{\OEr/\mathfrak{m}} with $c$ tame for \m{(\mathfrak{m},0,\chi)}. Then \m{\carrier{E}{m}} is rapidly decreasing and automorphic
of every level, and there is an order~$o$, positive first for the signs of \m{I_{2m,m}}, for which
\m{\carrier{E}{m}} is holomorphic of weight $0$ on the positive block and $3m$ on the negative block.
\item \m{\carrier{E}{m}\ne0}.
\end{enumerate}
\end{lemma}
\begin{proof}[Proof\acct]
(a) By the carrier identity (Input~3) the carrier is a finite combination of theta kernels of the pair
\m{(I_{2m,m},I_{3m})} with datum \m{\chi\circ\det}, which are rapidly decreasing (the datum is supported on invertible
matrices), invariant under every level through the tame character, and holomorphic of the stated weight because
the polynomial~$1$ is a holomorphic datum for a definite second space. (b) If the lift of the constant~$1$
vanished, then by the adjunction between lift and descent the integral of the constant~$1$ against the descent of
a suitable datum would vanish; for the polynomial \m{\prod_\tau\det(y^{-}_\tau)^{2m}} the archimedean matrix
coefficient is a product of holomorphic minors with non-zero value at~$1$, the datum is sharp, and the volume of
the compact group is not zero.
\end{proof}
\leanline{\lean{carrierFormspec}\sep \lean{carrierFormnezero}}

\begin{proposition}[the function F-natural approximates a test function]\label{p:approx}
In Setting~\ref{set:data}, let $j$ be a lift index and \m{\nu_D} the integer~$\nu$ of the data for which
Setting~\ref{set:iso} holds with \m{B=5b}. Then there are a function~$K_0$ on the archimedean group of~$W_j(D)$ and
real numbers $C$, $\varepsilon$ with
\[4\varepsilon\,(5b)^{7\nu_D}+\varepsilon<|K_0(1)|,\qquad|K_0(y)|\le C\ \text{for all }y,\]
such that for every \m{g\in\UU(W_j(D))(\AAA)} with \m{(5b\cdot b^{23})^{-1/2}\le\lambda_1(W_j(D)^{g})} there are
complex numbers \m{c_a}, \m{a\in\MemberIndex{j}}, with
\[\Bigl|\sum_{a}c_a\,\fNatural{j}{a}{g\,(y,1)}-K_0(y)\Bigr|\le\|y\|^{6\nu_D}\,\varepsilon\qquad\text{for all archimedean }y.\]
\end{proposition}
\begin{proof}[Outline of the proof]
The outline follows the formal proof step by step and names the result that each step uses. It does not reproduce
the estimates: the lower bound for the first minimum of the lattice of homomorphisms, the size of the
constant~$\varepsilon$ and the inequality between \m{\varepsilon(5b)^{7\nu_D}} and \m{|K_0(1)|} are closed in the formal
text by the inequalities between the scales and were not verified here.

\emph{The main term.} The padded space \m{W_j(D)\oplus I_j^{-}\oplus I_{3m-j}} has $m$ negative signs
(Lemma~\ref{l:counts}\,(c)). By Lemma~\ref{l:carrier} the carrier is rapidly decreasing, holomorphic of weight
\m{(0,3m)} and not zero; by the library's result on near maxima it has a point~$x$ with
\m{|\carrier{E(D)}{m}|\le2|\carrier{E(D)}{m}(x)|} everywhere. The archimedean matrix coefficient
\[C(u,y)=\ip{v}{\omega^{\infty}(u,y)v},\qquad v=v_{\linePoly{j}},\]
of the Fock vector of the line polynomial is a product of
holomorphic minors in each variable, with weights determined by the exponents of the line polynomial and
Lemma~\ref{l:counts}\,(a). Put
\[K_0(y)=\int\carrier{E(D)}{m}(x\,u)\,\operatorname{mh}_{w^{\circ}}(u)\,C(u,y)\,du,\]
the integral over the archimedean group of signature \m{(2m,m)}, with \m{w^{\circ}} the weight
\m{(-\mathsf{g}'(D),6m+1)} on the two blocks. It is bounded, by Lemma~\ref{l:mh} with \m{N=0}. At \m{y=1} the integrand is a
holomorphic function times the conjugate of a holomorphic minor of its weight, and the integral is its value at
\m{u=1} times \m{\int|\operatorname{mh}|^2>0}: \m{K_0(1)} is a positive multiple of \m{\carrier{E(D)}{m}(x)C(1,1)\ne0}.

\emph{The approximation.} Let $g$ be as in the statement. By Lemma~\ref{l:sparse}\,(a), the bound for the first
minimum at~$x$ from the near maximum and the hypothesis on~$g$, the lattice of
\m{\operatorname{Hom}} between the carrier space at~$xu$ and the padded space at \m{g(y,1)} has first minimum bounded below.
Lemma~\ref{l:sparse}\,(c), for the pair of the carrier space and the padded space, the polynomial \m{\linePoly{j}} and
\m{q=2\nu_D}, gives a function~$\Psi$ modulo~\m{\bar\pi_g^{2}} whose kernel differs from \m{C(u,y)} by at most
\m{\|u\|^{6\nu_D}\|y\|^{6\nu_D}} times a small constant. The reproducing formula (library), for the pair of the carrier
space and \m{I_{\FlatIndex{D}}}, expresses the pairing of the rapidly decreasing form \m{\carrier{E(D)}{m}\operatorname{mh}_{w^{\circ}}}
with theta kernels as a finite combination, with coefficients \m{c(i_2)}, of the kernels with the determinant
polynomials of the short rows of the line choices~$i_2$. The product of~$\Psi$ with the indicator of an inclusion
matrix is a finite combination of trace characters \m{\psi'_{i_1}}, with coefficients \m{f(i_1)} (finite Fourier
analysis on \m{\OEr/\bar\pi_g^{2}}). With \m{c_a=c(i_2)f(i_1)} for \m{a=(i_1,i_2)}, the seesaw for the three summands
of~$W'$ identifies \m{\sum_ac_a\fNatural{j}{a}{g(y,1)}} with the integral of the carrier against the kernel of~$\Psi$;
replacing this kernel by \m{C(u,y)} costs, by Lemma~\ref{l:mh} with \m{N=6\nu_D} and \m{|\carrier{E(D)}{m}|\le2|\carrier{E(D)}{m}(x)|}, at most
\m{\|y\|^{6\nu_D}\varepsilon}, with $\varepsilon$ a multiple of the small constant. The inequality between
\m{\varepsilon(5b)^{7\nu_D}} and \m{|K_0(1)|} is one of the inequalities between the scales.
\end{proof}
\leanline{\lean{existsfNaturalapprox}}

\subsection*{Theorems C and D}
\begin{proof}[Proof of Theorem C\acct]
Fix~$j$. The padded space \m{W_j(D)\oplus I_j^{-}\oplus I_{3m-j}\oplus I_{\FlatIndex{D}}\oplus S}, regrouped by
\m{\regroupEquiv{j}}, has the block signs of a bijection built from \m{\negIx{j}{\cdot}} and the row indices
(Lemma~\ref{l:counts}\,(c)), and for every twist~$\vartheta$ the characters satisfy
\m{\bigChar{\vartheta}{j}\,\upsilon^{m}=\twistedChi{\vartheta}{j}\,\carrierChar{j}}, by the definitions. Two results of
Input~3 give constants that do not depend on the twist. The explicit matching, for the pair of \m{W'\oplus S} and
\m{L^{+}}-type blocks, with the bijection \m{\posEquiv{j}} and the index sets \m{\DetIx{j}{\tau}}, whose sizes are the
required weights by Lemma~\ref{l:counts}\,(a), gives $C_1$ and, for the tame character
\m{\twistedChi{\vartheta}{j}\carrierChar{j}}, a number~$c_1(\vartheta)$ with \m{|c_1(\vartheta)|^2} equal to $C_1$ divided by a product
of $L$-values, and an identity between a class sum of kernels with the product of stacked determinants of the
matrices \m{\detRows{j}{i_2}{Y}{\tau}} and an integral over \m{[\UU(S)]} of kernels with the matched polynomial; the matched
polynomial is \m{\linePoly{j}} times the determinants of the short rows, as one sees by comparing the minors. The
carrier identity gives $C_2$ and, for the tame character \m{\twistedChi{\vartheta}{j}} of weight~$0$, a number
\m{c_2(\vartheta)} with \m{c_2(\vartheta)\clsum{c}{\KK{3m}}\cXi{m}{\twistedChi{\vartheta}{j}}{c}\cdot\theta(x,c)=\carrier{E(D)}{m}(x)}.
Now apply the matching summed against a form on a square block (Input~3) with the function
\m{\Xi=\cXi{m}{\twistedChi{\vartheta}{j}}{\cdot}}, the function \m{F=\carrier{E(D)}{m}}, the modulus
\m{\mathfrak{n}=\bar\pi_g^{2}}, the datum \m{\psi'_{i_1}} and the polynomial of \m{\fNatural{j}{i}{\cdot}}: its right side is
\m{\fNatural{j}{i}{g}}, and its left side is a double class sum of kernels with the datum
\m{\check\Psi_0(x'\mid x'')=\chi(\det x'')\sum_k\psi'_{i_1}((x''^{-1}x'\bmod\bar\pi_g^{2})k)}. By the passage from the
stabiliser to the full level (Input~3), for the matrix~$w$ and the function \m{y\mapsto\sum_k\psi'_{i_1}(yk)}, which is
invariant under the stabiliser of~$w$ because \m{uw=i_1}, the class sum at the full level is the class sum at the
level of the stabiliser with the column datum. The seesaw for the three blocks of~$L^{+}$ and the expression of the
product of determinants through its coordinates in the positive letters (unit \emph{Filling}) turn the inner sums
into \m{\sum_l\lineMultiplier{\vartheta}{j}{i}{h}{l}} times the $l$-th basis polynomial, that is into the kernel with
the polynomial of the line multiplier in the positive letters: the coset kernel sum. So the identity of the
theorem holds with \m{c(\vartheta)=c_1(\vartheta)\,c_2(\vartheta)/|\GL_{3m}(\OEr/\bar\pi_g^{2})|}, by the formula of the
matching summed against a form on a square block, and \m{|c(\vartheta)|^2} is a constant divided by a product of products of $L$-values at the characters
\m{(\twistedChi{\vartheta}{j}\carrierChar{j})^{-1}\upsilon^{\cdots}}, \m{\twistedChi{\vartheta}{j}^{-1}\upsilon^{m}} and
\m{\twistedChi{\vartheta}{j}\upsilon^{m/2}}; unfolding the characters, these are the three factors of
\m{\heckeValueProd{j}{\vartheta}}.
\end{proof}

\begin{lemma}[levels and holomorphy]\label{l:rapid}
In Setting~\ref{set:data}, let $j$ be a lift index and \m{F\in\Mgood{j}}.
\begin{enumerate}[label=\textup{(\alph*)},leftmargin=2.3em,itemsep=2pt]
\item For every integer~$a$ of \m{E(D)} divisible by~\m{\bar\pi_g}, $F$ is rapidly decreasing and invariant under the
principal level modulo~$a$.
\item For every choice \m{i\in\MemberIndex{j}}, the product \m{F\cdot\fNatural{j}{i}{\cdot}} is rapidly decreasing and
invariant under the principal level modulo~\m{\bar\pi_g^{2}}, and \m{\fNatural{j}{i}{\cdot}} is holomorphic, for a fixed
order that is positive first for the signs of~$W_j(D)$, of weight the difference of the weight of~$\bar F$ and the
weight of the kernels of \m{\tableTheta{j}{\cdot}{\cdot}}.
\end{enumerate}
\end{lemma}
\begin{proof}[Proof\acct]
(a) Good forms are cusp forms of principal level at the chosen primes. (b) By Theorem~C with \m{\ix a=\CarrierIndex{D}},
\m{u=1} and \m{w=i_1}, \m{\fNatural{j}{i}{\cdot}} is a finite combination of theta kernels of the pair of~$W_j(D)$ and
the big block with polynomials in the positive letters; such kernels are holomorphic, the Fock vector being a
highest weight vector (library), and the product of a rapidly decreasing form with a kernel of this kind is
rapidly decreasing of the level of the modulus.
\end{proof}
\leanline{\lean{goodFormsrapid}\sep \lean{mulfNaturalrapid}}

\begin{proof}[Proof of Theorem D\acct]
Suppose \m{F\ne0}. Its conjugate~$\bar F$ is holomorphic, of the weight~$k$ of the kernels of the adjoint lift,
because $F$ is in the image of the descent of the trivial twist (the adjoint lift is injective). By the library's
result on near maxima there is a point~$g_0$ with \m{\bar F(g_0)\ne0}, \m{|F|\le2|F(g_0)|} everywhere and first
minimum at least \m{(5b\cdot b^{23})^{-1/2}} wherever \m{|F|\ge|F(g_0)|}. Take $K_0$, $C$, $\varepsilon$ from
Proposition~\ref{p:approx}. The claim is that the hypothesis of Proposition~\ref{p:iso} holds for~$\bar F$, the
point~$g_0$, the weight~$w$ of the kernels of \m{\tableTheta{j}{\cdot}{\cdot}} and \m{B=5b}; its conclusion
\m{|K_0(1)|\le4\varepsilon(5b)^{7\nu_D}+\varepsilon} then contradicts Proposition~\ref{p:approx}. Let $x$ be a point with
\m{|F(g_0)|\le|F(x)|}, so that Proposition~\ref{p:approx} applies at~$x$, and put
\m{K(y)=\sum_ac_a\fNatural{j}{a}{x(y,1)}}; it is holomorphic of weight \m{k-w} by Lemma~\ref{l:rapid}\,(b) and within
\m{\|y\|^{6\nu_D}\varepsilon} of~$K_0$. The integral \m{\int F(x(y,1))K(y)\operatorname{mh}_w(y)\,dy} is, by the reproducing
formula for the pair of~$W_j(D)$ and the small block~$V^{\circ}$ (library), applied to the rapidly decreasing
function \m{F\cdot\sum_ac_a\fNatural{j}{a}{\cdot}}, a finite combination of the integrals over \m{[\UU(W_j(D))]} of
\m{\overline{F\,\fNatural{j}{a}{\cdot}}} against kernels of the small block with basis polynomials in the negative
letters at the class representatives. Each of these is, after writing a basis polynomial as the polynomial of a
coordinate vector~$v$, the integral of the hypothesis for the function~$\zeta$ equal to~$\bar v$ at one class and
zero elsewhere, hence zero.
\end{proof}
\leanline{\lean{eqzeroofintegral}}

\subsection*{Hecke values: Theorem B}
\begin{lemma}[flat pairs and the untwisted norm]\label{l:flat}
In Setting~\ref{set:data}:
\begin{enumerate}[label=\textup{(\alph*)},leftmargin=2.3em,itemsep=2pt]
\item Let \m{\ix n} be a finite index set with \m{4\le|\ix n|\le t(D)} and \m{|\ix n|} even, and $\chi$ a character of
weight~$0$ whose tame part is trivial outside the chosen primes. Then there is a pair~$p$ that is flat
for~$\ix n$, with \m{\nu_p\upsilon^{|\ix n|/2}=\chi} and with the tame part of~$\chi$.
\item Let $j$ be a lift index and $p$ a pair of rank one for~$j$. Then the constant family
\m{(\ell,\iota)\mapsto\|\triv{p}{1}\|^2} has height at most \m{B_0(D)-\sum_\tau\operatorname{js}_D(w_0(\tau))}.
\end{enumerate}
\end{lemma}
\begin{proof}[Proof\acct]
(a) Take \m{\nu_p=\chi\upsilon^{-|\ix n|/2}} and \m{z_p=|\ix n|/2}; the conditions on a pair hold by the inequalities
between the scales. (b) Take a flat pair~$q$ for the index set of~$p$ by~(a) and the witness
\m{H=\ialg{q}{1}/\|\triv{p}{1}\|^2}: then \m{H\cdot\|\triv{p}{1}\|^2=\ialg{q}{1}} is an algebraic integer, so the $\ell$-adic
exponents are~$0$, and the conjugates of~$H$ are bounded by the unit \emph{Exchange}, which computes the logarithm
of the conjugates of the squared norm of a trivial lift as a sum of jet sizes up to terms that the scales bound
by \m{B_0(D)}.
\end{proof}
\leanline{\lean{existsisFlat}\sep \lean{heightLenormtrivialLift}}

\begin{lemma}[two pigeonholes]\label{l:pigeon}
\textup{(a)} Let $E$ be a CM field with a CM type, $\ell$ a prime, \m{\iota\colon\CC\to\CC_\ell} a ring homomorphism,
\m{w=(w_i)} a family of finite places, all in a set~$S$ of finite places, $s$ a complex function on
\m{\UU(I_1)(\AAA_f)}, and \m{\Gamma}, $K$ subgroups such that $\Gamma$ contains the rational points, $K$ contains the
principal level at the places~$w_i$, the double coset space is finite and all class stabilisers are trivial. Then
there are a point~$z$ and a character~$\chi$ of weight~$0$ with tame part trivial outside~$S$ such that
\m{|\iota(s(1))|_\ell} is at most the $\ell$-adic absolute value of the mean over the classes of
\m{g\mapsto\chi(\det(gz))\,s(gz)}.
\textup{(b)} Let $G$ be a finite group, $\iota$ a ring homomorphism from~$\CC$ to a normed field, $a$ and~$b$ complex
functions on~$G$ and \m{V\in\CC} with \m{V^{|G|}=\prod_{\chi\in G}a(\chi)b(\chi)}. Then for every \m{\psi\in G} there are
\m{\phi,\chi\in G} with \m{\phi\chi=\psi} and \m{|\iota(V)|\le|\iota(a(\phi)b(\chi))|}.
\end{lemma}
\begin{proof}
(a) The class group of \m{\UU(1)} at this level is a finite abelian group, and every character of it is a character
of weight~$0$ with tame part trivial outside~$S$ (library). By orthogonality of characters, \m{s(1)} is the sum over
the characters of the means in question, for~$z$ the inverse of a representative of the trivial class; by the
ultrametric inequality one of them is at least as large. (b) Choose~$\chi$ maximising
\m{|\iota(a(\psi\chi^{-1})b(\chi))|}; the product over~$\chi$ of these numbers is \m{|\iota(V)|^{|G|}}, after the change of
variable \m{\chi\mapsto\psi\chi^{-1}} in~$a$.
\end{proof}
\leanline{\lean{existstwistedclassSum}\sep \lean{existsfactorisation}}

\begin{lemma}[test forms]\label{l:test}
In Setting~\ref{set:data}, let $\ell$ be a prime with an embedding \m{\iota\colon\CC\to\CC_\ell}.
\begin{enumerate}[label=\textup{(\alph*)},leftmargin=2.3em,itemsep=2pt]
\item Let $q$ be a pair that is flat for \m{\{1,\dots,m'\}}, $X$ an $L$-value datum, $\vartheta$ a twist with
$\ell$ prime to \m{N(Q(\vartheta))}, $\mathcal{N}$ a module over the valuation ring of~$\iota$ inside \m{\FormVx{O_X}}, and $F$
a function with values in~$\mathcal{N}$. Then \m{\xiSum{q}{\vartheta}{F}\in\mathcal{N}}.
\item Let $q$, $X$, $\vartheta$, $F$ be as in Definition~\ref{d:aux}\,(i), $g$ a point, \m{A,B\in\CC} and $\Lambda$ a
continuous conjugate-linear form on \m{\FormVx{O_X}} such that, in every coordinate and for all~$h$,
\[\Lambda(F(h))=B\cdot\bigl(A\cdot\text{the kernel form of }q\text{ at }(g,h)\bigr).\]
Then \m{\clU\,\kappa_0^{m'}\cdot\Lambda(\xiSum{q}{\vartheta}{F})=B\cdot\ialg{q}{\vartheta}\cdot A}, with \m{\kappa_0} the vacuum
scale and \m{\clU} the class sum of the constant~$1$ over \m{\UU(1)} at the chosen level of the data: the sum over
the classes of the inverses of the orders of their stabilisers. That level is free (unit \emph{Data}), so the
stabilisers are trivial and \m{\clU} is the number of classes.
\item Let $p$ be a pair of rank one for a lift index~$j$, let $q$, $q_2$, $q_3$ be pairs flat for
\m{\{1,\dots,2z_p-20\}}, \m{\{1,\dots,4\}} and \m{\{1,\dots,4\}}, $\vartheta$ a twist of~$D$ with $\ell$ prime to
\m{N(Q(\vartheta))}, $x_0$ a \m{12\times1} matrix modulo~\m{\bar\pi_g}, $\alpha$ an exponent whose total orders are the jet
degrees of the $L$-value datum~$X$ of~$p$ at~$\vartheta$, and $b$ a character whose product with the second characters
of the data of $q_2$ at~$\vartheta$, of $q_3$ at~$1$ and of~$q$ at~$1$ is the second character of~$X$. Then there is a
form \m{\Xi\in\FormVx{O_X}} that is a \kw{test form at~$\ell$}, integral at every embedding into~$\CC_\ell$ and tested
against the kernel of the descent for the label idempotent \m{\idem(X-1)}, with, for every point~$g$ of \m{\UU(1)},
\begin{align*}
\ip{\Xi}{\text{kernel form of }X\text{ at }g}\cdot\clU^{3}\kappa_0^{2z_p-12}
={}&\frac{\prod_\tau d_\tau!}{\alpha!}\cdot\ialg{q}{1}\,\ialg{q_3}{1}\,\ialg{q_2}{\vartheta}\\
&\cdot\tk{I_1}{I_{12}}{1}{b}{\bar\pi_g}{\mathbf{1}_{x_0}}{x^{\alpha}}(g,1),
\end{align*}
with \m{\clU} as in~(b) and $d_\tau$ the jet degrees.
\end{enumerate}
\end{lemma}
\begin{proof}[Proof\acct]
(a) The class sum is an integral combination, the orders of class stabilisers dividing~$M_D$; the factor
\m{\delta_q/N(Q(\vartheta))^{m'-1}} is $\iota$-integral since $\ell$ is prime to \m{N(Q(\vartheta))}, and the Gram form of~$1$
with \m{\delta_q} times the trivial lift is $\iota$-integral by the definition of the denominator bound. (b) Expand the
class sum and use that the pairing of the trivial lift with the kernel forms is the theta function of the trivial
lift (unit \emph{Exchange}); the sum of the squares of its values is the squared norm, whose normalisation is
\m{\ialg{q}{\vartheta}}. (c) Let $v$ be the integral vector of Input~2\,(d) for the monomial~\m{x^{\alpha}}, and let \m{F(h_1,h_2,h_3)} be
the point form (Input~2\,(e)) with value~$v$ at the block diagonal point \m{(1,h_2,h_3,h_1)} of the index set
\m{12+4+4+(2z_p-20)}, with the reductions prescribed by~$x_0$ and inclusion matrices. Put
\m{\Xi=\xiSum{q_2}{\vartheta}{h_2\mapsto\xiSum{q_3}{1}{h_3\mapsto\xiSum{q}{1}{F(\cdot,h_2,h_3)}}}}. The formula follows from~(b), applied
three times, once the pairing of the point form with a kernel form is computed: by the descent of a point form
and the seesaw for the four blocks it is the product of the kernel of the rank-$12$ block with the monomial, by
Input~2\,(d), and of the kernel forms of the three flat blocks, of jet degree~$0$. Integrality: by~(a) three times,
and because the point form of the integral vector~$v$ has gauge at most~$1$ at its point, $v$ being fixed by the
matrices that fix the columns of the monomial. Tested: a form in the kernel of the descent, in the image of the
label idempotent and killed by a Hecke test satisfies the hypotheses of the vanishing of periods (Input~2\,(e)),
for the splitting into the first three blocks and the last one, and the class sum over the last block that
defines~$\Xi$ pairs to zero with it.
\end{proof}
\leanline{\lean{xiSummem}\sep \lean{smulapplyxiSum}\sep \lean{existstestform}}

\begin{proof}[Proof of Theorem B\acct]
The argument is that of the formal proof; the powers of the class number, of the vacuum scale and of the
denominator bounds are not recomputed (§\ref{s:notes}). Write \m{\vartheta_{\ell,\iota}=\theta_{\ell,\iota}^{\,j}}, a twist,
and let $G$ be the finite group of the characters of weight~$0$ with tame part trivial outside the chosen primes
that occur in Lemma~\ref{l:pigeon}\,(a) for \m{\UU(1)} at the chosen level. That level is free (unit \emph{Data}), so
the class sum \m{\clU} of the constant~$1$ is the number of classes.

\emph{Choices.} For every \m{\chi\in G} take pairs \m{q_1(\chi)}, flat for \m{\{1,\dots,2z_p-20\}}, and \m{q_3(\chi)}, flat
for \m{\{1,\dots,4\}}, attached to~$\chi$ by Lemma~\ref{l:flat}\,(a), and let $q_2$ be the flat pair for
\m{\{1,\dots,4\}} attached to~\m{\chi_D}; it is ramified. The orders \m{\mathbf{m}_\tau=c_{p,\tau}-z_p}, the jet degrees
of~$p$, satisfy \m{w_D(\tau)\le\mathbf{m}_\tau\le w_D(\tau)+x^{42}}, because $p$ is of rank one; take \m{\tm} from
Input~2\,(b). Let $V$ be a complex number with
\[V^{|G|}=\prod_{\chi\in G}\ialg{q_3(\chi)}{1}\cdot\bigl(\ialg{q_1(\chi)}{1}\cdot\tm\cdot\ialg{q_2}{1}\bigr).\]
Then \m{V\ne0}, and \m{|\sigma(V)|\le e^{B_0(D)}} for all \m{\sigma\in\AutC}: by Lemma~\ref{l:pigeon}\,(b) for~$\sigma$ and
\m{\psi=1}, \m{|\sigma(V)|} is at most the absolute value of a conjugate of a product of three normalised values and
of~\m{\tm}, which Input~2\,(b), (c) and the scales bound.

\emph{The valuation.} Claim: for every index \m{(\ell,\iota)} of the family and every embedding
\m{\varphi\colon\CC\to\CC_\ell}, with \m{\vartheta=\vartheta_{\ell,\iota}} the twist of the index, there is a test form~$\Xi$
at~$\ell$ for the $L$-value datum of~$p$ at~$\vartheta$ with
\begin{equation}\label{e:val}
|\varphi(V)|_\ell\cdot\ell^{-F_D(\ell)}\le\bigl|\varphi\bigl(\clU^{2}\,\ip{\Xi}{\triv{p}{\vartheta}}/\rat{q_2}{\vartheta}{1}\bigr)\bigr|_\ell .
\end{equation}
In this claim the index fixes the twist and the embedding varies. Of the index the proof uses only that $\ell$ is
prime to \m{N(Q(\vartheta))}, which is clause~(1) of fitness; Input~2\,(b) and Lemma~\ref{l:pigeon} hold for every
embedding. Indeed, Input~2\,(b), for the modulus~\m{\bar\pi_g} and the embedding~$\varphi$, gives $\alpha$ and~$x_0$
such that the Taylor coefficient \m{s(1)}, where \m{s(g)} is the kernel of \m{(I_1,I_{12})} with the indicator
of~$x_0$ and the monomial~\m{x^{\alpha}} at~$g$, satisfies
\[\ell^{-F_D(\ell)}\le\Bigl|\varphi\Bigl(\kappa_0^{12}\prod_\tau\mathbf{m}_\tau!\cdot s(1)/\tm\Bigr)\Bigr|_\ell .\]
Lemma~\ref{l:pigeon}\,(a) gives $z$ and \m{\psi\in G} with \m{|\varphi(s(1))|_\ell} at most the absolute value of the
mean of \m{(\psi s)(gz)}, and Lemma~\ref{l:pigeon}\,(b) gives \m{\psi_3\psi_1=\psi} with
\[|\varphi(V)|_\ell\le\bigl|\varphi\bigl(\ialg{q_3(\psi_3)}{1}\,\ialg{q_1(\psi_1)}{1}\,\tm\,\ialg{q_2}{1}\bigr)\bigr|_\ell .\]
Lemma~\ref{l:test}\,(c), for the pairs \m{q_1(\psi_1)}, $q_2$, \m{q_3(\psi_3)} and the character~$b$ determined by them
and~$\psi$, gives~$\Xi$. Its pairing with the trivial lift, which is the class sum over \m{\UU(1)} of the kernel
forms, is by the formula of Lemma~\ref{l:test}\,(c) the mean of \m{(\psi s)(gz)} times
\m{\ialg{q_1(\psi_1)}{1}\,\ialg{q_3(\psi_3)}{1}\,\ialg{q_2}{\vartheta}\prod_\tau\mathbf{m}_\tau!/\alpha!} and powers of~\m{\clU}
and~$\kappa_0$. Since \m{\ialg{q_2}{\vartheta}/\rat{q_2}{\vartheta}{1}=\ialg{q_2}{1}} and \m{|1/\alpha!|_\ell\ge1}, the three
inequalities combine to~\eqref{e:val}.

\emph{The height.} Consider the family
\[y'_{\ell,\iota}=\bigl(\rat{q_2}{\vartheta_{\ell,\iota}}{1}\big/\|\triv{p}{\vartheta_{\ell,\iota}}\|\bigr)^2\big/\clU^{3}.\]
It has height at most \m{2B_0(D)+B_0(D)+\sum_\tau\operatorname{Ask}_D(w_D(\tau),\tau)}, with the witness \m{H=V\bar V}. The
conjugates of~$H$ are at most \m{e^{2B_0(D)}}, because \m{\sigma(\bar V)} is the image of~$V$ under the automorphism
\m{\sigma\circ\mathrm{conj}}. Fix an index \m{(\ell,\iota)} and write \m{\vartheta=\vartheta_{\ell,\iota}}. The
claim~\eqref{e:val} for this index, with \m{\varphi=\iota} and with \m{\varphi=\iota\circ\mathrm{conj}}, that is twice
for the same twist~$\vartheta$, gives test forms $\Xi$, $\Xi'$ with
\[|\iota(V\bar V)|_\ell\,\ell^{-2F_D(\ell)}\le\bigl|\iota\bigl(\clU\,\ip{\Xi}{\operatorname{pr}\Xi'}\bigr)\bigr|_\ell\cdot
\bigl|\iota\bigl(\clU^{3}\|\triv{p}{\vartheta}\|^2/\rat{q_2}{\vartheta}{1}^2\bigr)\bigr|_\ell ,\]
where \m{\operatorname{pr}} is the orthogonal projection to the line \m{\CC\,\triv{p}{\vartheta}}, so that
\[\ip{\Xi}{\operatorname{pr}\Xi'}=\ip{\Xi}{\triv{p}{\vartheta}}\cdot\ip{\triv{p}{\vartheta}}{\Xi'}\big/\|\triv{p}{\vartheta}\|^2;\]
here one uses that the ratio is real.

The first factor is at most \m{e^{\beta_\ell}} by the rank-one bound of the unit \emph{Selmer}: by clause~(2)
of fitness the line \m{\CC\,\triv{p}{\vartheta}} has denominators at most~$\beta_\ell$ at~$\iota$ for \m{\idem(X-1)}, the
class stabilisers of \m{\UU(1)} are trivial, and $\Xi$, $\Xi'$ are test forms at~$\ell$. Clause~(2) cannot be
applied to $\Xi$ and~$\Xi'$ directly: a test form is tested against the kernel of the descent, not against the
orthogonal complement of the line. The rank-one bound applies it to~$\Xi$ and to the sum of the translates
of~$\Xi'$ by the twisted central operators at representatives of the \m{\clU} classes of \m{\UU(1)}. This sum is
integral at~$\iota$ and tested against the orthogonal complement of the line, the trivial lift is fixed by these
operators, and the sum is where the factor~\m{\clU} comes from. So
\[|\iota(Hy'_{\ell,\iota})|_\ell\le\exp\bigl(\beta_\ell+2F_D(\ell)\log\ell\bigr),\]
and the sum of \m{2F_D(\ell)\log\ell} over the primes below a bound \m{\ge x_1} is
\m{\sum_\tau\operatorname{Ask}_D(w_D(\tau),\tau)} (unit \emph{Columns}).

The family
\m{\rat{q_2}{\vartheta_{\ell,\iota}}{1}^{-2}} has height at most~$2x^{999}$ by clause~(3) of fitness, for the ramified
pair~$q_2$ and the sign~$1$, and the constant family \m{\|\triv{p}{1}\|^2} has height at most
\m{B_0(D)-\sum_\tau\operatorname{js}_D(w_0(\tau))} by Lemma~\ref{l:flat}\,(b). The product of the three families is
\[\|\triv{p}{1}\|^2\big/\bigl(\clU^{3}\,\|\triv{p}{\vartheta_{\ell,\iota}}\|^2\bigr),\]
and \m{\rat{p}{\vartheta_{\ell,\iota}}{1}^{-1}} is this number times a natural number, a power of~\m{\clU} times a power of \m{N(Q(\vartheta_{\ell,\iota}))}, by the definition of the
normalised values; a height bound passes to a family whose $\ell$-adic absolute values are not larger (unit
\emph{Heights}). Heights add: \m{4B_0(D)+2x^{999}-\mathsf{U}(D)\le6B_0(D)-\mathsf{U}(D)}, since \m{x^{999}\le B_0(D)}.
\end{proof}
\leanline{\lean{heightLeinvratio}}

\begin{corollary}[the product of Hecke values]\label{c:hv}
In Setting~\ref{set:data}, let \m{\theta} be a family of $m$-fold twists of~$D$ that is fit for~$D$ and $j$ a lift
index. Then the family
\[(\ell,\iota)\longmapsto\Bigl(\frac{\heckeValueProd{j}{\theta_{\ell,\iota}^{\,j}}}{\heckeValueProd{j}{1}}\Bigr)^{-1}\]
has height at most \m{10\,m\,x^{999}+k(D)\,(6B_0(D)-\mathsf{U}(D))}.
\end{corollary}
\begin{proof}
Let $\vartheta$ be a twist. By its definition \m{\heckeValueProd{j}{\vartheta}} is the product, over the three line
pairs \m{(\nu_o,z_o)=\linePair{j}{o}} with their signs~$s_o$, of a power of the norm of~\m{T_\vartheta} and of the
squares \m{|L^{T_\vartheta}(z,(\nu_o\vartheta^{s_o})^{-1})|^2}, \m{1\le z\le z_o}. Since $T_1$ is the set of the chosen
primes, the ratio to the untwisted product is, up to an integer power of \m{N(Q(\vartheta))}, the square of the
product of the ratios
\[\rat{(\nu_o,z)}{\vartheta}{s_o},\qquad1\le z\le z_o,\]
over the three line pairs. The sign of each ratio is the sign~$s_o$ of its line pair: the line pair of the
sign~$-1$ contributes the ratios \m{\rat{(\nu_o,z)}{\vartheta}{-1}}, and not those for the sign~$1$. For the first
line pair, where \m{s_o=1}, and \m{7m/2<z\le\heckeCount{D}}, the pairs \m{(\nu_o,z)} are of rank one for~$j$, and
there are \m{\lfloor k(D)/2\rfloor} of them. All the other pairs are ramified, and there are \m{7m/2+m+m/2=5m} of
them. Now take \m{\vartheta=\theta_{\ell,\iota}^{\,j}}. Theorem~B for the first, clause~(3) of fitness with the
sign~$s_o$ for the others, the rules for products and squares of heights, and the rule for powers of
\m{N(Q(\vartheta))} (unit \emph{Heights}) give the bound
\m{2\,\bigl(\lfloor k(D)/2\rfloor\,(6B_0(D)-\mathsf{U}(D))+5m\,x^{999}\bigr)}, and \m{2\lfloor k(D)/2\rfloor=k(D)} because
\m{k(D)} is even.
\end{proof}
\leanline{\lean{heightLeinvheckeValueProd}}

\subsection*{The volume of a lift: Theorem A}\label{s:volume}
\begin{lemma}[the compressed Gram operator]\label{l:gram}
Let $T$, $U$, $V$ be finite-dimensional inner product spaces, \m{\operatorname{gl}\colon T\to V}, \m{u\colon U\to V} and
\m{A\colon T\to U} linear maps and \m{\gamma,\varkappa\in\CC}.
\begin{enumerate}[label=\textup{(\alph*)},leftmargin=2.3em,itemsep=2pt]
\item If \m{u^{*}\circ\operatorname{gl}=\gamma A}, then \m{\operatorname{gl}^{*}\circ u=\bar\gamma A^{*}}.
\item With $L$ the image of~$u$:
\m{\vf(u)^2\cdot\det\compressedGramOp{\varkappa}{L}{\operatorname{gl}}=\varkappa^{\dim L}\cdot\vf(\operatorname{gl}^{*}\circ u)^2},
where \m{\vf} is the volume factor.
\end{enumerate}
\end{lemma}
\begin{proof}
(a) Take adjoints. (b) If $u$ is not injective, both sides vanish: \m{\vf(u)=0}, and \m{\operatorname{gl}^{*}\circ u} is
not injective, so \m{\vf(\operatorname{gl}^{*}\circ u)=0}. Let $u$ be injective; it is then a bijection from~$U$
onto~$L$. The operator is $\varkappa$ times \m{w^{*}w} for \m{w=\operatorname{gl}^{*}|_L}, since on~$L$
\[\ip{z}{\operatorname{pr}_L\operatorname{gl}\operatorname{gl}^{*}z'}=\ip{\operatorname{gl}^{*}z}{\operatorname{gl}^{*}z'};\]
and \m{\operatorname{gl}^{*}\circ u=w\circ u}, with $u$ regarded as a bijection onto~$L$, so that
\[\vf(w\circ u)^2=\det(w^{*}w)\cdot\vf(u)^2\qquad\text{and}\qquad
\varkappa^{\dim L}\,\vf(\operatorname{gl}^{*}\circ u)^2=\det(\varkappa\,w^{*}w)\cdot\vf(u)^2 .\]
\end{proof}
\leanline{\lean{adjointcomp}\sep \lean{detcompressedGramOp}}

\begin{setting}[a comparison map]\label{set:gl}
In Setting~\ref{set:data}: $j$ is a lift index, $\vartheta$ a twist of~$D$,
\m{\operatorname{gl}\colon\operatorname{Tab}_D(j)\to\FormVx{O_D(\vartheta,j)}} a linear map, $L$ a linear subspace of
\m{\FormVx{O_D(\vartheta,j)}} and $\beta$ a real number, such that \m{\operatorname{gl}^{*}} is injective on~$L$ and, for every
\m{z\in L}, the form \m{\operatorname{gl}(\operatorname{gl}^{*}z)} is tested against~$L^{\perp}$ for the good idempotent
\m{e_0(D)}.
\end{setting}

\begin{proposition}[the determinant at the primes and at the conjugates]\label{p:det}
In Setting~\ref{set:gl}:
\begin{enumerate}[label=\textup{(\alph*)},leftmargin=2.3em,itemsep=2pt]
\item Let \m{(\ell,\iota)} be a prime with a field isomorphism from~$\CC$ to an algebraic closure
of~$\mathbb{Q}_\ell$. Suppose that for every table~$\zeta$ and every embedding \m{\varphi\colon\CC\to\CC_\ell}: if every
\m{\overline{\zeta_a}} has gauge at most~$1$ at the representative of its class for \m{\varphi\circ\mathrm{conj}}, then
\m{\operatorname{gl}(\zeta)} is integral at~$\varphi$; and that $L$ has denominators at most~$\beta$ at~$\iota$ for
\m{e_0(D)}. Then
\[\bigl|\iota\bigl(\det\compressedGramOp{\mathsf{h}_D(j)}{L}{\operatorname{gl}}\bigr)\bigr|_\ell\le\exp\bigl(\dim L\cdot(r\,\beta)\bigr).\]
\item Suppose that $L$ is the image of~$u_\vartheta$ and that, for all tables~$\zeta$ and all \m{\sigma\in\AutC},
\[\mathsf{h}_D(j)\,\nsq{\operatorname{gl}(\zeta)}{\sigma}\le e^{\beta}\sum_a\|\sigma(\zeta_a)\|^2 .\]
Then for every \m{\sigma\in\AutC}:
\[\bigl|\sigma\bigl(\det\compressedGramOp{\mathsf{h}_D(j)}{L}{\operatorname{gl}}\bigr)\bigr|\le\exp(\dim L\cdot\beta).\]
\end{enumerate}
\end{proposition}
\begin{proof}[Proof \textup{(}\textup{(a)}: an account of the formal proof; \textup{(b)}: complete\textup{)}]
(a) Move the class representatives of \m{\SmallClasses{D}} by rational points~\m{\delta_{\mathfrak{c}}} to points integral
above~$\ell$ (library). For an index \m{a=(i,\mathfrak{c})} and a coordinate~$e$ let \m{f_{a,e}} be the table that is zero
outside~$a$ and whose entry at~$a$ is the translate by~\m{\delta_{\mathfrak{c}}} of the vector dual to the coordinate~$e$
for the Fock pairing of the integral bases; then \m{\ip{f_{a,e}}{\zeta}} is \m{\mathsf{h}_D(j)} times the conjugate of the
coordinate~$e$ of the translate of \m{\overline{\zeta_a}}. The linear forms \m{z\mapsto\ip{f_{a,e}}{\operatorname{gl}^{*}z}} on~$L$
have no common zero, because \m{\operatorname{gl}^{*}} is injective on~$L$. By the extension of Auerbach's lemma (library), for the absolute
value \m{|\iota(\,\cdot\,)|_\ell}, some \m{\dim L} of them, with indices \m{c_1,c_2,\dots}, are the coordinates of a
basis \m{(z_b)} of~$L$ on which all the forms have \m{|\iota(\,\cdot\,)|_\ell\le1}. In this basis the matrix of the
compressed Gram operator has the entries \m{\ip{y_a}{\operatorname{pr}_Ly'_b}} with \m{y_a=\operatorname{gl}(f_{c_a})} and
\m{y'_b=\operatorname{gl}(\operatorname{gl}^{*}(\mathsf{h}_D(j)z_b))}. The forms~$y_a$ are integral at every embedding
into~$\CC_\ell$: the Fock pairings of the integral bases are integers, and the transport by~\m{\delta_{\mathfrak{c}}} and
the complex conjugation in the definition of the tables move the gauge to a representative that is integral
above~$\ell$, where a vector with integer coordinates has gauge at most~$1$. The forms~$y'_b$ are integral
at~$\iota$, because by the estimate of Auerbach's lemma the conjugates of the entries of
\m{\mathsf{h}_D(j)\operatorname{gl}^{*}z_b} have gauge at most~$1$ after the same transport, and they are tested by
Setting~\ref{set:gl}. The definition of
denominators at most~$\beta$ gives the bound, the rank of the first space being~$r$.
(b) The determinant is the product of the \m{\dim L} eigenvalues, so it suffices to show \m{|\sigma(\alpha)|\le e^{\beta}}
for an eigenvalue~$\alpha$ with eigenvector \m{z\in L}. Let \m{\zeta=\operatorname{gl}^{*}z}, which is not zero, put
\m{\tilde\sigma=\mathrm{conj}\circ\sigma\circ\mathrm{conj}}, and let \m{\zeta^{\sigma}} be the table with the entries
\m{\zeta^{\sigma}_a=\overline{\sigma^{-1}\bigl(\overline{\sigma(\zeta_a)}\bigr)}}. It is chosen so that
\[\tilde\sigma(\zeta^{\sigma}_a)=\sigma(\zeta_a)\quad\text{for all }a,\qquad\text{and}\qquad
\sigma\ip{\zeta^{\sigma}}{\zeta}=\sum_a\|\sigma(\zeta_a)\|^2,\]
without any assumption that $\sigma$ commutes with complex conjugation. Pairing
\m{\alpha z=\mathsf{h}_D(j)\operatorname{pr}_L\operatorname{gl}(\zeta)} with \m{\operatorname{gl}(\zeta^{\sigma})} and applying~$\sigma$ gives
\[|\sigma(\alpha)|\sum_a\|\sigma(\zeta_a)\|^2=\mathsf{h}_D(j)\,\bigl|\sigma\ip{\operatorname{gl}(\zeta^{\sigma})}{\operatorname{pr}_L\operatorname{gl}(\zeta)}\bigr|.\]
The norms \m{\|\cdot\|_\sigma} do not satisfy a Cauchy--Schwarz inequality at~$\sigma$ alone when $\sigma$ does not
commute with complex conjugation. The inequality of the library is
\m{|\sigma\ip{y}{x}|\le\|x\|_{\sigma}\,\|y\|_{\tilde\sigma}}; apply it to \m{y=\operatorname{gl}(\zeta^{\sigma})} and
\m{x=\operatorname{pr}_L\operatorname{gl}(\zeta)}. The size bound at~$\tilde\sigma$ for the table~\m{\zeta^{\sigma}} gives
\m{\mathsf{h}_D(j)\,\nsq{\operatorname{gl}(\zeta^{\sigma})}{\tilde\sigma}\le e^{\beta}\sum_a\|\tilde\sigma(\zeta^{\sigma}_a)\|^2=e^{\beta}\sum_a\|\sigma(\zeta_a)\|^2}.
The size bound at~$\sigma$ for~$\zeta$ gives the same bound for
\m{\mathsf{h}_D(j)\,\nsq{\operatorname{pr}_L\operatorname{gl}(\zeta)}{\sigma}}, because the projection to~$L$ does not increase
\m{\nsq{\cdot}{\sigma}} on the tested form \m{\operatorname{gl}(\zeta)=\operatorname{gl}(\operatorname{gl}^{*}z)} (unit \emph{Selmer};
$L$ is killed by the Hecke test at a split place off the level, which exists, since the lifts of good eigenforms
are, and $L$ is generated by such lifts, tuple by tuple, by Input~3). So the right side is at most the geometric
mean of the two bounds, \m{e^{\beta}\sum_a\|\sigma(\zeta_a)\|^2}, and the sum is positive.
\end{proof}
\leanline{\lean{normdetpadic}\sep \lean{normdetarch}}

\begin{proof}[Proof of Theorem A]
Write \m{n_0=\dim\operatorname{Row}(D)=\dim\Mgood{j}}, \m{\vartheta_{\ell,\iota}=\theta_{\ell,\iota}^{\,j}}, and note that \m{1^{j}=1}, so
\m{\lvol{1}{j}=\vf(u_1)^2} and \m{\lvol{\theta_{\ell,\iota}}{j}=\vf(u_{\vartheta_{\ell,\iota}})^2}. Take $\gamma$, $A$ and the
maps \m{\operatorname{gl}_\vartheta} of Input~1, and the numbers~$\beta_\ell$ of clause~(2) of fitness for~$j$. For a
twist~$\vartheta$ let \m{L_\vartheta} be the image of~$u_\vartheta$, of dimension~$n_0$ since $u_\vartheta$ is injective, and
\[\Delta(\vartheta)=\det\compressedGramOp{\mathsf{h}_D(j)}{L_\vartheta}{\operatorname{gl}_\vartheta}.\]
By Input~1\,(ii) and Lemma~\ref{l:gram},
\begin{equation}\label{e:volid}
\vf(u_\vartheta)^2\cdot\Delta(\vartheta)=\mathsf{h}_D(j)^{n_0}\,\vf(\bar\gamma(\vartheta)A^{*})^2=\mathsf{h}_D(j)^{n_0}\,|\gamma(\vartheta)|^{2n_0}\,\vf(A^{*})^2 .
\end{equation}
The right side is not zero: \m{\gamma(\vartheta)\ne0}, and $A^{*}$ is injective because $A$ is surjective. So
\m{\Delta(\vartheta)\ne0}; hence the compressed Gram operator is invertible and \m{\operatorname{gl}_\vartheta^{*}} is injective
on~$L_\vartheta$. For \m{z=u_\vartheta y\in L_\vartheta}, \m{\operatorname{gl}_\vartheta(\operatorname{gl}_\vartheta^{*}z)=\bar\gamma(\vartheta)\operatorname{gl}_\vartheta(A^{*}y)}
is tested against \m{L_\vartheta^{\perp}} for \m{e_0(D)} by Input~1\,(iv). So Setting~\ref{set:gl} holds for every twist. Dividing \eqref{e:volid} for~$1$ by
\eqref{e:volid} for \m{\vartheta_{\ell,\iota}} and using Input~1\,(i),
\[\frac{\lvol{1}{j}}{\lvol{\theta_{\ell,\iota}}{j}}=\frac{\Delta(\vartheta_{\ell,\iota})}{\Delta(1)}\cdot\Bigl(\Bigl(\frac{\heckeValueProd{j}{\vartheta_{\ell,\iota}}}{\heckeValueProd{j}{1}}\Bigr)^{-1}\Bigr)^{n_0}.\]
The first family has height at most \m{n_0\cdot3rB_0(D)}, with the witness \m{H=\Delta(1)}: its conjugates are at most
\m{\exp(n_0\cdot2rB_0(D))} by Proposition~\ref{p:det}\,(b) and Input~1\,(vi); and
\m{|\iota(\Delta(\vartheta_{\ell,\iota}))|_\ell\le\exp(n_0\,r\,\beta_\ell)} by Proposition~\ref{p:det}\,(a), Input~1\,(v), clause~(1) of
fitness and clause~(2) for the image of the adjoint lift; the sum of the~$\beta_\ell$ is at most \m{B_0(D)}. The second
family has height at most \m{n_0\,(10mx^{999}+k(D)(6B_0(D)-\mathsf{U}(D)))} by Corollary~\ref{c:hv}. Heights add, and
\[3rB_0(D)+10mx^{999}+6k(D)B_0(D)\le k(D)\log b\cdot B_0(D)\]
is one of the inequalities between the scales.
\end{proof}
\leanline{\lean{heightLeliftVolumeratio}}

\sect{Where the results are used}{s:used}
Theorem~A is the second of the two inputs of Theorem~C of the handout named in the head of the argument; the unit
\emph{Endgame} multiplies the $m$ families for \m{j=1,\dots,m}, adds their bounds, and combines the result with the
factorisation of the volume of the chain (Theorem~A of the first note). Theorems C and~D and Proposition~\ref{p:lineint} are the inputs of the comparison forms in the sister
note. Lemma~\ref{l:carrier} and the definitions of §\ref{s:setting} are used in this unit only.

\sect{Notes on the formal text}{s:notes}
\begin{enumerate}[label=\textup{\arabic*.},leftmargin=1.6em,itemsep=5pt]
\item \emph{Statements cut at a local definition.} The formal statements of Lemma~\ref{l:coords}\,(b) and~(c), of
Proposition~\ref{p:lineint} and of Lemma~\ref{l:test}\,(c) contain a local definition; the tools of this write-up
read such a statement only up to it and therefore list no argument for Proposition~\ref{p:lineint} and only the
point~$z$ for Lemma~\ref{l:coords}\,(b). The statements above were written from the Lean text: the first takes the
lift index, the second the lift index, the line choice, the matrix~$Y$ and the homomorphism~$\varphi$ besides~$z$.
\item \emph{Arguments that only a compiler can settle.} The text reading lists letters that the statements bind
themselves or that belong to other declarations of the same section, and that Lean presumably does not take: a
natural number~$h$ in the definitions of \m{\negIx{j}{i}}, \m{\DetIx{j}{\tau}}, \m{\linePoly{j}}, \m{\carrierChar{j}} and
\m{\detRows{j}{i}{Y}{\tau}} and in Lemma~\ref{l:counts}\,(b); a complex number~$k$ in Theorem~A; the line choice and the
matrix~$Y$ in Lemma~\ref{l:coords}\,(c), where they are quantified inside. The five definitions therefore carry no
notation made by the build, and the note writes them with the arguments that their uses show. None of this was
confirmed by asking Lean.
\item \emph{Constants that may vanish.} In Theorem~C the constant~$C$ is not shown to be non-zero; that the
numbers \m{c(\vartheta)} are not zero is proved in the sister note from Theorem~D.
\item \emph{Truncated and integer exponents.} \m{t(D)-9m-1-r}, \m{2z_p-20} and \m{3m-j} are differences of natural
numbers, truncated at zero in the formal text; \m{m/2} is an integer because $m$ is even; the exponent of
\m{\upsilon} in \m{\tableTheta{j}{\zeta}{g}} is the integer \m{m/2-j}, and that in \m{\carrierChar{j}} is an integer that may be
negative.
\item \emph{Theorem~A.} The formal bound is the product of \m{\dim\operatorname{Row}(D)} and of
$k(D)\cdot(\log b\cdot B_0(D)-\mathsf{U}(D))$, with \m{\log b} a field of the data. The family is indexed by the pairs of a prime and a field isomorphism from~$\CC$
to an algebraic closure of~$\mathbb{Q}_\ell$; the embeddings into~$\CC_\ell$ of the statements about gauges are obtained
from them.
\item \emph{What was not reconstructed.} In the proof of Theorem~B the exact normalising factors (the powers of
the class number, of the vacuum scale and of the denominator bounds) are balanced in the formal text by
simplification; the proof above tracks the absolute values that matter and says where a factor is harmless. In
Proposition~\ref{p:approx} the function~$K_0$ is not written in the formal text, Lean infers it; the formula above is
my reading of what it must be. The referee of this note compared it with the terms that the formal proof constructs
and found it in agreement.
\item \emph{The sign of the ratio.} The formal ratio of a pair takes three arguments: the pair, the tame character
and a sign, which has the default value~$1$ and multiplies the exponent of the tame character in the numerator.
Theorem~B and the claim in its proof are written in the formal text with the default; the note writes the sign~$1$
in the third slot. Clause~(3) of fitness is stated for every sign, and the proof of Corollary~\ref{c:hv} uses it with
the sign of each line pair, \m{-1} for the third.
\item \emph{The index and the embedding in the proof of Theorem~B.} The formal claim behind~\eqref{e:val} is stated
for an index of the family, a prime with a field isomorphism, and separately for an embedding into~$\CC_\ell$; the
twist is that of the index. It is applied twice to one index, at the embedding of the isomorphism and at its
composite with complex conjugation.
\item \emph{Letters that differ from the formal text.} The scalar of the compressed Gram operator is~$k$ in the
formal text and~$\varkappa$ here; the number of the indicator jet is~$t$ there and \m{\tm} here; the class sum of the
constant~$1$ over \m{\UU(1)} is~$N$ there and \m{\clU} here; the degrees of Lemmas \ref{l:coords} and~\ref{l:sparse}
are~$d$ there and~$\mathsf{d}$ here. The letters $k$, $t$, $N$ and~$d$ have other meanings in this write-up.
\end{enumerate}

\subhead{How this note was written.} From the dossier of the unit, whose reader calls the step claims of the long
analytic proofs a first reading and notes that no second reader had looked at them. After the dossier was frozen, the second reader of the write-up compared the readings of all statements of the unit
with the Lean text and sampled its structured proofs (her message to the writer of 26 November 2026); her
corrections are followed where they concern this note: the first minimum is cut off at~$1$ (confirmed in the Lean
definition); in the proof of isolation the approximant at angle zero is kept apart from the bounded target; the
second family in Proposition~\ref{p:det}\,(a) is \m{\operatorname{gl}(\operatorname{gl}^{*}(\mathsf{h}_D(j)z_b))} and the tables are
transported before integrality is read off; the auxiliary table in Proposition~\ref{p:det}\,(b) carries the two complex
conjugations; the size bound of Proposition~\ref{p:lineint}\,(ii) is stated for every tame character with empty set of
primes, not only for the trivial one; and in the proof of Theorem~A the power of the conductor has exponent zero
and the local exponents carry no factor \m{\log\ell}. Her exact identity for the ratio of the volumes is the one
derived above. I read completely the Lean
files of the height of the lift volume, of the Hecke values, of the separation and of the approximation, and the
statements and docstrings of the other three. The definitions of the line (Definitions \ref{d:index}
to~\ref{d:hv}) I took from the dossier's reading and compared with the docstrings; the referee of this note then
compared them with the Lean text of the two files (the sizes of the index sets, the two index maps, the index set
of the determinants, the exponent of the line polynomial, the rows of the determinants, the datum, the line
multiplier, the function \m{F^{\natural}}, the exponent and the characters, the line pairs, and the signs in the
product of Hecke values) and found them in agreement. The proofs of Theorem~B and of Proposition~\ref{p:approx} I read line by line in
the Lean text. For Theorem~B I could reconstruct the argument: the choice of~$V$, the two pigeonholes, the test
form, the use of the three clauses of fitness and the bookkeeping of heights; the powers of the normalising
constants I did not recompute. For Proposition~\ref{p:approx} I could reconstruct the architecture (the near maximum
of the carrier, the main term as an integral against holomorphic minors, the sparse theta sum, the reproducing
formula and the finite Fourier expansion), not the estimates, which the formal proof closes with the inequalities
between the scales. The proof of Theorem~A is complete relative to its inputs and follows the formal one; I
checked the identity of the volumes and the count of the heights by hand. I read the Lean definitions of a fit
family, of denominators, of heights, of the budget for congruences and of the cost of one direction. The program
check of the unit's dossier covered \UnitSteps{} leaf steps of \UnitProofs{} proofs.

\subhead{Revision.} This text was revised after the report of an independent referee (26 November 2026), who
compared the statements with the Lean text and re-derived the complete proofs. The changes: the ratio of a pair
carries its sign, as the formal definition does, in fitness, in Theorem~B and in Corollary~\ref{c:hv}, where the
sign~$-1$ is needed; the partial $L$-function is that of the base change, without the Euler factors at the
conjugate places as well; in the proof of Theorem~B the claim on valuations separates the index from the embedding
and the rank-one bound of the unit \emph{Selmer} is used where clause~(2) of fitness does not apply directly; in
Proposition~\ref{p:det}\,(b) the Cauchy--Schwarz inequality is the conjugated one; the proofs of
Lemmas \ref{l:gram}\,(b) and~\ref{l:sparse}\,(a) and of Proposition~\ref{p:iso} are completed; ``tested'' and
``denominators'' name their idempotent; the references of the background results are those of the catalogue with
the second reader's verdicts; and each proof says where it stands whether it is complete, an account or an
outline.

\head{Appendix A. Concordance with the Lean artifact}
{\fontsize{8.6}{11.6}\selectfont\renewcommand{\leanfont}{\ttfamily\fontsize{7.7}{10}\selectfont\color{leancol}}
\setlength{\LTleft}{0pt}\setlength{\LTright}{\fill}\setlength{\tabcolsep}{0pt}
\begin{longtable}{@{}p{30mm}p{56mm}p{43mm}r@{}}
\textbf{In this note} & \textbf{Lean declaration} & \textbf{File} & \textbf{Lines}\\\midrule\endhead
\ConcordanceRows
\end{longtable}}

\subhead{The four theorems as they are written in the artifact.}
\quoted{heightLeliftVolumeratio}
\quoted{heightLeinvratio}
\quoted{cosetKernelSumeqfNatural}
\quoted{eqzeroofintegral}
In these texts {\leanfont HeightLe} is ``height at most'', {\leanfont D.IsFitFamily} fitness, {\leanfont D.rowDim}
the dimension of the rows, {\leanfont D.logB} the number \m{\log b}, {\leanfont D.B₀} the budget for congruences,
{\leanfont D.dirCost} the cost of one direction, {\leanfont D.liftVolume} the lift volume, {\leanfont p.ratio} the
ratio \m{\rat{p}{\vartheta}{s}}, whose sign~$s$ is~$1$ where the formal text does not write it,
{\leanfont cosetKernelSum} the coset kernel sum and {\leanfont D.goodForms} the good forms.

\head{Notation}
\begingroup\fontsize{8.6}{10.9}\selectfont\raggedright
\setlength{\columnsep}{9mm}\setlength{\parskip}{0pt}
\begin{multicols}{2}
\nt{$P=(N_0,d)$, $K$, $D$, $E(D)$, $\OEr$}{parameters; the imaginary quadratic field; data; its field; integers}
\nt{$x$, $x_1$, $b$, $m$, $k(D)$, $r$, $t(D)$, $f(D)$, $g(D)$}{scales; number of lifts; ranks; step; number of
infinite places; second level}
\nt{$\bar\pi_g$, $M_D$, $b_1$, $b_2$, $\mathsf{g}(D)$, $\mathsf{g}'(D)$}{the integer of norm \m{g(D)}; the chosen group order;
auxiliary bounds; the gap and the reduced half gap}
\nt{$w_D(\tau)$, $w_j(\tau)$, $\operatorname{Ask}_D(z,\tau)$, $F_D(\ell)$}{nominal weight; weight of the $j$-th lift; prime sum;
forced part}
\nt{$B_0(D)$, $\mathsf{U}(D)$, $\operatorname{js}_D(n)$, $\Omega$, $\operatorname{budget}(D)$, $h_{\mathrm{main}}(D)$}{budget for
congruences; cost of one direction; jet size; period constant; budget; main amount}
\nt{$\chi_D$, $\omega_D$, $\mu^{\flat}_D$, $\mu^{0}_j$, $\upsilon$}{the fixed characters; the character of the Weil data}
\nt{$\tau$, $\sigma$, $\ell$, $p$ (prime), $\iota$, $\iota_p$, $\varphi$, $\CC_\ell$}{an infinite place; automorphisms and
embeddings; primes}
\nt{$j$, $\vartheta$, $Q(\vartheta)$, $N(Q(\vartheta))$, $\mathfrak{m}_D(\vartheta,c)$, $\theta$}{lift index; twist; its primes
and their norm; a modulus; a fit family}
\nt{$\beta_\ell$, $\beta$, $\mathcal{Z}$, $r_0$, $\operatorname{pr}_{\mathcal{Z}}$, $\operatorname{pr}$}{denominator exponents; a
subspace; a rank; orthogonal projections}
\nt{$e_0(D)$, $\idem(\mathbf{c})$, $\mathbf{c}$, $\idem(X-1)$, $e$}{the good idempotent; the idempotent of a tuple of
labels; the label idempotent; an idempotent}
\nt{$W_j(D)$, $\Mgood{j}$, $O_D(\vartheta,j)$, $\alift_D(\vartheta,j)$, $u_\vartheta$, $\lvol{\vartheta}{j}$}{objects of the first
note}
\nt{$\ix r,\ix n,\ix a,\ix L,\ix N$, $\ix n'$}{finite index sets}
\nt{$\UU(S)(\AAA)$, $\AAA$, $\AAA_f$, $G_\infty(s)$, $\KK{3m}$, $\clsum{h}{\Gamma}$, $\Gamma$}{groups; adeles; levels; class sums}
\nt{$\tk{S}{T}{\alpha}{\beta}{\mathfrak{m}}{\Psi}{P}(g,h)$, $\LP{T}{\nu}{d}$, $L^{T}(z,\nu)$, $T_\vartheta$}{theta kernel;
product of $L$-values; partial $L$-function of the base change, without the Euler factors at~$T$ and at the
conjugate places; the chosen primes and those of \m{Q(\vartheta)}}
\nt{$\kappa^{\pm}$, $\kappa$, $\operatorname{pol}_\kappa(\xi)$, $\operatorname{pol}^{-}$, $\|\xi\|_\kappa$}{weights; polynomials and norms
of coordinate vectors}
\nt{$\FormVx{O}$, $\gau{\iota}{\kappa}{\epsilon}{g}{\xi}$, $\nsq{\phi}{\sigma}$, $\cks{O}{b}{w}{e}{y}(g)$}{forms; gauge; norm
at~$\sigma$; coset kernel sum}
\nt{$\carrier{E}{m}$, $\cXi{m}{c}{h}$}{the carrier; its class sum}
\nt{$p$, $q$, $q_1$, $q_2$, $q_3$ (pairs), $\nu_p$, $c_p$, $z_p$, $\psi_p$}{pairs and their components}
\nt{$\rat{p}{\vartheta}{s}$, $s$ (sign), $\triv{p}{\vartheta}$, $\ialg{p}{\vartheta}$, $\delta_q$, $X$, $O_X$}{ratio of the
pair~$p$ for the tame character~$\vartheta$ and the sign~$s$; trivial lift; normalised value; denominator bound; an
$L$-value datum and its pair}
\nt{$\untwistedModulus{D}$, $\forgetTwist{\vartheta}{x}$, $\untwistedAuxPair{j}$}{Definition~\ref{d:index}\,(1)}
\nt{$\CarrierIndex{D}$, $\FlatIndex{D}$, $\SmallIndex{D}$, $\SmallClasses{D}$, $\mathfrak{c}$,
$h_{\mathfrak{c}}$}{index sets; classes and representatives}
\nt{$\PadIndex{j}$, $\BigIndex{j}$, $\LPlusIndex{j}$, $\padSigns{j}$}{index sets of the $j$-th lift; signs}
\nt{$\basisVec{D}{i}$, $\negIx{j}{i}$, $\posRowIndex{j}{h}{i}$, $\regroupEquiv{j}$, $\posEquiv{j}$}{indices and
bijections}
\nt{$\twistedChi{\vartheta}{j}$, $\bigCharExp{j}$, $\bigChar{\vartheta}{j}$, $\carrierChar{j}$}{characters}
\nt{$\LineChoice{D}$, $\shortRows{D}{i}$, $\MemberIndex{j}$, $i$, $i_1$, $i_2$, $i_c$, $i_l$}{line choices, short
rows, choices}
\nt{$\DetIx{j}{\tau}$, $\linePoly{j}$, $\detRows{j}{i}{Y}{\tau}$, $\linePolyCoords{j}{i}{Y}$, $\Delta(j,i,Y)$,
$\Delta_{i_2}$}{determinants and coordinates of the line}
\nt{$\lPlusDatum{\vartheta}{j}{u}{z}$, $\lineMultiplier{\vartheta}{j}{i}{h}{l}$, $\fNatural{j}{i}{g}$, $\psi'$,
$\psi'_t$}{datum, line multiplier, \m{F^{\natural}}; additive characters}
\nt{$S$, $W'$, $L^{+}$, $I_j^{-}$, $I^{\flat}$, $V''$, $V^{\circ}$, $\mathbf{1}_I$, $\mathbf{1}_{x_0}$,
$\Theta^{+}_{l}$}{spaces of the construction; indicators; a kernel}
\nt{$\eSign{D}{s}$, $\heckeCount{D}$, $\linePair{j}{o}$, $\heckeValueProd{j}{\vartheta}$, $\tableTheta{j}{\zeta}{g}$}{Definition~\ref{d:hv}}
\nt{$\xiSum{q}{\vartheta}{F}$, $\compressedGramOp{\varkappa}{L}{\operatorname{gl}}$, $\varkappa$, $\vf$}{Definition~\ref{d:aux};
a complex scalar; the volume factor}
\nt{$\operatorname{Tab}_D(j)$, $\zeta$, $\zeta_a$, $\gamma$, $A$, $\operatorname{gl}_\vartheta$, $\mathsf{h}_D(j)$}{objects of Input~1}
\nt{$\|h\|$, $\operatorname{mh}_k$, $\lambda_1$, $\nu$, $\nu_D$, $B$, $f_E$, $o$, $w$, $k$ (weights)}{objects of
Setting~\ref{set:iso}}
\nt{$K_0$, $K$, $C$, $\varepsilon$, $C(u,y)$, $w^{\circ}$, $c_a$, $\Psi$, $\mu$, $q$ (integer), $c_N$,
$\jmath$}{objects of Propositions \ref{p:iso} and~\ref{p:approx} and Lemma~\ref{l:sparse}}
\nt{$G$, $V$, $\clU$, $s(g)$, $\mathbf{m}$, $\mathbf{m}_\tau$, $\tm$, $\alpha$, $x_0$, $\kappa_0$, $d_\tau$, $\Xi$, $\Xi'$,
$y'$}{objects of the proof of Theorem~B}
\nt{$\Delta(\vartheta)$, $L_\vartheta$, $n_0$, $H$, $T$, $U$, $u$}{objects of the proof of Theorem~A; a witness of a height;
spaces and maps}
\nt{$z_\vartheta$, $\chi_\vartheta$, $W_\vartheta$, $K^{(\vartheta)}$, $I$, $n_j$}{objects of the proof of isolation, with
$\vartheta$ an angle; an exponent}
\nt{$\rowpar{\tau}{q}$, $\tau_\sigma$, $\operatorname{sgn}(\sigma)$, $\tilde\sigma$}{row parameters; the place and the sign of
an embedding; the automorphism \m{\mathrm{conj}\circ\sigma\circ\mathrm{conj}}}
\nt{$N$, $t$, $\mathsf{d}$}{an integer in Lemmas \ref{l:mh} and~\ref{l:sparse}; a matrix in Lemma~\ref{l:sparse}; a
degree}
\nt{$c(\vartheta)$, $c_1$, $c_2$, $C_1$, $C_2$, $u$, $w$ (matrices), $\phi$, $\check\Psi_0$, $\mathfrak{n}$}{objects of
Theorem~C}
\nt{$F$, $g_0$, $y$, $z$, $a$, $e$, $l$, $h$, $d$, $c$, $c'$, $v$, $Y$, $Q$, $H$ (integer), $A$, $B$ (numbers),
$\Lambda$, $\mathcal{N}$, $m'$, $\delta_{\mathfrak{c}}$, $f_{a,e}$, $z_b$, $\zeta^{\sigma}$, $\psi$, $\chi$, $\epsilon$}{further
forms, points, indices, vectors, constants}
\nt{$L$}{in ``$L$-value'' and ``$L$-function''; also a subspace}
\nt{$\mathbb{Q}$, $\mathbb{Q}_\ell$}{the rationals; the $\ell$-adic numbers}
\end{multicols}
\endgroup

\end{document}
